% concatenated from Revised_final_NComm.tex
\documentclass[%
 reprint,
%superscriptaddress,
%groupedaddress,
%unsortedaddress,
%runinaddress,
%frontmatterverbose, 
% preprint,
% preprintnumbers,
%nofootinbib,
%nobibnotes,
%bibnotes,
 amsmath,amssymb,
 % aps,
 aip,
% pra,
%prb,
%rmp,
%prstab,
%prstper,
%floatfix,
% linenumbers
]{revtex4-2}


\usepackage{graphicx} 
\usepackage{amsmath,amssymb,amsfonts}

\usepackage[normalem]{ulem}

\usepackage{bm}% bold math
\def \bF {\pmb{F}}
\def \bM {\pmb{M}}
\def \bH {\pmb{H}}
\def \bD {\pmb{D}}
\def \bd {\pmb{d}}
\def \bG {\pmb{G}}
\def \bg {\pmb{g}}
\def \bL {\pmb{L}}
\def \bd {\pmb{d}}
\def \bu {\pmb{u}}
\def \bU {\pmb{U}}
\def \bx {\pmb{x}}
\def \by {\pmb{y}}
\def \bX {\pmb{X}}
\def \bY {\pmb{Y}}
\def \bz {\pmb{z}}
\def \bZ {\pmb{Z}}
\def \bw {\pmb{w}}
\def \bW {\pmb{W}}
\def \bs {\pmb{s}}
\def \bk {\pmb{k}}
\def \bv {\pmb{v}}
\def \bb {\pmb{b}}
\def \bp {\pmb{p}}
\def \bq {\pmb{q}}
\def \bB {\pmb{B}}
\def \bV {\pmb{V}}
\def \blambda {\pmb{\lambda}}
\def \balpha {\pmb{\alpha}}
\def \btheta {\pmb{\theta}}
\def \bbeta {\pmb{\beta}}
\def \boldeta {\pmb{\eta}}
\def \bone {\pmb{1}}
\def \diag {\text{diag}}

\usepackage{xcolor}
\usepackage{float}
\usepackage{mathtools}

\let\proof\relax
\let\endproof\relax


\usepackage{amsthm}

\newtheoremstyle{italichead}% <name>
  {6pt}%      Space above
  {6pt}%      Space below
  {\normalfont}% Body font
  {}%         Indent amount
  {\itshape}% Head font
  {.}%        Punctuation after theorem head
  {0.5em}%    Space after theorem head
  {}%         Head specification

\theoremstyle{italichead}
\newtheorem{theorem}{Theorem}
\newtheorem{lemma}[theorem]{Lemma}
\newtheorem{proposition}[theorem]{Proposition}
\newtheorem{corollary}[theorem]{Corollary}


\newtheorem{conjecture}[theorem]{Conjecture}
\newtheorem{assumption}[theorem]{Assumption}
\newtheorem{definition}[theorem]{Definition}




\newcommand{\argmin}[1]{\underset{#1}{\operatorname{arg}\,\operatorname{min}}\;}

\theoremstyle{definition}
%\newtheorem{definition}[theorem]{Definition}
\newtheorem{remark}{Remark}
\newtheorem{example}{Example}


\usepackage{algorithm}
% \usepackage{algorithmic}
\usepackage{algpseudocode}% http://ctan.org/pkg/algorithmicx
 \renewcommand{\algorithmicrequire}{\textbf{Input:}}
 \renewcommand{\algorithmicensure}{\textbf{Output:}}
\usepackage{varwidth}% http://ctan.org/pkg/varwidth
\usepackage{xcolor} % Include this package in your preamble for color support

\newcommand{\MA}[1]{\textcolor{purple}{#1}}

\usepackage{comment}

\usepackage{enumitem}

\usepackage{xurl}


\begin{document}


\author{Amirhossein Nazerian}
\affiliation{Department of Mechanical Engineering, University of New Mexico, Albuquerque, NM 87131, USA}
\author{Sahand Tangerami}
\affiliation{Department of Mechanical Engineering, K. N. Toosi University of Technology, Tehran, Iran}
\author{Malbor Asllani}
\affiliation{Department of Mathematics, Florida State University,
1017 Academic Way, Tallahassee, FL 32306, USA}
\author{David Phillips}
\affiliation{Department of Mechanical Engineering, University of New Mexico, Albuquerque, NM 87131, USA}
\author{Hernan Makse}
\affiliation{Levich Institute and Physics Department, City College of New York, New York, NY 10031, USA}
\author{Francesco Sorrentino}
\altaffiliation{Corresponding author. Email: fsorrent@unm.edu}
%\thanks{Corresponding author. E-mail address: fsorrent@unm.edu}}
%\email{fsorrent@unm.edu}
\affiliation{Max Planck Institute for the Physics of Complex Systems, 01187 Dresden, Germany}
\affiliation{Department of Mechanical Engineering, University of New Mexico, Albuquerque, NM 87131, USA}



% % \date{}
% \maketitle

%\title{Extreme vulnerability to intruder attacks in dynamical networked systems}

\title{Extreme vulnerability to intruder attacks destabilizes network dynamics}
%Intruders trigger dynamical instabilities in networked systems}

\begin{abstract}
Consensus, synchronization, formation control, and power grid balance are examples of desirable dynamical states that arise in networks. Here we investigate how such states can be destabilized by an intruder agent within an otherwise functioning network. We see that a single adversarial node, coupled through adversarial connections to one or more other nodes, is sufficient to destabilize the entire network. We further show that concentrating the attack on a single low-indegree node induces the greatest instability, challenging the common assumption that hubs are the most critical nodes. This leads to a new characterization of network vulnerability, identifying low-indegree nodes as the most vulnerable components. Although derived for linear systems, our results extend to nonlinear networks, including the Kuramoto model. These findings reveal an intrinsic vulnerability of technological, social, and biological networks.
\end{abstract}
%(\textbf{MA: Highly probable will change again, but just to summarize a bit everything.)}}
%    We study the network intruder attack, which occurs when an intruder node infiltrates a network and attempts to destabilize its dynamics. Under very general conditions, and for a broad variety of the dynamics such as consensus, formation control, and power grid power balancing, we show that a single intruder connected to the rest of the network can destabilize the network dynamics. We also show that the largest instability is achieved when the intruder connects to a single node of an existing network and when this one node has the lowest indegree. This contrasts with the common expectation that an attack will be more severe when applied to one the network hubs with high degree. Our work is particularly relevant to the security of complex technological networks, such as coupled autonomous systems, power grids, and sensor networks and points out to an intrinsic vulnerability of such networks, which could be exploited to destabilize them.
%\end{abstract}

\maketitle
%\noindent\textit{a) Corresponding author: Mail address: fsorrent@unm.edu}

\begin{comment}
\section*{Significant Statement}
{A large literature has studied how cascading processes can destabilize network dynamical states \cite{PhysRevE.66.065102, Buldyrev_2010, PhysRevE.69.045104, motter2004cascade,DENG2007714, motter2004cascade,requiao2015fast}. Here we consider an alternative route to instability in networked systems,  based on the presence of adversarial interactions between the nodes of a network. Within this context, a fundamental question}
%A fundamental question, with significant societal and technological implications, 
is how desired dynamical network states (such as coordinated dynamical states in networks of drones or autonomous vehicles) can be destabilized by the purposeful insertion of one or a few {adversarial agents.} Our work tackles this problem and points out a main vulnerability of complex dynamical networks, namely a single node connecting with adversarial connections to one or more other nodes is sufficient to destabilize the entire network dynamics. Moreover, surprisingly, the highest instability is generated when the attacked node has low in-degree. This work has direct implications to the attainment of consensus in opinion dynamics, synchronization in oscillator networks, and desired patterns in coupled autonomous systems. 
\end{comment}


\section{Introduction} 





Complex networks are ubiquitous in technological, biological, and social systems. It is often the case that the nodes of these networks evolve over time to reach a well-defined dynamical state, such as consensus in opinion dynamics, synchronization in oscillator networks, or a desired pattern in coupled autonomous systems \cite{6341809, 9356608, 1545539}. Ensuring the robustness of such networks against failures and attacks is a critical research challenge \cite{artime2024robustness}. Robustness has been studied extensively from both structural and dynamical perspectives \cite{granovetter1978threshold, pastor2001epidemic, castellano2009statistical, centola2010spread}. Traditional approaches to studying network robustness, including structural attacks such as cascading failures via percolation theory and node or edge removals, have offered valuable insights.  Structural interventions, such as cascade failures \cite{PhysRevE.66.065102, Buldyrev_2010, PhysRevE.69.045104, motter2004cascade}, node or edge removals \cite{DENG2007714, motter2004cascade}, and network fragmentation \cite{requiao2015fast}, examine how disruptions alter the network's connectivity. 


Alternatively, network instabilities may arise from the presence of adversarial connections between the nodes of a network \cite{Altafini2013Consensus,Chen2017spectral,ahmadizadeh2017eigenvalues}. In contrast to cascading failures, which are triggered by the removal of nodes or edges, adversarial interactions correspond to the addition of links with negative weights. Such interactions have been reported in brain networks as inhibitory connections between neurons \cite{ahmadizadeh2017eigenvalues}, in social networks as antagonistic relationships between agents \cite{Altafini2013Consensus}, and in distributed control systems as faulty or corrupted communication processes \cite{Chen2017spectral}.

While both cascading failures \cite{Buldyrev_2010,reis2014avoiding,albert2000error,PhysRevE.66.065102} and instabilities induced by adversarial interactions can disrupt network functionality, they operate through fundamentally different mechanisms. Cascading failures typically unfold as discrete, far-from-equilibrium processes, whereas adversarial interactions affect an existing equilibrium through continuous dynamical effects. 
In this work, we focus on the latter mechanism, namely the destabilization induced by adversarial connections. Although distinct in nature, both edge removals and adversarial interactions can ultimately lead to network instability.




%{Alternatively, network instabilities may originate from the presence of adversarial connections between the nodes of a network \cite{Altafini2013Consensus,Chen2017spectral,ahmadizadeh2017eigenvalues}. {The presence of links with associated negative weights has been reported in brain networks as inhibitory connections between  inhibitory neurons \cite{ahmadizadeh2017eigenvalues}, in social networks as foe relationship between agents \cite{Altafini2013Consensus}, and in distributed control networks as faulty communication processes between systems \cite{Chen2017spectral}.} While both cascading failures \cite{Buldyrev_2010,reis2014avoiding,albert2000error,PhysRevE.66.065102} and instabilities induced by adversarial interactions can destabilize the dynamics of networks, they operate through fundamentally different mechanisms: the former unfold through discrete, far-from-equilibrium sequences, whereas the latter result from the destabilization of  an equilibrium. This paper focuses on the effects of adversarial connections. We emphasize that, although being completely different processes, both edge removals and adversarial connections can lead to network instabilities.} 
%Dynamical approaches, on the other hand, explore how perturbations such as the introduction of a new opinion or the spread of an infection propagate through the network and impact its dynamics \cite{granovetter1978threshold, pastor2001epidemic, castellano2009statistical, centola2010spread}. 


 
\begin{comment}
Attacks against networked systems %, often in the form of cyberattacks, 
are common. 
An example is the Stuxnet worm cyberattack  
which led to the destruction of around 1,000 centrifuges used for uranium enrichment
 of an Iranian  nuclear facility.
The worm manipulated the centrifuges' speeds, causing them to malfunction while sending normal readings to monitoring systems.
%Stuxnet is considered one of the first known cyberattacks to cause physical damage to infrastructure, and it highlighted the vulnerability of critical industrial control systems to cyber intrusions. 
Another example is the 2015 Ukrainian Power Grid Attack. The attackers gained access to the network of an energy company and %coordinated an attack that shut down substations. B
by manipulating the communication protocols between network components, disrupted synchronization between generators and grid operations, leading to large-scale power outages. 
\end{comment}

A fundamental question, with significant societal and technological implications, is  how desired dynamical network states (such as coordinated dynamical states in networks of drones or autonomous vehicles) can be
disrupted by the purposeful insertion of one or a few
{adversarial} agents. %The intruders are structurally similar to the other network agents, which is consistent with the
In particular, this is relevant the 
case of cyber-attacks (CAs)  against
cyber-physical systems (CPSs) and critical infrastructures, such as power grids \cite{pagani2013power}, autonomous vehicle fleets \cite{ren2007information}, and industrial networks \cite{hearnshaw2013complex}. 
Since CAs target the cyber (software) level, while leaving the physical (hardware) level unaltered, they can be effectively modeled as network nodes that have the same structure as the others, but act adversarially.  Because CPSs rely on synchronization and coordination among such agents to ensure stability and proper functionality \cite{zhang2011optimal}, the presence of even a small number of adversarial nodes can directly undermine these mechanisms, rendering CPSs particularly vulnerable to attacks that disrupt synchronization \cite{pasqualetti2013attack}.  
%integrate physical processes with computational control and communication, which is at the core of critical infrastructures like
 %A notable example is that of the WannaCry ransomware attack in May 2017 \cite{kao2018dynamic}, which disrupted industries worldwide—particularly healthcare—causing financial losses, data breaches, and service disruptions. %The growing frequency of cyberattacks exacerbates the risks to multi-agent systems, highlighting the urgency of studying and analyzing network vulnerabilities.
%Complex technological systems, such as coupled autonomous systems, power grids, and sensor networksare distributed cyber-physical systems. 
%Notable examples include the Stuxnet worm cyberattack   which led to the destruction of around 1,000 centrifuges used for uranium enrichment of an Iranian  nuclear facility, and originated from infection of a single machine \cite{buchanan2022cyber}. 
 % \textbf{{Amir: How is the Stuxnet worm relevant to our theory? From what I read in \cite{buchanan2022cyber}, the worm was a  computer software that was operating on Microsoft windows systems and targeted controllers. It shut down the alarms and sent malicious controlling signals to other units. One can say it acted as a smart malicious human operator. If there was a single "bad" centrifuge that was placed in the network of physically coupled centrifuges without any extra controlling power, then mentioning this worm would had relevance. I suggest we remove the reference to a specific worm/virus.}}
%are the 2007 cyber assault on Estonia's government and financial systems, the 2015 cyberattack on Ukraine's power grid \cite{case2016analysis, and the 2021 Colonial Pipeline ransomware attack.  %. Consequently, in the case of a cyberattack, a
%A consequence of the fact that cyber-attacks are only active at the cyber layer is that it may be non-trivial to distinguish adversarial agents from trustworthy agents. Hence, 
Thus an essential step in order to develop resilience against  
cyber-attacks is to understand the mechanisms by which one or a few
intruder agents within %infiltrating
%adversarial nodes intruding 
an otherwise functioning network may compromise its dynamics. Our work also extends to attacks against social networks and biological networks, which we demonstrate using the classical Kuramoto model \cite{KuraBOOK}.
%For example, one may think that the number of intruders needed to destabilize a network is a function of the number of network nodes, $N$. However, in what follows we show that this assumption is incorrect, as only one intruder is always sufficient to destabilize a network, independent of $N$.
%This poses the challenge that  of complex technological systems, such as coupled autonomous systems, power grids, and sensor networks relies on the ability to distinguish adversarial agents from trustworthy agents.
%is the ability to 
%it may become nontrivial 
%A critical problem related to cyberattacks is that they infringe the 
%make it it may not always be possible to 
 %This is especially relevant for the security of complex technological systems, including coupled autonomous systems, power grids, and sensor networks.

{Many real-world networks are highly directed, non-normal, and hierarchically organized, structural features known to strongly shape their dynamical response \cite{johnson2017looplessness,Asllani2018Structure,asllani2018topological,MUOLO2019Patterns,muolo2020synchronization,o2021hierarchical,Duan2022Network,nazerian2023commphys,Ramon, nazerian2024efficiency, muolo2024persistence, nazerian2026frequency}. In such systems, asymmetries in the interaction pattern generate preferred source-to-sink directions of flow, while the network architecture is often close to acyclic, with only sparse or weak feedback loops. As a consequence, perturbations do not need to spread isotropically: they can be transiently amplified, strongly biased downstream, and funneled toward terminal or sink components \cite{o2021hierarchical,Ramon,nazerian2026frequency}. These features naturally call for spectral measures that capture both stability and transient amplification when assessing vulnerability to adversarial attacks.}

{For all the cases of consensus, synchronization, formation control, power grid balance, and Kuramoto-type dynamics, we will see that the most vulnerable components of a network are the low-indegree nodes, {where the indegree of node $i$ is defined as the sum of the weights of the connections directed to node $i$,} as opposed to the hubs on which much previous work has focused \cite{artime2024robustness}. Thus %we conclude that  
while 
{attacks targeting the hubs usually lead to more severe cascading failures, attack strategies that leverage adversarial interactions are more effective when targeting nodes with low indegree. }

The general theme of this paper is  illustrated in Fig.\ \ref{fig:main} through an example. %As we will see in what follows, our work goes beyond the case of consensus dynamics \eqref{consensus} and has implications to a broad variety of network dynamical processes, described by the general set of Eqs.\ \eqref{general}. As an example of this,
Fig.\,\ref{fig:main} a) displays a network of drones attempting to attain a desired formation.  The intruder attack consists in a single intruder node establishing adversarial connections with some of the other nodes, which is shown in b). The bottom panels of Fig.\,\ref{fig:main} are pictorial representations of the spatial trajectories of the drones before and after the attack: in c)  they converge to a set of desired target positions;  in d) they diverge from the target positions, due to the interactions with the intruder.






%{Our analysis encompasses both directed and undirected networks, with a particular focus on transient responses. Directed networks, which are prevalent in real-world systems, are found to exhibit distinct vulnerabilities under such attacks. Universal scaling properties emerge in the destructive power of the intruder, with small perturbations eliciting uniform global responses and larger perturbations exposing specific vulnerabilities, particularly in low-indegree nodes. We specifically demonstrate that similar mechanisms extend to other types of dynamics, including formation control, power grid stability, and synchronization in oscillator networks. }

\begin{comment}

{Besides technological and social systems, intruder attacks can also impact ecological networks, where stable interactions among species or populations are crucial for maintaining functionality and diversity. Both competitive dynamics and targeted disruptions can destabilize these networks, leading to cascading effects. A notable example is the introduction of the American grey squirrel \emph{Sciurus carolinensis} to Europe and the UK. This invasive species, with its larger size, adaptability and superior ability to exploit resources, outcompeted native red squirrels \emph{Sciurus vulgaris}, causing significant population declines \cite{mackinnon1978competition, gurnell2004alien}. This example demonstrates how the introduction of a single adversarial agent can propagate disruptions through an ecological network, reshaping its balance and stability. Such disruptions are not limited to competitive dynamics but also extend to mutualistic networks. %, where the loss of highly connected species can weaken critical interactions and destabilize the entire system. 
Within this context,  Ref.\ \cite{Morone_2018}  investigated the removal of species within critical $k$-cores of mutualistic ecosystems and showed that they can reduce resource availability and disrupt essential relationships between species. If too many species are lost within these cores, the network crosses a tipping point, making recovery impossible and leading to irreversible collapse \cite{Morone_2018}. Competitive and mutualistic disruptions highlight the fragility of ecological networks and emphasize the importance of understanding how adversarial influences can destabilize interactions critical to maintaining ecological resilience.}
\end{comment}



%In this work, we analyze the robustness of network dynamics against intruder attacks, focusing on their efficiency and potential to exploit intrinsic vulnerabilities. This paper introduces the problem of the intruder attack in a complex network and studies conditions for such attack to destabilize the network dynamics. We uncover simple mathematical  conditions for the emergence of instability in the presence of an intruder attack and quantify the level of instability achieved as a result. We first study  
 %the case of the network consensus dynamics and then show how an intruder attack may produce similar effects in other network synchronization dynamics, such as synchronization, formation dynamics, and power grid dynamics.


% Both natural and artificial complex networks evolve in time, through either the addition or removal  of nodes and links. The Internet, the world wide web, the network of scientific collaborations, and... are all examples of continuously growing networks. %, to become what they are now. %This time evolution is ongoing and never ending. 
% Growth is one of the two hallmarks (the other one being preferential attachment) of the 
%  Barabasi-Albert model, the most popular network generative model,  which synthesizes networks with the same power-law degree distribution observed in many real networks. In general, the process by which nodes are added to an existing network can be either deterministic or random or a mix of the two, which will finally determine the emergent structure and properties of the network. One of the issues with growing networks is that some of the newly added nodes may act as `intruders', and attempt to destabilize the network dynamics. We show here that in the case of an intruder attack, a single intruder node connected to the rest of the network by a single connection can destabilize the network dynamics, independent of the size of the network. 

% Many existing biological, technological, and social networks are quite large, as they may be composed of thousands, millions and sometimes billions of nodes. An important question is how a careful insertion of one or a few additional nodes can be used to intentionally impact these networks and possibly transform them. Such precise insertions can either be helpful when they are intended to increase robustness and stability, or adversarial when they are designed to compromise safety and functionality. The latter case has strong implications %relevance  %tightly related to 
% %the fundamental concepts 
% in terms of %assessing 
% network robustness and security, as it is for instance in the case of a network of autonomous vehicles coupled to
% achieve a stable formation, for which a relevant question is 
% whether the purposeful addition of %a single
% one or more adversarial vehicles is sufficient to destabilize the network.
% Within this context, a design question that arises is how to connect the new node(s) to those that are already part of the network, in order to maximize  the negative effects
% %helpful or hurtful impact %consequences 
% of the insertion.
% This is precisely the subject of this paper.


%






%The rest of the paper is organized as follows: Section \ref{sec:undirected} studies the adversarial agent problem for undirected graphs, Section \ref{sec:directed} studies this problem for directed graphs, Section \ref{sec:applications} provided applications of the adversarial agent problem, and conclusions are given in Section \ref{sec:conclusions}. The mathematical details of our results and derivations are available in Section \ref{sec:methods}.



\section{RESULTS}

This paper investigates the case of intruder attacks on networks,
where %the intruder is an adversarial agent that camouflages in a trustworthy agent. % and how this may affect the entire network dynamics. 
%These examples point out to the importance of characterizing how attacks can be perpetrated against networked systems and possible prevention strategies.
%This paper studies networks of coupled dynamical systems under the effect of an intruder attack. 
%Modern biological, technological, and social networks often consist of thousands, millions, or even billions of nodes. This raises the critical question of how the intentional addition of one or a few nodes can impact these networks, by compromising their safety and functionality. %under the effects of what we call an intruder attack 
%A particularly concerning form of attack involves an adversarial agent intruding on the network with the goal of destabilizing its dynamics. 
% Complex networks are ubiquitous components of biological, social, and technological systems. Often the nodes of these networks evolve in time %based on the states of the individual nodes and of their neighbors. A large literature has investigated conditions under which the states of the nodes can 
% to converge on a `well behaved' dynamical state, such as consensus, the agreement state in opinion dynamics, the synchronous evolution in oscillator networks, and  a desired formation in the case of coupled autonomous systems  [refs]. Both the nodes and the connections of a network can be subject to either spontaneous failures or intentional attacks, with many papers investigating the effects of these perturbations in terms of both the network structure [refs] and its dynamics [refs]. Another type of attack occurs when an adversarial agent intrudes a network with the goal of disrupting %a certain well behaved 
% the network dynamics. This case holds much relevance in terms of the security of complex technological networks, such as coupled autonomous systems, power grids, and sensor networks.
%{(\textbf{MA: problem statement?!})} {Modern biological, technological, and social networks often consist of thousands, millions, or even billions of nodes. This raises the critical question of how the intentional addition of one or a few nodes can impact these networks, either by enhancing their robustness and stability or, conversely, by compromising their safety and functionality. The latter scenario has significant implications for network security, such as in autonomous vehicle networks designed to maintain stable formations, where the addition of a single or a few adversarial vehicles could potentially destabilize the system. Within this context, a key design challenge is determining how such nodes should be connected to maximize their destructive impact.}
%{In particular, we analyze the robustness of network dynamics against what we call an `intruder attack'. %, focusing on their efficiency and ability to exploit intrinsic network vulnerabilities. 
our definition of an intruder is a network node that obeys the same individual dynamics  and communicates with the same output function as the other network nodes, but acts maliciously in order to disrupt the network dynamics.  In terms of the mathematical model, this means that an intruder node obeys the same individual (uncoupled) dynamics  and it uses the same output function as the other nodes in the network, but it interacts with the remaining nodes in a purposefully adversarial %destructive 
way. %An analogy could be the following: the human body can be attacked by viruses and bacteria that are inherently different from the cells of the body or by cancer cells which are fundamentally the same as the other cells in the body but act differently. The latter would be the case of an intruder attack.
%Because networks are inherently distributed systems,  each one of their nodes provides a possible entry point for an attack that from there can then propagate to the rest of the nodes.
\begin{comment}{Fundamental questions that arise are how many intruders are needed to destabilize the dynamics of a possibly large network and how the particular position of these intruders over the network can either minimize or maximize the effects of an attack. }
\end{comment}
{In what follows, we mathematically formulate the problem of intruder attacks on networks and derive rigorous conditions to quantify their destabilizing effects on the dynamics. {Although our results extend to both undirected and directed networks, we focus principally on the latter as real-world systems are usually characterized by directional interactions \cite{Asllani2018Structure, o2021hierarchical, Ramon, johnson2020digraphs}.} We first focus on the important case of the linear consensus problem in networks and then consider other types of network dynamics, including formation control, power grid stability, and synchronization in oscillator networks. Our work reveals an intrinsic vulnerability in technological systems, such as autonomous networks, power grids, and sensor systems, emphasizing the need for a deeper understanding of these weaknesses and the development of specific strategies to mitigate them.}

%This section discusses our methodology and results. We briefly outline here our work that follows.  We first introduce a general set of equations that describe the dynamics of the systems that form a network, for which we formulate two types of intruder attacks, with an assigned budget, which we call bidirectional and unidirectional. We then define optimization Problems 1 and 2 that aim to find the most destructive intruder attfack and we prove that for both problems, the optimal attack is one for which the entire budget is spent on one individual network node; we then show by analyzing the attacks to different individual nodes of the network, that the node with lowest indegree is the optimal node to attack. 

\begin{comment}
The outline of this section is as follows: 
% 1- a general set of equations that describe the dynamics of the networked systems are introduced; 2- the two types of budgeted attacks, bidirectional and unidirectional, are formally defined; 3- the notion of algebraic connectivity (a generalization of the Fiedler value) is defined; 4- we define Problems 1 and 2 that aim to find the most destructive attack; 5- we prove that in both problems, all budget must be focused on one connection to one node of the network for optimality; 6- by analyzing attacks to individual nodes of the network, we show that the node with lowest indegree is the optimal node to attack.
\begin{enumerate}[label=\roman*.]
    \item a general set of equations that describe the dynamics of the networked systems are introduced;
    \item the two types of budgeted attacks, bidirectional and unidirectional, are formally defined;
    \item the notion of algebraic connectivity (a generalization of the Fiedler value) is defined;
    \item Problems 1 and 2 are defined that aim to find the optimal attack;
    \item for both problems, Propositions \ref{prop1} and \ref{prop3} prove that the optimal attack is to concentrate all budget on the attack on one node;
    \item by analyzing concentrated attacks on individual nodes of the network in subsections \,\ref{sec:balanced} and \ref{sec:general}, it is shown that the node with the lowest indegree is the optimal node to attack;
    \item finally, the subsection \ref{sec:applications} revisits step i. by providing three numerical examples as applications of the attacked dynamical networks to demonstrate the effect of the attack on wide ranges of dynamics.
\end{enumerate}

\end{comment}



We consider a general set of equations for the network dynamics,
\begin{equation}
    \dot{\bx}_i (t)= \bF(\bx_i (t))-\sum_{j=1}^N L_{ij} \bH(\bx_j (t) - \bbeta_j), \quad i=1,...,N, \label{general}
\end{equation}
where  $\bx_i(t)$ is the $m$-dimensional state of node $i$ at time $t$ and $N$ is the number of nodes. The individual nodal dynamics is given by $\bF(\bx_i(t))$, and the output of node $j$ is given by the function $\bH(\bx_j(t) - \bbeta_j)$ where $\bbeta_j$ is an $m$-dimensional constant vector that acts as a coordinate shift ($\bbeta_j$ may be zero). {The network connectivity is described by a digraph with adjacency matrix $A=[A_{ij}]$, where  $A_{ij}$ is the strength of the directed coupling from node $j$ to node $i \neq j$, $A_{ij}>0$ indicates a cooperative interaction (also sometimes called mutualistic or attractive), $A_{ij}<0$ indicates an adversarial interaction (also sometimes called competitive or repulsive), and $A_{ij}=0$ indicates no interaction.} The indegree of node $i$ is equal to $d_i^{in}=\sum_{j \neq i} A_{ij}$ and the outdegree of node $i$ is equal to $d_i^{out}=\sum_{j \neq i} A_{ji}$. The Laplacian matrix is denoted by $L=[L_{ij}]$ where $L_{ij}= (\delta_{ij} d_i^{in} - A_{ij})$, and $\delta_{ij}$ is the Kronecker delta. By construction, $\sum_{j=1}^N L_{ij} = 0, \forall i$. {For clarity, we note that two sign conventions for the Laplacian matrix are used in the literature: one with off-diagonal entries equal to the adjacency matrix and diagonal entries equal to the negative indegrees, and one with off-diagonal entries equal to the negative adjacency matrix and diagonal entries equal to the indegrees. In this work, we adopt the latter convention.}
We proceed under the assumption that the Laplacian matrix $L$ is `proper', i.e., all the eigenvalues of  $L$ have non-negative real parts and there is only one eigenvalue equal to $0$.
Conditions for the Laplacian matrix of a signed digraph to be proper have been investigated in the literature \cite{Altafini2013Consensus,Chen2017spectral,ahmadizadeh2017eigenvalues}.


Reference \cite[Definition 2.32]{wu2007synchronization} defines, for a given Laplacian matrix $L$, the following algebraic connectivity $f \in \mathbb{R}$ which may be considered as a generalization of the Fiedler value \cite{fiedler1973algebraic}  for digraphs:
\begin{equation} \label{eq:algeb}
    f(L) := \min_{\bX \neq \pmb{0}, \bX \perp \pmb{1}} \dfrac{\bX^\top L \bX}{\bX^\top \bX} = \lambda_{\min} \left( V^\top \dfrac{L + L^\top}{2} V \right).
\end{equation}


Here, $V \in \mathbb{R}^{N \times N-1}$ is an orthonormal basis for the null subspace of $\bone^\top$, i.e., $V$ is a matrix whose columns are normal and orthogonal to one another and have zero column sums, and $\lambda_{\min}(S)$ indicates the smallest eigenvalue of the symmetric matrix $S$.
Note that if $L$ is symmetric (the graph is undirected), the {Fiedler value of the Laplacian $\alpha = f \geq 0$,} where the strict inequality is achieved if and only if the graph is connected.
    The algebraic connectivity $f$ is relevant to studying the synchronization and consensus stabilities of directionally coupled networks. 
    Also, $f$ appears in studying the contractivity of such dynamics and their respective Lyapunov functions, e.g., see \cite[Corollary 4.23]{wu2007synchronization}. 
    Note that for $f<0$, $-{f}$ measures the time scale of the unstable dynamics.







We consider the presence of an adversarial agent, labeled as node $N+1$, and that this agent is described by the same functions $\bF$ and $\bH$ as the remaining nodes.
As a result, the network equations can be written,
\begin{equation}
    \dot{\bx}_i (t)= \bF(\bx_i (t))-\sum_{j=1}^{N+1} {L_{aug}}_{ij} \bH(\bx_j (t) - \bbeta_j),
\end{equation}
$i=1,...,N+1$, where the augmented Laplacian matrix $L_{aug}$ is either equal to
\begin{equation} \label{eq:Lnewbi}
    L_{aug}^b = \begin{bmatrix}
        L + \text{diag}(\bb) & -\bb \\
        -\bb^\top & -c
    \end{bmatrix}
\end{equation}
in the case that the intruder is bidirectionally coupled to the network nodes or equal to
\begin{equation} \label{eq:Lnewuni}
    L_{aug}^u = \begin{bmatrix}
        L + \text{diag}(\bb) & -\bb \\
        \pmb{0}^\top & 0
    \end{bmatrix}
\end{equation}
in the case that the intruder is unidirectionally coupled to the network nodes.
%\begin{equation} \label{eq:Lnew}
%    L_{aug} = \begin{bmatrix}
%        c & -\bd^\top \\
%        -\bb & L + \text{diag}(\bb)
%    \end{bmatrix},
%\end{equation}
In both \eqref{eq:Lnewbi} and \eqref{eq:Lnewuni} the vector $\bb=[b_i \leq 0]$, $\bb \neq \textbf{0}$ (i.e., at least one $b_i<0$) contains the weights with which the additional node connects to nodes $1,...,N$. {A negative (zero) weight $b_i$ indicates the presence (absence) of an adversarial connection between the intruder and node $i$.  Furthermore, we introduce the attack budget, defined as the positive quantity $c=-\sum_{i=1}^N b_i${, which is taken to be equal to the sum of the negatives of the weights of all the connections established by the intruder with the existing network nodes.}} %the vector \textcolor{red}{$\bd=[d_i \leq 0] $} contains the weights with which nodes $1,...,N$ connect to the additional node and $c=\sum_{i=1}^N d_i $.





\begin{figure*}
    \centering
    %\includegraphics[width=0.7\linewidth]{figs/Drones.png}
    \includegraphics[width=0.8\linewidth]{figs/Illustration_Intruder3.png}
    \caption{\textbf{Illustration of an intruder attack on a network of drones attaining a given formation.} The top panels show a network of coupled drones, before (A) and after (B) an intruder attack. The bottom panels display the trajectories of the drones before the attack in (C), showing convergence to the target positions, and after the attack in (D), showing divergence from the target position. The target positions are represented as stars in (C) and (D).}
    \label{fig:main}
\end{figure*}

{
A foundamental realization of the general dynamics \eqref{general}, which we will analyze in detail in what follows, is that of the single-integrator consensus dynamics, described by the equation,
\begin{equation} \label{consensus}
    \dot{\bX} (t) = -L \bX(t),
\end{equation}
in the $N$-dimensional state vector $\bX=[x_1,x_2,...,x_N]$ ($x_i$ here is a scalar). When the Laplacian matrix is proper, Eq.\ \eqref{consensus}
reaches consensus asymptotically independent of its initial condition, i.e., $\lim_{t \rightarrow \infty} x_i(t) = \bar{x}$, $i=1,...,N$, where $\bar{x}$ is the consensus state \cite{olfati2007consensus}. 
%Consensus dynamics refers to the process by which agents in a network adjust their opinions or states through Laplacian interactions, in order to converge toward a common opinion .
In Supplementary Note 1, we provide an example of single-integrator consensus dynamics \eqref{consensus} on a network subject to a single intruder attack to further motivate the relevance of the algebraic connectivity $f$. In this example, the intruder unidirectionally connects to different network nodes, resulting in different values of the algebraic connectivity $f$ as the Laplacian matrix $L_{aug}^u$ changes.
In all cases, we see that after the intruder connects to one of the other nodes, the individual trajectories $x_i (t)$ diverge from one another.


{
We characterize consensus in terms of the time evolution of the norm of the perturbations about the consensus state 
\begin{equation}
    \| \delta \bX(t) \| := \sqrt{\sum_{j=1}^{N} \left(x_j (t) - \frac{1}{N} \sum_{k=1}^{N} x_k (t)  \right)^2}. 
\end{equation} Consensus is achieved if $\|\delta \mathbf{X}(t)\| \to 0$ as $t \to \infty$. }

\begin{definition}\textit{{Destabilization of the single-integrator consensus dynamics.}}
By construction the Laplacian matrix has one eigenvalue equal zero. The consensus dynamics (6) is: (i) stable iff all the remaining eigenvalues have positive real part and (ii) unstable iff at least one of the remaining eigenvalues has negative real part. Consequently, we call a destabilization a transition from (i) to (ii).  
\end{definition}
\color{black}
We take the original Laplacian matrix $L$ to be proper, which guarantees that consensus is achieved prior to the attack. However, the addition of a single adversarial connection with the intruder implies that the resulting augmented Laplacian matrices $L_{aug}^b$ and $L_{aug}^u$ are no longer proper, as they are known to possess at least one eigenvalue with negative real part (see e.g., \cite{Chen2017spectral,ahmadizadeh2017eigenvalues}) and, as a consequence, consensus is not achieved. 

%the consensus dynamics becomes unstable and perturbations grow in time. We refer to this transition from a stable regime (before the adversarial connection) to an unstable one (after its introduction) as a \emph{destabilization}.



While the the addition of a single adversarial connection with the intruder produces a destabilization, in this paper we focus not on $\alpha(L_{aug})$, the asymptotic rate of instability for $\| \delta \bX(t) \|$, but rather on its transient rate of instability $|f(L_{aug})| \geq |\alpha(L_{aug})|$ immediately following the attack. As detailed in Subsection A of the Methods, a larger transient growth arises from the non-normal dynamics  characteristic of directed networks \cite{Asllani2018Structure, trefethen1991pseudospectra}. We emphasize the transient rate for three main reasons: (i) when Eq.~(1) is obtained via linearization about an unstable synchronous solution \cite{nazerian2024efficiency}, it accurately describes the dynamics only shortly after the perturbation, rendering the asymptotic rate less relevant; (ii) the transient rate provides a worst-case estimate, as it is greater than or equal to the asymptotic rate; and (iii) larger transient growth increases the likelihood that a perturbation produces a significant alteration of the system dynamics \cite{asllani2018topological, Asllani2018Structure, MUOLO2019Patterns, muolo2020synchronization, nazerian2024efficiency}.

%In our setting, when the adversarial connection is absent, perturbations decay and the network dynamics are stable. When the adversarial connection is introduced, perturbations instead grow, leading to instability. We refer to the transition from a stable to an unstable regime induced by the addition of an intruder node as \emph{a destabilization}.

{%In what follows, we quantify this effect through the growth rate of perturbations, which is directly related to the algebraic connectivity $f$.
%%Also, the norm of the perturbations 
%\begin{equation}
%    \| \delta \bX(t) \| := \sqrt{\sum_{j=1}^{N+1} \left(x_j (t) - \frac{1}{N+1} \sum_{k=1}^{N+1} x_k (t)  \right)^2} 
%\end{equation}
%decays when the adversarial connection is off ({which is the case of stable dynamics}) and grows when the adversarial connection is on  ({which is the case of unstable dynamics}.)  {In what follows we are particularly interested in the transition from stable dynamics to unstable dynamics, a process which we call destabilization, as a result of adding a single intruder node to a stable network. The rate of instability, i.e., the rate at which $ \| \delta \bX(t) \|$ grows directly correlates with the algebraic connectivity $f$.}
Supplementary Note 1 also reviews some of the relevant properties of the algebraic connectivity $f$ from \cite{wu2007synchronization}.}}




\begin{comment}
    
{
For both the cases of the augmented Laplacian matrices $L_{aug}^b$ and $L_{aug}^u$,
it is known (see e.g., \cite{ahmadizadeh2017eigenvalues} ) that they have at least one negative real part eigenvalue, i.e., they are not proper, which results in asymptotic instability of the consensus dynamics \eqref{consensus}. In this paper, we use this result; however rather than being interested in the asymptotic rate of instability $\alpha( L_{aug})$, 
we focus on the transient rate of instability $f( L_{aug}) \geq \alpha( L_{aug})$ that is produced by the addition of the adversarial node immediately after the attack. Subsection A of the Methods exactly defines and compares these two rates of instability, where
{a larger transient growth arises from the non-normal dynamics characteristic of directed networks %as quantified by the \emph{numerical abscissa}, which represents the maximum possible instantaneous growth rate of a perturbation 
\cite{Asllani2018Structure, trefethen1991pseudospectra}.} There are three main reasons why we decide to focus on the transient rate of instability: (i)  in the case in which the dynamics (1) is linearized about an unstable synchronous solution \cite{nazerian2024efficiency}, the linearization is only a valid representation of the dynamics immediately after the attack, thus making the asymptotic rate of growth of a perturbation from the linearized dynamics irrelevant; (ii) the transient rate of growth is larger than or equal to the asymptotic rate of growth, thus it provides a `worst-case' rate of growth; (iii) the higher is the transient growth, the 
%The importance of transient growth in the dynamics of non-normal linearly stable systems lies in its ability to estimate the
higher is the likelihood that a perturbation will cause a significant alteration of the system dynamics %escape the original attractor 
\cite{asllani2018topological, Asllani2018Structure, MUOLO2019Patterns, muolo2020synchronization, nazerian2024efficiency}.}

\end{comment}




{
In principle, one could consider the presence of more than one intruder; however, we show in what follows that under appropriate conditions, a single attacker is able to destabilize the network dynamics, so we focus on the case of a single attacker.} {The case of multiple intruders is studied in Supplementary Note 2.}



\begin{comment}

\section{Adversarial Agent: Undirected Graphs} \label{sec:undirected}


The traditional study of steady state behavior of the system usually is performed through the eigenvalues of the linearized dynamical system which measure the asymptotic rate of growth or decay of the states of the system.
Although a system may be stable, it can show instability characteristics in transient.
We put a higher emphasis on short-term (transient) stability since natural and artificial networked systems are often nonlinear and the steady-state stability analysis is only valid locally. 
The local approach has a drawback in that if the system leaves the linear regime of the solution, the long-term stability is no longer predicted via the linearization-based methods.
Thus, it is vital to ensure that not only the system is stable in long-term, but also stable in short-term as well where the latter is quantified via reactivity.
If the system remains non-reactive in the short term, the local stability approach is guaranteed to remain valid.
There are abundant studies that relate the concept of reactivity to networked systems [refs].
In simple terms, the reactivity measures the instantaneous rate of instability (the instantaneous rate of growth or decay of the norm of the state vector.)
The question that interests us is how the addition of an adversarial agent affects the short-term stability (reactivity) of the system.
We look into the cases of directed and undirected graphs and study the adversarial agent problem.



In this section, we assume that a new agent (node) is being added to the existing network with a difference that the sign of the connections between the new agent (node) and the existing agents (nodes) are different from that of among the existing agents (nodes). We assume the sum of the weights of the new connections is a given budget.

To motivate the importance of studying the adversarial agent problem, we demonstrate the effect of having mixed signed eigenvalues of the Laplacian matrix via examples of single-integrator consensus dynamics and complete synchronization of coupled chaotic oscillators.

\textbf{Example:} The Consensus Dynamics.
The state of the system $i$ at time $t$ is denoted as $x_i(t)$ and $\dot{x}_i (t) = - \sum_{j=1}^N L_{ij} x_j$,
% \begin{equation}
%     \dot{x}_i (t) = - \sum_{j=1}^N L_{ij} x_j,
% \end{equation}
where the matrix $L = [L_{ij}]$ is a Laplacian matrix, i.e., such that $\sum_j L_{ij}=0$.
In compact form, 
\begin{equation} \label{eq:consensus}
    \dot{\bx} (t) = - L \bx(t),
\end{equation}
where $\bx (t) = [x_1 (t) \ \ x_2 (t) \ \cdots \ x_N (t)]^\top$.

We randomly generate the Laplacian matrix $L$ and add an adversarial node with budget $c = -1$ to the node with the minimum degree to obtain $L_{aug}$:
\begin{align}\label{eq:Lconsensus}
\begin{split}
    L = & \begin{bmatrix}
    1 & 0 & -1 & 0 & 0 \\ 
    0 & 2 & -1 & 0 & -1 \\ 
    -1 & -1 & 3 & -1 & 0 \\ 
    0 & 0 & -1 & 1 & 0 \\ 
    0 & -1 & 0 & 0 & 1 
    \end{bmatrix}, \\
    L_{aug} = &  \begin{bmatrix}
    -1 & 1 & 0 & 0 & 0 & 0 \\ 
    1 & 0 & 0 & -1 & 0 & 0 \\ 
    0 & 0 & 2 & -1 & 0 & -1 \\ 
    0 & -1 & -1 & 3 & -1 & 0 \\ 
    0 & 0 & 0 & -1 & 1 & 0 \\ 
    0 & 0 & -1 & 0 & 0 & 1 
\end{bmatrix}.
\end{split}
\end{align} 
The set of the eigenvalues of the matrices $L$ and $L_{aug}$ are $\Lambda = \{0, 0.51, 1, 2.31, 4.17\}$ and $\Lambda_{new} = \{ -1.69, 0, 0.45, 0.83, 2.28, 4.12 \}$, respectively. 
We simulate the consensus dynamics in Eq.\,\eqref{eq:consensus} for both Laplacian matrices $L$ and $L_{aug}$. 
The initial conditions are chosen randomly with zero mean as $\bx(0) = [-6.27, 0.75, 0.04, 2.69, 2.77]$ for the case with $L$ and $\bx(0) = [0, -6.2735, 0.7521, 0.0452, 2.6989, 2.7773]$ for the case with $L_{aug}$.

Figure \ref{fig:consensus}A shows the consensus dynamics with the Laplacian matrix $L$ is stable, while panel B shows the instability of the dynamics with the Laplacian matrix $L_{aug}$.
This difference in stability resulted from having mixed sign eigenvalues of the Laplacian matrix $L_{aug}$ due to the addition of the adversarial agent.

\begin{figure}
    \centering
    \text{A) Original consensus} \hspace{0.15\linewidth} \text{B) Adversarial agent added} \\
    \includegraphics[width=0.48\linewidth]{figs/L.pdf} 
    \hspace{0.01\linewidth}
    \includegraphics[width=0.48\linewidth]{figs/Lnew.pdf}
    \caption{The consensus dynamics in Eq.\,\eqref{eq:consensus} for the two Laplacian matrices $L$ and $L_{aug}$ from Eq.\,\eqref{eq:Lconsensus}. Panels A and B show the state trajectories for each system with the Laplacian matrices $L$ (original consensus) and $L_{aug}$ (adversarial agent added), respectively. The dashed black line shows the trajectory associated with the newly added adversarial agent.}
    \label{fig:consensus}
\end{figure}

% Next, we study the effect of $\lambda_{\min} < 0$, the smallest eigenvalue of the Laplacian matrix $L_{aug}$ from Eq.\,\eqref{eq:Lnew}, over the consensus dynamics.
% We define the instability time $T$ that measures the time it takes for the consensus dynamics to reach a given threshold $\Delta$, i.e., $\exp(-\lambda_{\min} T)= \Delta$, where $\Delta$ is the threshold and $T$ is the instability time. Thus, 
% \begin{equation}
%     T=\dfrac{\log(\Delta)}{-\lambda_{\min}}.
% \end{equation}




Now we formally introduce the adversarial agent problem:

\textit{Problem Statement 1:} Given the $N$-dimentional symmetric Laplacian matrix $L = L^\top \succeq 0$ and a budget $c$, find the weights $\bb = [b_i]$, $b_i \leq 0$, $i = 1, \hdots, N$, with $\sum_i b_i = c \leq 0$ such that the minimum eigenvalue of the matrix 
\begin{equation} \label{eq:Lnew}
    L_{aug} = \begin{bmatrix}
        c & -\bb^\top \\
        -\bb & L + \text{diag}(\bb)
    \end{bmatrix}
\end{equation}
is minimized.




\begin{proposition} \label{prop1}
    Given $c \leq 0$, the minimum eigenvalue of the matrix $L_{aug}$, $\lambda_{\min} (L_{aug})$, is bounded as $2c\leq \lambda_{\min} (L_{aug}) \leq c$.
\end{proposition}
\begin{proof}
    See Sec.\,\ref{sec:prop1}.
\end{proof}

\begin{proposition} \label{prop2}
    Given $c \leq 0$, the minimum eigenvalue of the matrix $L_{aug}$, $\lambda_{\min} (L_{aug})$, is achieved when $b_{i^*} = c$ and $b_i = 0, \ \forall i \neq i^*$ where $i^* = \text{arg}\min_{i} L_{ii}$ (dictatorial). Also, given $c \leq 0$, the maximum of $\lambda_{\min}(L_{aug})$ is achieved when $b_i = c/N, \ \forall i$ (democratic).
\end{proposition}
\begin{proof}
    See Sec.\,\ref{sec:prop2}.
\end{proof}

Next, we analyze the effect of varying the budget $-b \leq 0$ when the connection from the adversarial node is only made to one of the current nodes of the graph.
We implemented matrix perturbation theory with a limit of very small and very large $b$ and obtained:
\begin{itemize}
    \item for small budget $b$: the eigenvalue $\lambda_{\min}(L_{aug}) = -b$, and
    \item for large budget $b$: the eigenvalue $\lambda_{\min}(L_{aug}) = -2b + L_{ii}/2$, and
\end{itemize}
For details of the derivations see Sec.\,\ref{sec:undirect}.
We have summarized the results of this section in Fig. \ref{fig:schematic}.

\begin{figure}
    \centering
    \includegraphics[width=0.6\linewidth]{figs/Schematic_adverserial.pdf}
    \caption{Schematic showing $\lambda_{\min} (L_{aug})$, the minimum eigenvalue of the matrix $L_{aug}$ in \eqref{eq:Lnew}, when only one node (either node $i$, $j$, or $k$) is pinned with the budget $-b < 0$. The dashed black line shows the line $\lambda_{\min} (L_{aug}) = -b$. The colored dashed lines show $\lambda_{\min} (L_{aug}) = -2b + L_{xx}/2$ for $x = i, j,k$.  The schematic is not drawn in scale to enhance visualization.}
    \label{fig:schematic}
\end{figure}



    
\end{comment}






%In this section, we focus on directed graphs, for which, the Laplacian matrix is possibly asymmetric. 
%For directed graphs an another definition of connectivity is needed, alternative to the Fielder eigenvalue for undirected graphs.







\begin{comment}
To motivate the relevance of the algebraic connectivity $f$, we have the following single-integrator consensus dynamics example.
We study the effect of unidirectional agent addition in consensus dynamics through the change in the algebraic connectivity $f$. 

Given a graph with Laplacian matrix $L$, we set
\begin{equation*}
    \tilde{L}_{off} = \begin{bmatrix}
        0 & \pmb{0}^\top\\
        \pmb{0} & L
    \end{bmatrix}, \quad \tilde{L}_{on} = \begin{bmatrix}
        0 & \pmb{0}^\top \\
        -\bb_{i^*} & L + \diag (\bb_{i^*})
    \end{bmatrix}.
\end{equation*}
The Laplacian matrix 
$\tilde{L}_{off}$ describes a graph in which an adversarial agent is added to a network, but no connection between this agent and the original network is made.
The Laplacian matrix 
$\tilde{L}_{off}$ describes the graph after the unidirectional connection from the adversarial agent to the original network is made.
The vector $\bb_{i^*}$ has all zero entries except for the $i^*$th, which equals $-1$.

Consider the consensus problem
\begin{equation} \label{eq:consensusadver}
    \dot{\bx} (t) = - \tilde{L}(t) \bx(t), \quad \tilde{L}(t) = \begin{cases}
        \tilde{L}_{off} \quad  & 0\leq t < T, \\
        \tilde{L}_{on} \quad & t \geq T.
    \end{cases}
\end{equation}
in which the effective Laplacian matrix $\tilde{L}$ switches from $\tilde{L}_{off}$ to $\tilde{L}_{on}$ at $t = T$, i.e., the unidirectional connection from the adversarial agent to the original network is made at time $T$.
This indicates for $0\leq t < T$, the consensus dynamics is stable and for $t \geq T$, it becomes unstable.

Figure \ref{fig:consensusadver} shows the effect of adversarial agent addition to a directed tree of 7 nodes. 
It is clear that the adversarial agent destroys the global consensus since the individual trajectories $x_i (t)$ do not coincide after the unidirectional connection from the adversarial agent is turned on.
Also, the norm of the perturbations 
\begin{equation}
    \| \delta \bx(t) \| := \sqrt{\sum_j \left(x_j (t) - \frac{1}{N} \sum_k x_k (t)  \right)^2} 
\end{equation}
decays when the adversarial connection is off and grows when the adversarial connection is on.
The rate at which $ \| \delta \bx(t) \|$ grows directly correlates with the algebraic connectivity $f$.

\begin{figure*}
    \centering
    \includegraphics[width=0.9\linewidth]{figs/Fig.pdf}
    \caption{
    The consensus dynamics as the adversarial agent (black node) is connected to different nodes in the original network.
    Plots show the individual trajectories of the nodes $x_i (t)$ and the norm of the perturbations about the consensus state $\| \delta \bx(t) \|$. The dashed black line shows the instance in which the unidirectional connection from the adversarial agent to the original network is made, i.e., the line shows the instance $t = T$ in Eq.\,\eqref{eq:consensusadver}.
    The algebraic connectivity $f$ is shown for the time the adversarial connection is turned on.
    The consensus dynamics and the algebraic connectivity $f$ for pinning nodes 6 and 7 are not shown since they are identical to node 5 due to the symmetry of the problem.
    }
    \label{fig:consensusadver}
\end{figure*}
\end{comment}



Next, we introduce our main two problems, which we call attack via bidirectional connections and attack via unidirectional connections. 





{Problem 1 Statement: Attack via bidirectional connections.} Given the $N$-dimensional possibly asymmetric Laplacian matrix $L$, and a budget $-c$, find the weights $\bb = [b_i]$, $b_i \leq 0$, $i = 1, \hdots, N$, with $\sum_i b_i = -c \leq 0$ such that the algebraic connectivity $f$ of the new graph with Laplacian matrix $L_{aug}^b$ in Eq.\,\eqref{eq:Lnewbi} is minimized.
% \begin{proposition} \label{prop1}
%     Given $c \leq 0$ and the Laplacian matrix $L_{aug}$ in \eqref{eq:Lnewbi}, the minimum algebraic connectivity $f$ is achieved when $b_{i^*} = c$ and $b_i = 0, \ \forall i \neq i^*$ where $i^* = \text{arg}\min_{i} L_{ii}$. {\textbf{Incomplete statement.} If two different nodes $i$ and $j$ have the same indegree, $L_{ii} = L_{jj}$, the $f$ may not be the same for pinnig either node $i$ or $j$. However, the budget must still be focused on one single connection to either $i$ or $j$ for optimality.}
% \end{proposition}
% \begin{proof}
%     {No proof at the moment.}
% \end{proof}


\begin{proposition} \label{prop1}
    Given $-c \leq 0$ and the Laplacian matrix $L_{aug}^b$ in Eq.\,\eqref{eq:Lnewbi}, the minimum algebraic connectivity $f$ is achieved when $b_{i^*} = -c$ for one node $i^*$ and $b_i = 0$ for all other nodes $i \neq i^*$, i.e., the entire budget must be allocated to one node.
    % where $i^* = \text{arg}\min_{i} L_{ii}$.
\end{proposition}
\begin{proof}
    See Methods \ref{sec:prop1}.
\end{proof}




\begin{proposition} \label{prop2}
    Given a budget $-c \leq 0$ in Problem 1, the algebraic connectivity is $f \leq 0$.
\end{proposition}
\begin{proof}
    See Methods \ref{sec:prop2}.
\end{proof}




{Problem 2 Statement: Attack via unidirectional connections.} Given the $N$-dimensional possibly asymmetric Laplacian matrix $L$, and a budget $-c$, find the weights $\bb = [b_i]$, $b_i \leq 0$, $i = 1, \hdots, N$, with $\sum_i b_i = -c \leq 0$ such that the algebraic connectivity $f$ of the new graph with the Laplacian matrix $L_{aug}^u$ in Eq.\,\eqref{eq:Lnewuni} is minimized.
% \begin{proposition} \label{prop3}
%     Given $c \leq 0$ and the Laplacian matrix $L_{aug}$ in \eqref{eq:Lnewuni}, the minimum algebraic connectivity $f$ is achieved when $b_{i^*} = c$ and $b_i = 0, \ \forall i \neq i^*$ where $i^* = \text{arg}\min_{i} L_{ii}$. {\textbf{Incomplete statement.} If two different nodes $i$ and $j$ have the same indegree, $L_{ii} = L_{jj}$, the $f$ may not be the same for pinnig either node $i$ or $j$. However, the budget must still be focused on one single connection to either $i$ or $j$ for optimality.}
% \end{proposition}
% \begin{proof}
%     {No proof at the moment.}
% \end{proof}



\begin{proposition} \label{prop3}
    Given $-c \leq 0$ and the Laplacian matrix $L_{aug}^u$ in Eq.\,\eqref{eq:Lnewuni}, the minimum algebraic connectivity $f$ is achieved when $b_{i^*} = -c$ for one node $i^*$ and $b_i = 0$ for all other nodes $i \neq i^*$, i.e., the entire budget must be allocated to one node.
    % where $i^* = \text{arg}\min_{i} L_{ii}$.
\end{proposition}
\begin{proof}
    See Methods \ref{sec:prop3}.
\end{proof}









\begin{proposition} \label{prop4}
    Given a budget $-c \leq 0$ in Problem 2, the algebraic connectivity is $f \leq 0$.
\end{proposition}
\begin{proof}
    See Methods \ref{sec:prop4}.
\end{proof}

{Propositions \ref{prop1} and \ref{prop3} are used in the next subsections \,\ref{sec:balanced} and \ref{sec:general} to characterize the effect of the particular selection of the network nodes to which the intruder should connect, $i=1,..,N$, in order to maximize the severity of the attack.}


Supplementary Note 3 discusses the case of bidirectional attacks in undirected networks with symmetric Laplacian matrices for which we also prove that the optimal attack is to allocate the entire budget to one node.




\begin{comment}
\textbf{Amir: I prefer not to keep the following blue information, but if all think it is better to keep, we can certainly keep it.}

Specifically, given the initial Laplacian matrix $L$ and budget $-c <0$, we aim to see the effect of node $i$ under the attack via perturbing the network described by the Laplacian matrix
\begin{equation*}
    \tilde{L} = \begin{bmatrix}
        L & \pmb{0} \\
        \pmb{0}^\top & 0
    \end{bmatrix},
\end{equation*}
via the addition of the following attack matrices:
\begin{gather*} 
\small
    P_i^b = 
    \begin{bmatrix}
    0 & \hdots & 0 & 0 & 0 & \hdots & 0 & 0 \\ 
    \vdots & \ddots & \vdots & \vdots & \vdots & \ddots & \vdots & \vdots \\ 
    0 & \hdots & 0 & 0 & 0 & \hdots & 0 & 0 \\ 
    0 & \hdots & 0 & -c & 0 & \hdots & 0 & c \\ 
    0 & \hdots & 0 & 0 & 0 & \hdots & 0 & 0 \\ 
    \vdots & \ddots & \vdots & \vdots & \vdots & \ddots & \vdots & \vdots \\ 
    0 & \hdots & 0 & 0 & 0 & \hdots & 0 & 0 \\ 
    0 & \hdots & 0 & c & 0 & \hdots & 0 & -c
    \end{bmatrix}, \\
    \small
    P_i^u = 
    \begin{bmatrix}
    0 & \hdots & 0 & 0 & 0 & \hdots & 0 & 0 \\ 
    \vdots & \ddots & \vdots & \vdots & \vdots & \ddots & \vdots & \vdots \\ 
    0 & \hdots & 0 & 0 & 0 & \hdots & 0 & 0 \\ 
    0 & \hdots & 0 &  -c & 0 & \hdots & 0 & \, \ c \ \, \\ 
    0 & \hdots & 0 & 0 & 0 & \hdots & 0 & 0 \\ 
    \vdots & \ddots & \vdots & \vdots & \vdots & \ddots & \vdots & \vdots \\ 
    0 & \hdots & 0 & 0 & 0 & \hdots & 0 & 0 \\ 
    0 & \hdots & 0 & 0 & 0 & \hdots & 0 & 0
    \end{bmatrix}.
\end{gather*}
The matrices $P_i^b, P_i^u \in \mathbb{R}^{N+1 \times N+1}$ and the entry $(i,i)$, $i = 1, \hdots, N$, in these matrices is equal to $-c$.
The bidirectional (unidirectional) attack on node $i$ is characterized by the matrix $P_i^b$ ($P_i^u$) such that the attacked Laplacian matrix becomes $L_{aug}^b = \tilde{L} + P_i^b$ ($L_{aug}^u = \tilde{L} + P_i^u$.)
We will see that the node $i^*$ with the lowest indegree, i.e., $i^* = \arg \min_i L_{ii}$, is the optimal node to attack.

\end{comment}



\subsection{Balanced Directed Graphs} \label{sec:balanced}

Next, we investigate in detail the case of interest that the directed graph is balanced, i.e., for each one of its nodes, the indegree equals the outdegree, $d_i=d_i^{in}=d_i^{out}, i=1,..,N$ ($d_i$ is simply the degree of node $i$.) Undirected graphs form a particular class within the broader class of balanced directed graphs. %Although this case is rather rare to observe in real networks, it offers some valuable insight.
{
Based on Propositions \ref{prop1} and \ref{prop3}, we know that the optimal attack is to allocate the entire budget on one connection only. %For easier demonstration, we denote the budget as $-c \leq 0$ with $c \geq 0$.
In what follows, we study the attack to node $i$ of a given network, i.e., we set $\bb = [b_j], b_i = -c, b_{j\neq i} = 0$.
}
We are interested in how the algebraic connectivity $f$ of a balanced directed graphs is affected by an intruder attack. 
%The indegree and outdegree of each node in a balanced graph are equal.
By using matrix perturbation theory, we study {Problem 1} in the limit of very small and very large budget $c$ and obtain that:
\begin{itemize}
    \item for small budget $c$: the algebraic connectivity $f(L_{aug}^b) = -1.1c$, and
    \item for large budget $c$: the algebraic connectivity $f(L_{aug}^b) = -2c + L_{ii}/2$.
\end{itemize}
For details on the derivations see Sec.\,\ref{sec:uni}.
Note the dependence on the degree of  node $i$, $d_i=L_{ii}$.

%\subsection{Balanced Directed Graphs: Bidirectional Connection}

Using a similar approach, we can also study {Problem 2} and obtain that:

%We implemented matrix perturbation theory with a limit of very small and very large $b$ and obtained:
\begin{itemize}
    \item for small budget $c$: the algebraic connectivity $f(L_{aug}^u) = -0.1c$, and
    \item for large budget $c$: the algebraic connectivity $f(L_{aug}^u) = -\frac{\sqrt{2}+1}{2}c + (3+2\sqrt{2})  L_{ii}$.
\end{itemize}
For details of the derivations see Sec.\,\ref{sec:bi}. Also in this case, note the dependence on the degree of  node $i$, $d_i=L_{ii}$.


% The above results are summarized in Fig.\,\ref{fig:directed}.



Figure\,\ref{fig:directed} summarizes the results of this section, by providing a schematics for how the algebraic connectivity $f$ varies with the magnitude of the budget $c$. 
The top (bottom) panel is for the case of Problem 1 (Problem 2.) 
{The figure shows that our calculations in the small and large budget limits can be used to sketch the dependence of the algebraic connectivity on $c$. In both panels, it is shown that the slopes in the small- and large-c limits are constant, while the intercept of the large-$c$ asymptote increases with the node degree. This indicates that the algebraic connectivity is minimized when the node degree is minimal.} {Let us further remark that the above result should be interpreted in the context of the adopted Laplacian. More precisely, in the appropriate cost limit, the optimal node ranking is determined by the diagonal entries $L_{ii}$ of the adopted coupling operator. For example, for the symmetrically normalized Laplacian \cite{chung1997spectral}, $L_{ii}=1$ for all $i$, and therefore this distinction between nodes vanishes.}


%For the standard graph Laplacian associated with diffusive coupling considered throughout this work, $L_{ii}=-d_i$, implying that nodes with smaller degree are the most vulnerable. More generally, whenever the adopted coupling operator has identical diagonal entries, either through the choice of edge weights or the Laplacian definition (e.g., the random-walk Laplacian), this node distinction disappears. In particular, for the random-walk Laplacian, all diagonal entries are identical (in our sign convention, $L_{ii}=-1$).}

%an attack via a  bidirectional (unidirectional) connection.




\begin{figure}
    \centering
    \includegraphics[width=0.8\linewidth]{figs/Fig2.pdf}
    \caption{\textbf{Schematic dependence of the algebraic connectivity on the attack budget.}
Panels A and B schematically illustrate the algebraic connectivity, $f(L_{aug}^b)$ and $f(L_{aug}^u)$, of the asymmetric Laplacian matrices $L_{aug}^b$ and $L_{aug}^u$ defined in Eqs.~\eqref{eq:Lnewbi} and \eqref{eq:Lnewuni}, corresponding to directed graphs with bidirectional (panel A) and unidirectional (panel B) adversarial connections, respectively. In each case, a single node (either $i$, $j$, or $k$) is attacked with budget $-c<0$. The dashed black lines show the asymptotic relations $f(L_{aug}^b)=-1.1c$ (panel A) and $f(L_{aug}^u)=-0.1c$ (panel B). The colored dashed lines in panel A correspond to
$f(L_{aug}^b)=-2c+L_{xx}/2$, whereas those in panel B correspond to
$f(L_{aug}^u)=-(\sqrt{2}+1)c/2+f_{\mathrm{inter}}^x$, for $x=i,j,k$, where the intercept is
$f_{\mathrm{inter}}^x=(3+2\sqrt{2})L_{xx}$. The schematics are not drawn to scale and are intended only to illustrate the qualitative behavior.}
    %Schematic showing $f (L_{aug}^b)$ and $f (L_{aug}^u)$, the algebraic connectivity of asymmetric Laplacian matrices $L_{aug}^b$ and $L_{aug}^u$ in Eqs.\,\eqref{eq:Lnewbi} (panel A: directed graphs with bidirectional adversarial connections) and Eq.\,\eqref{eq:Lnewuni} (panel B: directed graphs with unidirectional adversarial connections), when only one node (either node $i$, $j$, or $k$) is pinned with the budget $-c < 0$. The dashed black line shows the line $f (L_{aug}^b) = -1.1c$ for panel A and shows the line $f (L_{aug}^u) = -0.1c$ for panel B. The colored dashed lines in panel A show $f (L_{aug}^b) = -2c + L_{xx}/2$ for $x = i, j,k$. The colored dashed lines in panel B show $f (L_{aug}^u) = -\frac{\sqrt{2}+1}{2} c + f_{inter}^x$ for $x = i, j,k$ where {the intersection with the ordinate axis} $f_{inter}^x = (3+2\sqrt{2})L_{xx}$. The schematics are not drawn in scale to enhance visualization.}
    \label{fig:directed}
\end{figure}









\subsection{General Directed Graphs} \label{sec:general}

The overall behavior of the algebraic connectivity $f$ for a general digraph (i.e., a digraph that is not balanced) subject to an adversarial node addition is more difficult to characterize. 
In the limit of a large budget $c$, we derive the same relationships we had previously obtained in the case of balanced digraphs: %the only difference that in this case the dependence is on the indegree of a node, rather than its degree. In fact, the following relations hold for a general digraph in the limit of a large budget $c$:
    $f(L_{aug}^b) = -2c + L_{ii}/2$ in the case of a bidirectional adversarial node addition and
     %the algebraic connectivity
    $f(L_{aug}^u) = -\frac{\sqrt{2}+1}{2}c + (3+2\sqrt{2})  L_{ii}$ in the case of a unidirectional adversarial node addition,
with the only difference that in this case the dependence is on the indegree $d_i^{in}=L_{ii}$ of a node, rather than its degree.
The above relations  indicate that for any directed graph, in the limit of a large budget $c$, the algebraic connectivity $f$ is defined by the indegree of the targeted node. 




In general, the curve $f(L_{aug})$ vs $c$ in the limit of small $c$, i.e., the initial slope of the curve for small $c$, depends non-trivially on the choice of the targeted node. 
However, certain relations can still be derived for the average slope in directed graphs, which is discussed next.








We direct our focus to the change in the algebraic connectivity in the limit of small $c$,
\begin{equation}
    \frac{d f}{d c} := \lim_{c \to 0^-} \dfrac{\tilde{f} - f_{new}}{c}.
\end{equation}
Here, $c \geq 0$ is the budget, and $f_{new}$ ($\tilde{f}$) is the algebraic connectivity corresponding to the Laplacian matrix $L_{aug}$ before the attack, i.e., to the matrix $\tilde{L}$:
\begin{align}
\begin{split}
    \tilde{L} = &
    % \begin{bmatrix}
    %     L & 0 \\
    %     0 & 0
    % \end{bmatrix}, 
    \begin{bmatrix}
    L & \pmb{0} \\ 
    \pmb{0}^\top &  0
    \end{bmatrix}.
\end{split}
\end{align}
We consider the two cases of (i) an attack via bidirectional connection ($L_{aug}^b$ in Eq.\,\eqref{eq:Lnewbi}), and (ii) an attack via unidirectional connection ($L_{aug}^u$ in Eq.\,\eqref{eq:Lnewuni}).
We target all nodes of the network one by one and calculate $d f / d c$ for each targeted node, i.e., we set $\bb_i = -c \pmb{e}_i$ where $\pmb{e}_i$ has all zero entries except for entry $i$ which is equal to 1.
Then, we calculate 
% the most negative slope $\min (d f / d c)$ and 
the average slope $< df / d c >$:
\begin{equation}
    % \min \left(\frac{d f}{d c} \right) := \min_i \frac{d f_i}{d c}, \quad 
    \left< \frac{d f}{d c} \right> := \dfrac{1}{N} \sum_{i=1}^N \frac{d f_i}{d c},
\end{equation}
where $f_i$ is the algebraic connectivity when $\bb_i$ is used in either Eqs.\,\eqref{eq:Lnewbi} or \eqref{eq:Lnewuni}.


Next, we investigate the algebraic connectivity for selected real networks,
 such as animal, biological, social, neural, trade, metabolic, and genetic networks.
The network data is obtained from \cite{konect}. {All the adjacency matrices have positive entries, with certain networks being weighted and certain being unweighted (detailed information on these networks can be found in Supplementary Note 4.)} 
For each network dataset, we apply our analysis to their largest strongly connected component. %, %s for each original network and set that as our network of interest. 
%which is a necessary condition for the emergence of synchronization (consensus) before the attack. 

{We emphasize that the slope $d f_i / d c$ is not easy to characterize for non-balanced digraphs since the slope $d f_i / d c$ and the indegree of the node $i$ may not have a monotonic relationship.
This is shown in Supplementary Note 4 using two examples of animal networks; for one, the slope is directly related to the indegree, while for the other, the slope is not correlated with the indegree.}


% \textcolor{red}{First, we emphasize that the slope $d f_i / d c$ is not easy to characterize for non-balanced digraphs. 
% This is shown in Figure~\ref{fig:animal}, where the top and bottom panels are plots of the slope $d f_i / d c$ versus the indegree of the node $i$ for two selected real networks. In particular, the bottom panel shows a non-monotonic relationship between two quantities.}



% \begin{figure}
%     \centering
%     \hspace{0.1\linewidth}\text{Animal network with $N=31$ nodes}
%     \includegraphics[width=0.7\linewidth]{figs/Animal_N31.pdf} \\
%     \vspace{0.01\linewidth}
%     \hspace{0.1\linewidth}\text{Animal network with $N=26$ nodes}
%     \includegraphics[width=0.7\linewidth]{figs/Animal_N26.pdf}
%     \caption{
%     The slope $d f_i / d c$ vs the node indegree $L_{ii}$ for two selected real animal networks with $N=31$ (top) and $N=26$ (bottom) nodes with bidirectional (blue diamonds) and unidirectional (red circle) adversarial connections. 
%     The top panel shows a plot in which the slope is directly related to the indegree, while the bottom panel shows a plot in which the slope is not correlated with the indegree.}
%     \label{fig:animal}
% \end{figure}



% Fig.\,\ref{fig:realnet} shows the results of the most negative slope $\min (d f / d c)$ and the mean slope $< df / d c >$ for a collection of real networks under bidirectional and unidirectional attacks. 
% The top (bottom) row panels show the results of bidirectional (unidirectional) attacks.
% The left (right) column panels show  $\min (d f / d c)$ ($< df / d c >$.) 
% It does not seem that the top row panels and the bottom left panel show a pattern.
% Surprisingly, the bottom right panel in Fig.\,\ref{fig:realnet} shows that the mean slope $< df / d c >$ scales exactly with the size of the network $N$ when the adversarial connection is unidirectional and the relation is $< df / d c > = -1/N$. 
% For the detailed derivation of this relation see Sec.\,\ref{sec:meanslope}


Figure\,\ref{fig:realnet} shows the mean slope $< df / d c >$, averaged over all the network nodes, for a collection of real networks under unidirectional attacks. 
We see the emergence of a clear scaling relation with the size of the network $N$, namely $< df / d c > = -1/N$. 
This result is %confirmed by our theoretical analysis see 
analytically proven in
Sec.\,\ref{sec:meanslope}. We conclude that larger networks are less reactive in average to intruder attacks.
We have also computed the most negative slope $\min (d f / d c):= \min_i d f_i / d c$ and the mean slope $< df / d c >$ under bidirectional and unidirectional attacks, and those plots are provided in Supplementary Note 4.


\begin{figure}
    \centering
    % \includegraphics[width=0.48\linewidth]{figs/min_bi_n.pdf} \hspace{0.01\linewidth}
    % \includegraphics[width=0.48\linewidth]{figs/ave_bi_n.pdf} \\
    % \vspace{0.03\linewidth}
    % \includegraphics[width=0.48\linewidth]{figs/min_uni_n.pdf} \hspace{0.01\linewidth}
    \includegraphics[width=0.6\linewidth]{figs/ave_uni_n.pdf}
    \caption{\textbf{Mean slope $\langle df/dc \rangle$ for unidirectional attacks on real-world networks.}
The mean slope $\langle df/dc \rangle$ for unidirectional (Eq.\,\eqref{eq:Lnewuni}) connections of the adversarial agent for selected real-world networks of size $N$. The dashed black curve shows the analytically predicted relation $\langle df/dc \rangle = -1/N$.
   %\textbf{Variation of the mean slope $< df / d c >$ for selected real networks with the network size $N$.}
    % The most negative slope $\min (d f / d c)$ and the mean slope $< df / d c >$ for both bidirectional (Eq.\,\eqref{eq:Lnewbi}) and unidirectional (Eq.\,\eqref{eq:Lnewuni}) connection of the adversarial agent for selected real networks with size $N$. The relation between $< df / d c >$ and $N$ in lower right panel is $< df / d c > = -1/N$.
    %We plot the mean slope $< df / d c >$ for unidirectional (Eq.\,\eqref{eq:Lnewuni}) connections of the adversarial agent for selected real networks with size $N$. The dashed black curve is the analytically predicted relation 
     %$< df / d c > = -1/N$.
    } 
    
    \label{fig:realnet}
\end{figure}


\subsection{Applications} \label{sec:applications}

Although our results are derived for the case of the consensus dynamics \eqref{consensus},  they have direct implications in the case of several real-world applications, described by the general set of equations \eqref{general}.
Below we discuss three applications of interest in the realm of man-made networks. %seek evidence 
{A detailed asymptotic stability analysis for each application is provided in the Methods. Specifically, the synchronization dynamics is analyzed in Methods Sec.\ \ref{Sec:MAI}, the power grid dynamics in Methods Sec.\ \ref{Sec:MAII}, and the formation control dynamics in Methods Sec.\ \ref{Sec:MAIII}.}

%\section{Applications}  \label{sec:applications}




As a first application, we first consider  synchronization of coupled chaotic oscillators.
A model for the emergence of synchronization in a network of coupled oscillators is the following,
\begin{equation} \label{eq:synch}
    \dot{\bx}_i (t) = \bF(\bx_i(t)) - \sigma \sum_{j = 1}^N L_{ij} \bH(\bx_j (t)),
\end{equation}
where $\bx_i(t) \in \mathbb{R}^m$ is the state of oscillator $i$ at time $t$, the functions $\bF: \mathbb{R}^m \to \mathbb{R}^m$ and $\bH: \mathbb{R}^m \to \mathbb{R}^m$ describe the local dynamics and the coupling functions, respectively. 
The matrix $L = [L_{ij}]$ is the Laplacian matrix that describes the network connectivity, and the scalar $\sigma \geq 0$ is the coupling strength. For this example, we set  the local dynamics $\bF$ to be the chaotic Lorenz oscillator. The functions $\bF$ and $\bH$ are
\begin{equation} \label{eq:Lorenz}
        \bx = \begin{bmatrix}
        x \\ y \\ z
    \end{bmatrix}, \ \
    % lorenz
    \bF(\bx) = \begin{bmatrix} 10(y-x) \\
    x(28 - z) - y \\
    xy - 2 z \end{bmatrix}, \ \  \bH(\bx) = \begin{bmatrix}
        0  \\
        y  \\
        0 
    \end{bmatrix}.
\end{equation}

We take the directed network shown in Fig.\,\ref{fig:application} A and set $\sigma=6$.
Figure \ref{fig:application}B shows the synchronous time evolution of the state components $x_i(t)$ for each oscillator before the attack. Panel C shows the loss of synchronization observed after the attack.
The loss of stability is the direct consequence of having mixed sign eigenvalues of the Laplacian matrix after the attack. 
The plots of the $y$ and $z$ components show similar behavior as the $x$ component and can be found in Supplementary Note 5. 
In Supplementary Note 5 we also provide an explanation for why the synchronous state becomes unstable when an attack occurs in terms of the {master stability function approach (MSF) \cite{fujisaka1983stability,pecora1998master}}. 


Reference \cite{sorrentino2007controllability} has used the master stability function to study stability of the synchronous solution in the case of the pinning control problem, for which a reference synchronous trajectory is selected to which all the network nodes need to synchronize. 
{There are, however, two main differences between the two approaches.} { {The first concerns the setting of the problem:} in pinning control one {considers} the role of a reference agent in driving the synchronization of a network of coupled systems {toward} the reference trajectory, whereas in our work {we consider} the role of an intruder agent in disrupting synchronization. In the former case the coupling of the reference agent to the rest of the network is cooperative, %(also sometimes called mutualistic or attractive), 
while in the latter it is adversarial. % (also sometimes called competitive or repulsive). 
{In particular, the insertion of an intruder effectively alters the coupling structure so that the relevant Laplacian has at least one negative part eigenvalue; as a result, the asymptotic stability analysis naturally involves the evaluation of the Master Stability Function at negative real arguments, corresponding to these eigenvalues, which is not the common approach in the literature. } The Methods Sec.\ I provides master stability function plots, over both negative and positive values of its argument, for the cases of the Lorenz, Rossler, and Chua oscillators, under various node-to-node couplings. For all the cases considered, we see that the MSF is positive when its argument is negative, implying instability of the synchronous solution.
{The second difference concerns the observable used to characterize the effect of the intervention: while Ref.\ \cite{sorrentino2007controllability} focuses on asymptotic stability through Lyapunov exponents, here we focus on the transient rate of instability. }} %Following the Reviewer's input we have now added an extensive discussion {of these points} in the revised manuscript. Both these points are important and should be kept in mind when attempting to address the question posed by the Reviewer.}






%{The first concerns the setting of the problem:} in pinning control one {considers} the role of a reference agent in driving the synchronization of a network of coupled systems {toward} the reference trajectory, whereas in our work {we consider} the role of an intruder agent in disrupting synchronization. In the former case the coupling of the reference agent to the rest of the network is cooperative, while in the latter it is adversarial. {In particular, the insertion of an intruder effectively alters the coupling structure so that the relevant Laplacian spectrum is shifted toward negative values; as a result, the stability analysis naturally involves the evaluation of the master stability function at negative arguments, corresponding to these eigenvalues.}  {The second difference concerns the observable used to characterize the effect of the intervention: while \cite{sorrentino2007controllability} focuses on asymptotic stability through Lyapunov exponents, here we focus on the transient rate of instability.}





%The results are sensitive to the particular choice of this reference solution, as well as to the particular choice of the individual nodal function $\bF$  and of the output function $\bH$. Different from \cite{sorrentino2007controllability}, here we derive conditions that are sufficient to destabilize the network dynamics and are independent from the choice of any particular reference trajectory and of the functions $\bF$ and $\bH$. \textcolor{red}{add more here}


\begin{figure*}
    \centering
    \includegraphics[width=0.8\linewidth]{figs/Fig_application.pdf}
    \caption{
    \textbf{Effects of intruder attacks in different applications.}
    The top row (A--C) illustrates the application to directly coupled chaotic oscillators, the middle row (D--F) to the power grid described by the nonlinear swing equation, and the bottom row (G--I) to linear single-integrator robot formation control. The left column (A, D, G) shows the corresponding network topologies. Panel A displays a directed network, whereas panels D and G show undirected networks. In all networks, black and orange links denote edge weights of $-1$ and $1$, respectively. The middle column (B, E, H) shows the dynamics in the absence of the adversarial agent (node A), while the right column (C, F, I) shows the dynamics in its presence.
In panel D, blue square nodes represent generators and orange circular nodes represent load nodes; the adversarial agent is also a load node. Panels B and C show the states $x_i(t)$ of the Lorenz oscillators; panels E and F show the angular positions $\vartheta_i(t)$ of the generators and loads in the swing equation; and panels H and I show the robot positions in the $xy$ plane, where diamonds denote the initial positions and stars denote the desired final positions.
    %  The top row panels (A-C) demonstrate the application of directly coupled chaotic oscillators, the middle row panels (D-F) depict the application of the power grid (nonlinear swing equation), and the bottom row panels (G-I) show the application of linear single-integrator robot formation control. The left column panels (A, D, G) show the topologies of their respective application. Panel A shows a directed network, and panels D and G show undirected networks. In all shown networks, black and orange links have a weight of $-1$ and $1$, respectively. The middle column panels (B, E, H) show the dynamics when the adversarial agent (denoted as node A) is not present, and the right column panels (C, F, I) show the dynamics when the adversarial agent is present. In panel D, the square blue nodes show the generators and the circle orange nodes are the load nodes (also the adversarial agent is a load node.) Panels B and C show the state $x_i (t)$ of the Lorenz oscillators; panels E and F show the angular position $\vartheta_i (t)$ of the loads/generators from the swing equation, and panels H and I show the position of the robots in $xy$ plane (diamonds label initial positions and stars label desired final positions.
    }
    \label{fig:application}
\end{figure*}


Next, we study the effect of a unidirectional attack on a scale-free network of Lorenz oscillators coupled in the $y$-variable, Eq.\ \eqref{eq:Lorenz}. {For this case, the MSF analysis for Eq.~\eqref{eq:Lorenz}, presented in Sec.\ I of the Methods, shows that, for negative values of its argument, the master stability function becomes increasingly negative as the magnitude of the argument increases.}
We use the algorithm from \cite{goh2001universal} to randomly generate an undirected and unweighted scale-free graph, which is shown in Fig. \ref{fig:synch_node} A. 
The network has $N=100$ nodes and the degree distribution follows a power law $P_D(k) \propto k^{-\gamma}$ with $\gamma = 2.33$ with an average degree of $10$.
We aim to study the effect of the attack on different nodes; we set $\sigma = 10$, $c = 1$, and vary the node $i$ that is targeted.
\begin{comment}
The matrix $L(t) = [L_{ij}(t)]$ is the Laplacian matrix that describes the network connectivity and is given by
\begin{equation} \label{eq:Lsynch}
    L(t) = \begin{cases}
        \begin{bmatrix}
            L_0 & \pmb{0} \\
            \pmb{0}^\top & 0
        \end{bmatrix}, \quad & t < T, \\
        \\
        \begin{bmatrix}
            L_0 & \bb^i \\
            \pmb{0}^\top & 0
        \end{bmatrix}, \quad & t \geq T.
    \end{cases}
\end{equation}
Here, $\bb^i$ is the vector that characterizes the attack from the intruder to node $i$.
We set the time of the attack to be $T=4s$.

We study the effect of attacking nodes with different indegrees.
The choice of targeted nodes are defined by, 
\begin{equation} \label{eq:bsynch}
    % L_0 = \begin{bmatrix}
    % -1 & 1 & 0 \\ 
    % 1 & -2 & 1 \\ 
    % 1 & 1 & -2 
    % \end{bmatrix}, \quad \bb^1 = \begin{bmatrix}
    % -2 \\ 
    % 0 \\ 
    % 0 
    % \end{bmatrix}, \quad \bb^2 = \begin{bmatrix}
    % 0 \\ 
    % -2 \\ 
    % 0 
    % \end{bmatrix}.
    \bb^i = \begin{bmatrix}
        0 & \hdots & 0 & -c & 0 & \hdots & 0 
    \end{bmatrix}^\top,
\end{equation}
where the $i$th component of $\bb^i$ is equal to $-c$. 
\end{comment}
We select the nodes $i = 1,5,17,94,100$ 
with degrees $82, 47, 23, 13$, and $8$, respectively, and associate to these nodes the colors blue, orange, yellow, purple, and green, respectively. 
%based on their degrees, with node 1 having the highest degree and node 100 having the lowest degree.
For more details on this example, see Supplementary Note 6.

The perturbation vector for each node is calculated as $\delta \bx_i(t) = \bx_i (t) - \frac{1}{N} \sum_{j = 1}^N \bx_j \in \mathbb{R}^3$. 

We then form the transverse perturbations vector $\delta \bX(t) = [\delta \bx_1(t)^\top, \ \delta \bx_2(t)^\top, \hdots, \delta \bx_{N}(t)^\top]^\top \in \mathbb{R}^{3N}$ and calculate its norm:
\begin{equation} \label{eq:trans_synch}
    \| \delta \bX(t) \| := \sqrt{\delta \bX(t) ^\top \delta \bX(t)}.
\end{equation}
\color{black}
If the oscillators synchronize, then $\| \delta \bX(t) \| \to 0$, otherwise $\| \delta \bX(t) \| \gg 0$.

We see from Fig.~\ref{fig:synch_node} B that before the attack, i.e., on the left side of the dashed black line, the norm of the transverse perturbations approaches zero, indicating that the Lorenz oscillators have synchronized. However, after attacking either one of the nodes $1,5,17,94,100$,
we see an immediate loss of synchronization, with a faster rate of transient instability when the nodes with lower degrees are targeted. For example, the attack on node 100 with degree 8 (the green curve) results in much larger $\| \delta \bX(t) \|$ than the attack on node 1 with degree 82 (the blue curve).
To conclude, this example shows that in a nonlinear networked system, lower indegree nodes are more vulnerable to intruder attacks. {Note that this is consistent with the MSF analysis for the case of Eq.~\eqref{eq:Lorenz}, presented in Sec.\ I of the Methods, since attacking nodes with lower indegrees corresponds to evaluating the master stability function at more negative values of its argument, for which the associated maximum Lyapunov exponent is larger.}




%In the Supplementary Note 6 we provide MSF plots associated with Lorenz, Rossler, and Chua's circuit oscillators under various couplings. We show that the MSF is positive for the negative range of its argument; thus, the synchronization is locally unstable.


% \begin{figure}
%     \centering
%     \includegraphics[width=0.6\linewidth]{figs/sync_adver_node.pdf}
%     \caption{
%     We consider the dynamics of a network of coupled Lorenz oscillators, Eq.\,\eqref{eq:synch}, and plot
%     the norm of the transverse perturbations $ \| \delta \bX(t) \|$ as a function of time $t$. The blue curve corresponds to the case that  node $i=1$ is attacked, i.e., $\bb^1$ is chosen in Eq.\,\eqref{eq:bsynch}.  The red curve corresponds to the case that  node $i=2$ is attacked, i.e., $\bb^2$ is chosen in Eq.\,\eqref{eq:bsynch}. 
%      The dashed black line denotes the start time of the attack at $t = T = 4s$. The initial conditions of the oscillators are chosen to be the same in the two cases that nodes $1$ and $2$ are attacked.}
%     \label{fig:synch_node}
% \end{figure}



\begin{figure}
    \centering
    % \text{A} \\
    % \vspace{0.01\linewidth}
    % \includegraphics[width=0.7\linewidth]{figs/synch_adver_BA_graph.pdf} \\
    % \vspace{0.03\linewidth}
    % \text{B} \\
    % \includegraphics[width=0.9\linewidth]{figs/sync_adver_node_BA.pdf}
    \includegraphics[width = \linewidth]{figs/Fig5_new.pdf}
    \caption{\textbf{Effect of the attacked node on network destabilization.}
Panel A shows a randomly generated scale-free network under attack through either of the colored nodes (blue $i=1$, orange $i=5$, yellow $i=17$, purple $i=94$, and green $i=100$). Panel B considers the dynamics of a network of coupled Lorenz oscillators, Eq.~\eqref{eq:synch}, and plots the norm of the transverse perturbations $\|\delta \mathbf{X}(t)\|$ as a function of time $t$ for attacks targeting different nodes $i$. Each curve corresponds to a different attacked node, all with the same attack budget $c=1$. The dashed black line indicates the onset of the attack at $t=4\,\mathrm{s}$. The initial conditions of the oscillators are identical in all simulations, so that the observed differences are solely due to the attacked node. The degree of each attacked node, $d_i$, is also indicated in the plot.}
    %Panel A shows a randomly generated scale-free network under attack through either of the colored nodes (blue $i=1$, orange $i=5$, yellow $i=17$, purple $i=94$, and green $i=100$).  In panel B, we consider the dynamics of a network of coupled Lorenz oscillators, Eq.\,\eqref{eq:synch}, and plot the norm of the transverse perturbations $ \| \delta \bX(t) \|$ as a function of time $t$ and the attacked node $i$. Different curves correspond to different attacked nodes, all with the same budget $c = 1$.  The dashed black line denotes the start time of the attack at $t = 4s$. The initial conditions of the oscillators are chosen to be the same in all cases where different nodes are attacked. The plot also provides the degree of node $i$ as $d_i$.}
    \label{fig:synch_node}
\end{figure}








% \begin{figure*}
%     \centering
%     \text{A) $L$} \\
%     \includegraphics[width=0.32\linewidth]{figs/Lx.pdf}
%     \includegraphics[width=0.32\linewidth]{figs/Ly.pdf}
%     \includegraphics[width=0.32\linewidth]{figs/Lz.pdf} \\
%     \text{B) $L_{aug}$} \\
%     \includegraphics[width=0.32\linewidth]{figs/Lnewx.pdf}
%     \includegraphics[width=0.32\linewidth]{figs/Lnewy.pdf}
%     \includegraphics[width=0.32\linewidth]{figs/Lnewz.pdf}
%     \caption{The synchronization dynamics in Eq.\,\eqref{eq:synch} for the two Laplacian matrices $L$ and $L_{aug}$ from Eq.\,\eqref{eq:Lconsensus}. Panels A (top row) and B (bottom row) show the state trajectories for each system with the Laplacian matrices $L$ and $L_{aug}$, respectively. The dashed black line shows the trajectory associated with the newly added adversarial agent.}
%     \label{fig:synch}
% \end{figure*}

% Our results generalize to the case of networks characterized by nonlinear dynamics. Supplementary Note 3 shows how the severity of an intruder attack depends on the degree of the attacked node in both the cases of synchronization of coupled chaotic oscillators and phase-locking in Kuramoto oscillators. In both cases, we see that attacks targeting low-degree nodes result in a higher rate of transient instability in these nonlinear networks. 
% \textcolor{red}{Supplementary Note 3 also considers the case of an undirected scale free network, for which we clearly see that the rate of instability generated by an intruder attack anticorrelates with the degree of the node targeted by the intruder.} 
%In particular, attacking nodes in a scale-free network of chaotic oscillators are shown to be more destructive when lower-degree nodes (outliers) are targeted, unlike the common belief that the hubs (higher-degree nodes) are generally considered to be more critical.


%Supplementary Note 5 also provides an example of phase-locking in Kuramoto oscillators. Also in this case, we see that attacks targeting low-degree nodes result in a higher rate of transient instability.


As a second application we consider power grid dynamics. The swing equation is  a classical model for the dynamics of a power grid, \cite{nishikawa2015comparative,bhatta2021modal}
\begin{equation} \label{eq:swingsimple}
     \ddot{\vartheta}_i (t) = p_i - {\gamma} \dot{\vartheta}_i (t) - \sum_{j= 1, j \neq i}^N A_{ij} \sin(\vartheta_i (t) - \vartheta_j (t))
\end{equation}
where $\vartheta_i (t)$ is the angular displacement of node/rotor $i$, the scalar $p_i$ denotes the power consumption ($p_i <0$ for loads) or power generation ($p_i >0$ for generators,) the scalar ${\gamma >0}$ denotes the damping ratio, and the symmetric adjacency matrix $A = [A_{ij}]$ describes the connectivity of the grid, i.e.,  $A_{ij}=A_{ji}=1$ ($0$) if there is (is not) a transmission line between between nodes $i$ and $j$. 
{Note that although Eq.\ \eqref{eq:swingsimple} is not a particular case of Eq.\ \eqref{general}, the linearized swing equation (see Supplementary Note 7) on which our analysis is based is.}

%{Each node (either a generator or a load) can be seen as an oscillator and the network is typically modeled as an undirected graph, where edges represent the presence of transmission lines between the nodes.In what follows, we will perform a modal decomposition of the linear swing equation, add the adversarial agent, and show that at least one mode of the system becomes unstable, which will lead to total instability.}

% The swing equation is 
% \begin{equation}
%     \dfrac{2 H_i}{\omega_R} \ddot{\vartheta}_i (t)+ \dfrac{D_i}{\omega_R} \dot{\vartheta}_i (t) = b_i - \sum_{\substack{j= 1, \\ j \neq i}}^N \tilde{A}_{ij} \sin(\vartheta_i (t) - \vartheta_j (t)) 
% \end{equation}
% where $\vartheta_i (t)$ is the displacement of node $i$, $\omega_R$ is the reference frequency, $H_i$ is the inertia, $D_i$ is the damping constant, and $b_i$ is the power generated ($b_i >0$ for generators) or the power consumed ($b_i < 0$ for loads.)
% The connectivity if given via the symmetric adjacency matrix $\tilde{A} = [\tilde{A}_{ij}]$ where if nodes $i$ and $j$ are connected, then $\tilde{A}_{ij} = \tilde{A}_{ji} \neq 0$, otherwise $\tilde{A}_{ij} = 0$.
% By defining $\gamma_i := \frac{D_i}{2 H_i}$ and $J_i := \frac{\omega_R}{2 H_i}$, the above equation is rewritten as
% \begin{equation} \label{eq:swingsimple}
%      \ddot{\vartheta}_i (t) = - \gamma_i \dot{\vartheta}_i (t) + p_i - \sum_{j= 1, j \neq i}^N A_{ij} \sin(\vartheta_i (t) - \vartheta_j (t))
% \end{equation}
% where $p_i = J_i b_i$ and $A_{ij} = J_i \tilde{A}_{ij}$.




{We emphasize that the swing equation \eqref{eq:swingsimple} conveniently demonstrates how an intruder oscillator can be mathematically incorporated into the Laplace matrix through a negative sign. By constraining the intruder oscillator to remain in anti-phase with the rest of the network—i.e., \(\sin(\vartheta_i(t) - \vartheta_j(t) \pm \pi) = -\sin(\vartheta_i(t) - \vartheta_j(t))\)—a negative sign naturally emerges in the coupling kernel. This negative sign is then absorbed into the adjacency matrix as a negative link connection, represented by \(-A_{ij}\).} 

To illustrate the effect of the addition of an adversarial agent, we provide a numerical example. 
Here, we set %$\gamma_i = \gamma = 0.9$
$\gamma = 0.9$ and the network shown in Fig.\,\ref{fig:application} D shows the grid topology before and after the addition of the adversarial agent. 
% the matrices
% \begin{align}
% \footnotesize
% \begin{split}
%     A & = \begin{bmatrix}
%     0 & 1 & 1 & 1 & 0 & 0 & 0 \\ 
%     1 & 0 & 0 & 0 & 1 & 0 & 0 \\ 
%     1 & 0 & 0 & 0 & 0 & 1 & 0 \\ 
%     1 & 0 & 0 & 0 & 0 & 0 & 1 \\ 
%     0 & 1 & 0 & 0 & 0 & 0 & 0 \\ 
%     0 & 0 & 1 & 0 & 0 & 0 & 0 \\ 
%     0 & 0 & 0 & 1 & 0 & 0 & 0 
%     \end{bmatrix}, \
%     A_{new}  = \begin{bmatrix}
%     0 & 0 & 0 & 0 & 0 & 0 & 0 & -1 \\ 
%     0 & 0 & 1 & 1 & 1 & 0 & 0 & 0 \\ 
%     0 & 1 & 0 & 0 & 0 & 1 & 0 & 0 \\ 
%     0 & 1 & 0 & 0 & 0 & 0 & 1 & 0 \\ 
%     0 & 1 & 0 & 0 & 0 & 0 & 0 & 1 \\ 
%     0 & 0 & 1 & 0 & 0 & 0 & 0 & 0 \\ 
%     0 & 0 & 0 & 1 & 0 & 0 & 0 & 0 \\ 
%     -1 & 0 & 0 & 0 & 1 & 0 & 0 & 1 
%     \end{bmatrix},
% \end{split}
% \end{align}
% Note that the adversarial agent is set as the last node (node $8$) of the graph with the adjacency matrix $A_{new}$ and the budget (the weight of the connection) is assumed to be $-1$.
Nodes $1,2,3$ and $4$ are loads, nodes $5,6$ and $7$ are generators, and the adversarial node is denoted as node A with a black connection with weight $-1$.
The vector $\bp = [
-0.3, \ -0.3, \ -0.3, \ -0.3, \ 0.4, \ 0.4, \ 0.4 
]^\top$ describes the power generated/consumed in the original power grid.
Without loss of generality, for the adversarial agent, we set $p_{A} = -0.1$, so the adversarial agent acts as a load.


% Figure \ref{fig:application} E demonstrates the addition of an adversarial agent to a power grid. 
As seen in Fig.\,\ref{fig:application} E (Fig.\,\ref{fig:application} F), the dynamics of the power grid converges (does not converge) to a fixed point before (after) the attack.
%After the attack, the dynamics no longer converges as shown in Fig.\,\ref{fig:application} F.
In Supplementary Note 7, we show analytically how the addition of the adversarial agent results in at least one unstable mode of the linearized swing equation, which in turn destabilizes the nonlinear dynamics.

\begin{comment}
In what follows, we prove that by the addition of an adversarial agent, the dynamics of the power grid as described by the swing equation no longer converges to fixed points.


\begin{figure}
    \centering
    \includegraphics[width=0.7\linewidth]{figs/swing_original.pdf} \\
    \vspace{0.03\linewidth}
    \includegraphics[width=0.3\linewidth]{figs/swing_network.pdf} \\
    \vspace{0.03\linewidth}
    \includegraphics[width=0.7\linewidth]{figs/swing_adversarial.pdf}
    \caption{{
    Example of a power grid dynamics based on the swing equation \eqref{eq:swingsimple}. Plots show the oscillator phase $\vartheta_i (t)$ vs time $t$ for different nodes in the network in the middle panel. Blue square nodes are generators and the orange circle nodes are loads. The adversarial agent is a load and is shown as node $A$ with hashed black/orange lines. The orange links have weights equal to 1 and the black link has a weight of $-1$.  The top (the bottom) panel shows the dynamics before (after) the addition of the adversarial agent.  }}
    \label{fig:swing}
\end{figure}

Linearizing Eq.\,\eqref{eq:swingsimple} about the fixed-point $\vartheta^*_i$ results in
\begin{equation} \label{eq:swinglinear}
    \ddot{\theta} (t) = - \gamma_i \dot{\theta} (t) + p_i - \sum_{j=1}^N L_{ij} \theta_j (t)
\end{equation}
where $\theta_i (t) := \vartheta_i (t) - \vartheta^*_i$ is the deviations about the fixed-point $\vartheta^*_i$ and the matrix $L = [L_{ij}]$ is a symmetric Laplacian matrix where $L_{ij} = \delta_{ij} \sum_{k=1}^N A_{ik} \cos(\vartheta_k^* - \vartheta_i^*) -  A_{ij} \cos(\vartheta_j^* - \vartheta_i^*)$. 
Here, $\delta_{ij}$ is the Kronecker delta. 
In what follows, we assume $\gamma = \gamma_i, \forall i$, for simplicity.
By defining $\btheta := [\theta_1 , \ \theta_2, \hdots, \ \theta_N]^\top$ and $\bp := [p_1, \ p_2, \hdots, \ p_N]^\top$, the compact form of Eq.\,\eqref{eq:swinglinear} is obtained:
\begin{equation} \label{eq:swingcompact}
    \ddot\btheta(t) = - \gamma \dot{\btheta}(t) - L \btheta(t) + \bp.
\end{equation}
Since the Laplacian matrix $L$ is symmetric, it can be diagonalized via its eigenvectors. 
The diagonal matrix $\Lambda = V L V^\top$ has the real eigenvalues $\lambda_1, \lambda_2, \hdots, \lambda_N$ of the Laplacian matrix $L$ on its diagonal.
The matrix $V$ has the eigenvectors of the Laplacian matrix $L$ as its columns.
Assuming the network is connected, there is only one zero eigenvalue which we denote via $\lambda_1 = 0$, and all other eigenvalues $\lambda_i \neq 0, \ i > 1$.
By setting $\boldeta(t) := V^\top \btheta(t)$ and $\bq := V^\top \bp$, the coupled compact equation \eqref{eq:swingcompact} is then decomposed as
\begin{equation}
    \ddot \boldeta (t) = - \gamma \dot{\boldeta} (t) - \Lambda \boldeta(t) + \bq.
\end{equation}
The following equation describes the dynamics of each mode of the oscillation 
\begin{equation}
    \ddot \eta_i (t) = -\gamma \dot \eta_i(t) - \lambda_i \eta_i(t) + q_i, \quad i = 1, \hdots, N.
\end{equation}
It is now clear that if any of $\lambda_i < 0$ then the mode $\eta_i (t)$ of the system becomes unstable.
Since the overall behavior of the system is the linear combination of the modes, having an unstable mode results in overall instability of the system, i.e., if $\exists i: \eta_i(t) \to \infty$, then $\theta_j (t) \to \infty, \forall j$.
Based on Proposition \ref{prop1}, we know if the adversarial agent attacks the system, then we have at least one negative eigenvalue that resulting in destroying the stability of the linearized swing equation \eqref{eq:swinglinear} which in turn suggests that the nonlinear dynamics of the power grid in Eq.\,\eqref{eq:swingsimple} have unstable fixed points.
\end{comment}


As a third application we consider  distributed formation control. %We present the last example of the application of an adversarial agent in distributed formation control.
%Formation control of multi-agent systems is mainly the cooperation between several robots to perform a certain task.
The formation control of robots is used in both civilian and military applications, see \cite{oh2015survey,wang2021formation,yu2022decentralized}.
Here, we focus on positioning a set of robots  on a given desired pattern, e.g., the robots are positioned around a circle with a known diameter.

We take the shape-based formation control strategy from \cite{Morbidi2022Functions}:
\begin{equation} \label{Eq_formation_control}
    \dot{\bX} (t) = (- F(L) \otimes I_2 ) (\bX(t) - \bbeta),
\end{equation}
where the state vector $\bX(t) = [x_1(t), y_1 (t), x_2 (t), y_2 (t), \hdots, x_N (t), y_N (t)]^\top \in \mathbb{R}^{2N}$ is the concatenation of the $xy$ coordinates of $N$ robots at time $t$.
The matrix $L$ is the symmetric Laplacian matrix that describes the connectivity between the robots, the function $F(L): \mathbb{R}^{N \times N} \to \mathbb{R}^{N \times N}$ is a function of the Laplacian matrix $L$.
The network topology for this example with $N=5$ is shown in Fig.\,\ref{fig:application} G.
We set $F(L) = I_N - \exp (-3L)$ as in  \cite{Morbidi2022Functions}.
The vector $\bbeta \in \mathbb{R}^{2n}$ is a constant vector of target positions, which we
 position around a circle with radius $1$ in equidistant phases, i.e., $\bbeta = [\sin (0), \cos (0), \sin (\frac{2 \pi}{5}), \cos (\frac{2 \pi}{5}),\hdots, \sin (\frac{8 \pi}{5}), \cos (\frac{8 \pi}{5})]^\top$.
For the adversarial agent, we set its corresponding $\bbeta$ to the origin $[0,0]^\top$.




Figure \ref{fig:application} H shows the formation control in the absence of the adversarial agent, where each robot converges to its respective target position in less than $20$ seconds. Note that this can be seen as a case of cluster synchronization, since different agents converge to different states.
However, Fig.\,\ref{fig:application} I shows that in the presence of the adversarial agent, the robots no longer reach the target positions and the closed-loop system becomes unstable, position and velocity of the robots tend to infinity.)





\subsection{Phase synchronization based on Kuramoto dynamics}

Here we consider the case of phase synchronization in a network of coupled Kuramoto oscillators \cite{KuraBOOK},
\begin{equation} \label{eq:Kuramoto}
    \dot{\theta}_i (t) = \omega_i +\dfrac{1}{N+1} \sum_{j=1}^{N+1} A_{ij} (t) \sin(\theta_j (t) - \theta_i (t)),
\end{equation}
$i = 1,\hdots, N+1$, where the scalar $\theta_i (t)$ is the phase angle of oscillator $i$ at time $t$, the scalar $\omega_i$ is the frequency of  oscillator $i$. 
We require different frequencies for different oscillators, i.e., $\omega_i \neq \omega_j$, $\forall i \neq j$.
The time-varying connectivity is described by the adjacency matrix $A(t) = [A_{ij}(t)]$ where if node $i$ receives a connection from node $j$ at time $t$, then $A_{ij}(t) \neq 0$, otherwise $A_{ij}(t) = 0$.

We assume that a unidirectional adversarial connection from the intruder to one of the network nodes is  
established at some time $T > 0$. 
More precisely, the time-varying adjacency matrix is given by,
\begin{equation} \label{eq:Akuramoto}
    A(t) = \begin{cases}
        \begin{bmatrix}
            A_0 & \pmb{0} \\
            \pmb{0}^\top & 0
        \end{bmatrix}, \quad & t < T, \\
        \\
        \begin{bmatrix}
            A_0 & \bb^i \\
            \pmb{0}^\top & 0
        \end{bmatrix}, \quad & t \geq T.
    \end{cases}
\end{equation}
Here, $\bb^i$ is the vector that characterizes the attack from the intruder to node $i$ and $A_0$ is the binary $N$-dimensional  adjacency matrix describing the network connectivity, apart from the intruder.
In what follows, we use both the norm of the vector of phase perturbations and the order parameter (both defined next) to
characterize the effect of the attack on the nonlinear dynamics of coupled Kuramoto oscillators \eqref{eq:Kuramoto}.

The norm of the vector of phase perturbations is equal to,
\begin{equation}
    \| \delta \theta(t) \| := \sqrt{\sum_{j=1 }^{N}\left(\theta_j (t) - \frac{1}{N} \sum_{k=1}^{N} \theta_k (t)  \right)^2}. 
\end{equation}
It approaches a constant value in time if oscillators $i=1,\hdots, N$ are phased-locked and varies in time otherwise.
Also, we calculate the order parameter defined as,
\begin{equation} \label{eq:order}
{
    \rho(t) = \left| \dfrac{1}{N} \sum_{j=1}^{N} e^{i \theta_j(t)} \right|
}
\end{equation}
where $i = \sqrt{-1}$. The order parameter provides a normalized index of synchronization among all the oscillators: it is $\rho \approx 1$, when the oscillators are phase-locked, and $\rho \approx 0$ when there is no synchrony.
\begin{comment}
In our example, we set  $T=15s$, and randomly select the frequencies $\omega_1 =     0.5010, \omega_2 = 0.5030, \omega_3 = 0.4962, \omega_4 = 0.4860$. 
The topology and the attacks are defined by 
\begin{equation} \label{eq:bKuramoto}
    A_0 = \begin{bmatrix}
    0 & 1 & 0 \\ 
    1 & 0 & 1 \\ 
    1 & 1 & 0 
    \end{bmatrix}, \quad \bb^1 = \begin{bmatrix}
    -1 \\ 
    0 \\ 
    0 
    \end{bmatrix}, \quad \bb^2 = \begin{bmatrix}
    0 \\ 
    -1 \\ 
    0 
    \end{bmatrix}.
\end{equation}
and study the effect of the attack on nodes $1$ and $2$ when once $\bb^1$ is used in Eq.\,\eqref{eq:Akuramoto} and once $\bb^2$ is used.
We calculate the phase perturbations
\begin{equation}
    \| \delta \theta(t) \| := \sqrt{\sum_{j=1 }^{N}\left(\theta_j (t) - \frac{1}{N} \sum_{k=1}^{N} \theta_k (t)  \right)^2} 
\end{equation}
that will remain constant in time if oscillators $i=1,\hdots, N$ are phased-locked and will vary in time otherwise.

\ref{fig:kuramoto} shows that before the attack, on the left side of the dashed black line, the phase perturbations approaches a constant value which indicates that the Kuramoto oscillators are phased locked.
The figure shows a faster transient instability rate immediately after attacking node $1$ than attacking node $2$.
This is due to node $1$ having a lower indegree than node $2$.
Thus, this example shows lower indegree nodes in nonlinear networked systems are more vulnerable to intruder attacks.

\begin{figure}
    \centering
    \includegraphics[width=0.6\linewidth]{figs/Kuramoto.pdf}
    \caption{
    The phase perturbation $ \| \delta \theta(t) \|$ in time $t$ for a network of Kuramoto oscillators, Eq.\,\eqref{eq:Kuramoto}, as node $i=1$ and node $i=2$ are attacked. The attacks on nodes $1$ and $2$ are characterized by $\bb^1$ and $\bb^2$ in Eq.\,\eqref{eq:bKuramoto}. The dashed black line denotes the start time of the attack at $t = T = 15s$. The initial condition of the oscillators is kept the same when attacking nodes $1$ and $2$.}
    \label{fig:kuramoto}
\end{figure}
\end{comment}

In our first example, we set  $T=15s$, and randomly select the frequencies $\omega_i$ from a normal distribution with a mean $0.5$ and a standard deviation $0.01$.
The initial conditions $\theta_j(0), \ j=1,\hdots, N$, are drawn from a standard normal distribution and then normalized such that $\sum_{j=1}^{N} \theta_j(0)^2 = 1$.
We set the initial condition of the intruder to be $\theta_{N+1} (0) = 0$.
The network topology and the attacks are given by, 
\begin{equation} \label{eq:bKuramoto}
    A_0 = \begin{bmatrix}
    0 & 1 & 0 \\ 
    1 & 0 & 1 \\ 
    1 & 1 & 0 
    \end{bmatrix}, \quad \bb^1 = \begin{bmatrix}
    -1 \\ 
    0 \\ 
    0 
    \end{bmatrix}, \quad \bb^2 = \begin{bmatrix}
    0 \\ 
    -1 \\ 
    0 
    \end{bmatrix}.
\end{equation}
We compare the two cases of an attack on node $1$ with indegree one and node $2$ with indegree two.
The former (latter) case corresponds to using $\bb^1$ ($\bb^2$) in Eq.\,\eqref{eq:Akuramoto}.
% \ref{fig:Kuramoto_order} shows the order parameter in time as either node 1 or 2 is attacked.
% It is evident that before the attack, the oscillators became phase-locked.
% We see a greater divergence rate in the order parameter when node 1 is attacked in comparison to when node 2 is attacked.
% Thus, it is shown once again that nodes with lower indegrees (here node 1) are more vulnerable to intruder attacks.
% \begin{figure}
%     \centering
%     \includegraphics[width=0.6\linewidth]{figs/Kuramoto_order.pdf}
%     \caption{
%     The order parameter $ \rho$ in time $t$ for a network of Kuramoto oscillators, Eq.\,\eqref{eq:Kuramoto}, as node $i=1$ and node $i=2$ are attacked. The attacks on nodes $1$ and $2$ are characterized by $\bb^1$ and $\bb^2$ in Eq.\,\eqref{eq:bKuramoto}. The dashed black line denotes the start time of the attack at $t = T = 15s$. The initial condition of the oscillators is kept the same when attacking nodes $1$ and $2$.}
%     \label{fig:Kuramoto_order}
% \end{figure}
% Next, we run the above example of Kuramoto oscillators for multiple choices of random initial conditions and random frequencies.
% We draw the frequencies from a normal distribution with a mean $0.5$ and a standard deviation $0.01$.
% The initial conditions $x_j(0)$ are drawn from a standard normal distribution and then normalized such that $\sum_j x_j(0)^2 = 1$.
% We keep the rest of the parameters the same as before.
We repeat this numerical experiment 500 times and evaluate both $ \| \delta \theta(t) \|$ and $\rho$.

Figure \ref{fig:Kuramotorandom} shows that before the attack the oscillators become phase-locked and that after the attack phase locking is lost. 
We see a greater divergence rate both in the norm of the phase perturbation vector and in the order parameter, see panels \ref{fig:Kuramotorandom} A and B, respectively, when node 1 is attacked, compared to the case when node 2 is attacked.
%This indicates that nodes with lower indegrees (here node 1) are more vulnerable to intruder attacks.
This indicates that attacks targeting low-degree nodes result in a higher rate of transient instability.

\begin{figure}
    \centering
    % \text{A} \\
    % \includegraphics[width=0.6\linewidth]{figs/Kuramoto_example_band_delta.pdf} \\
    % \text{B} \\
    % \hspace{0.02\linewidth}
    % \includegraphics[width=0.58\linewidth]{figs/Kuramoto_example_band_rho.pdf}
    \includegraphics[width=0.7\linewidth]{figs/Fig6_new.pdf}
    \caption{\textbf{Effect of the attacked node on the dynamics of a network of Kuramoto oscillators.}
We consider a network of Kuramoto oscillators and plot the norm of the phase perturbation vector, $\|\delta\theta(t)\|$ (panel A), and the order parameter, $\rho(t)$ (panel B), as functions of time $t$. The blue curves correspond to an attack on node $i=1$, i.e., $\mathbf{b}^{1}$ is selected in Eq.~\eqref{eq:bKuramoto}, whereas the red curves correspond to an attack on node $i=2$, i.e., $\mathbf{b}^{2}$ is selected. The dashed black line marks the onset of the attack at $t=T=15\,\mathrm{s}$. Prior to the attack, the blue and red curves are identical since the same initial conditions are used in both simulations. The results show the averages of $\|\delta\theta(t)\|$ (panel A) and $\rho(t)$ (panel B) over 500 realizations, each corresponding to a different random choice of the initial conditions and the natural frequencies of the oscillators, as described in the text. The shaded regions indicate one standard deviation across the realizations.}
    %We consider the dynamics of a network of Kuramoto oscillators and plot the norm of the phase perturbation vector $ \| \delta \theta(t) \|$ (panel A) and the order parameter $ \rho(t)$ (panel B) as a function of time $t$. The blue curve corresponds to the case that node $i=1$ is attacked, i.e., $\bb^1$ is chosen in Eq.\,\eqref{eq:bKuramoto}.  The red curve corresponds to the case that node $i=2$ is attacked, i.e., $\bb^2$ is chosen in Eq.\,\eqref{eq:bKuramoto}.  The dashed black line denotes the start time of the attack at $t = T = 15s$. Before the attack, the blue and red curves are identical. The initial conditions of the oscillators are kept the same in the two cases that nodes $1$ and $2$ are attacked. The plots are averages of $ \| \delta \theta(t) \|$ (panel A) and of $ \rho(t)$ (panel B) over 500 realizations, where each realization corresponds to a different random choice of the initial conditions and of the frequencies of the oscillators, as explained in the text. The shaded background shows the standard deviation over different choices of the initial conditions and of the frequencies.}
    \label{fig:Kuramotorandom}
\end{figure}







We repeat the experiment above for a randomly generated scale-free network of coupled Kuramoto oscillators, which we choose to be the network of Fig. \ref{fig:synch_node} (reproduced for clarity in \ref{fig:Kuramoto_SF} A.)
Panels B and C show that before the attack the oscillators become phase-locked and that after the attack phase locking is lost. When  nodes of lower indegree are attacked,
we see a greater divergence rate both in the norm of the phase perturbation vector and in the order parameter, see panels \ref{fig:Kuramoto_SF} A and B, respectively. 
This confirms once again that nodes with lower indegrees (e.g., node 100) are more vulnerable to intruder attacks.

{A detailed asymptotic stability analysis for the case of Kuramoto dynamics is provided in the Methods Sec.\ \ref{Sec:MAIIII}.}


\begin{figure*}
    \centering
    % \text{A} \hspace{0.35\linewidth} \text{B} \hspace{0.35\linewidth} \text{C} \\
    % \vspace{0.01\linewidth}
    % \includegraphics[width=0.24\linewidth]{figs/synch_adver_BA_graph.pdf} 
    % \includegraphics[width=0.35\linewidth]{figs/dtheta_SF.pdf} 
    % \includegraphics[width=0.35\linewidth]{figs/rho_SF.pdf}
    \includegraphics[width = \linewidth]{figs/Fig7_new.pdf}
    \caption{\textbf{Effect of the attacked node on destabilization of Kuramoto dynamics.}
    Panel A shows a randomly generated scale-free network under attack through either one of the colored nodes with different indegrees (blue $i=1$ with degree $82$, orange $i=5$ with degree $47$, yellow $i=17$ with degree $23$, purple $i=94$ with degree $12$, and green $i=100$ with degree $8$). 
    In panel B, we integrate the Kuramoto dynamics of Eq.\,\eqref{eq:Akuramoto} on this network, and plot the norm of the phase perturbations vector $ \| \delta \theta(t) \|$ as a function of time $t$ and for different choices of the attacked node $i$. 
    In panel C, we plot the order parameter $ \rho(t)$, Eq.\,\eqref{eq:order}, as a function of time $t$ and for different choices of the attacked node $i$. Different curves correspond to different attacked nodes, all with the same budget $c = 1$. 
    The dashed black line denotes the start time of the attack. 
    The initial conditions of the oscillators are chosen randomly from the uniform distribution $[0, \ 0.1]$. 
    The plot also provides the degree of node $i$ as $\deg_i$.
    The highlighted background denotes the standard deviation over 20 realizations.
    The natural frequencies of the oscillators are randomly drawn from a Normal distribution with a mean of 0.1 and a standard deviation of 0.001.
    }

    \label{fig:Kuramoto_SF}
\end{figure*}




\begin{comment}
 
\subsection{Mutualistic Networks}

In this section, we study the effect of the intruder attack in mutualistic dynamical networks.
We consider the classical Lotka–Volterra model \cite{valdovinos2019mutualistic,bascompte2006asymmetric} that is described by the following equations:
\begin{align} \label{eq:lotka}
\begin{split}
    \dot{x}_i (t) & = r_i x_i (t) - s_i x_i^2 (t) + \sum_{j = 1}^m \alpha_{ij} x_i (t) y_j (t), \\
    \dot{y}_j (t) & = q_j y_j (t) - T_j y_j^2 (t) + \sum_{i = 1}^N \beta_{ji} y_j (t) x_i (t),
\end{split}
\end{align}
$i = 1, \hdots, N$, $j = 1, \hdots, m$ where $N$ and $m$ are the number of plant and animal species, respectively. 
Here, $x_i (t)$ and $y_i (t)$ are the population of plant $i$ and animal $j$ at time $t$, the positive scalars $r_i$ and $q_j$ are the growth rates, and the positive scalars $s_i$ and $T_j$ are the intraspecific competition rates.
The per capita effect of animal $j$ on plant $i$ is given by $\alpha_{ij}$ and the per capita effect of plant $i$ on animal $j$ is given by $\beta_{ji}$.

The simplification of having $r_i = r, \forall i$, $s_i = s, \forall i$, $q_j = q, \forall j$, $T_j = T, \forall j$, and $\alpha_{ij} = \alpha, \beta_{ji} = \beta, \ \forall i,j$ allows to analytically prove the existence of positive \textcolor{magenta}{stable?} fixed point
\begin{align}
\begin{split}
    x^* & = \dfrac{rT + m \alpha q}{Ts - m \beta N \alpha}, \quad
    y^* = \dfrac{qs + N \beta r}{Ts - m \beta N \alpha}.
\end{split}
\end{align}
if and only if all the parameters are positive and 
\begin{equation}
    \alpha \beta < \dfrac{sT}{nm},
\end{equation}
see \cite{bascompte2006asymmetric} for in depth details. 
Note if $N=m=1$, then the above inequality also guarantees the stability of the positive fixed point.
Note that in case of having strong connections ($\alpha \beta$ becoming large), the fixed point may no longer exist.




We denote the population of the intruder/adversarial species with $z(t)$ and denote its respective connection via the subscript $A$. 
We present our model in a general fashion in which the intruder may be either plant or animal.
The attacked mutualistic network equations becomes
\begin{align} \label{eq:lotka2}
\small
\begin{split}
    \dot{z} (t) & = w z(t)-u z^2(t) + \sum_{j=1}^m \alpha_{Aj} z(t) y_j (t) + \sum_{i=1}^N \beta_{Ai} z (t) x_i (t), \\
    \dot{x}_i (t) & = r_i x_i (t) - s_i x_i^2 (t) + \sum_{j = 1}^m \alpha_{ij} x_i (t) y_j (t) + \beta_{iA} z(t) x_i(t), \\
    \dot{y}_j (t) & = q_j y_j (t) - T_j y_j^2 (t) + \sum_{i = 1}^N \beta_{ji} y_j (t) x_i (t) + \alpha_{jA} z(t) y_j(t), 
\end{split}
\end{align}
where $w$ is the intruder's growth rate, $u$ is the intruder's competition rate.
For all $k$ and $k'$, the connection weights $\alpha_{kA}, \alpha_{Ak}, \beta_{k'A}, \beta_{Ak'}$ are all non-positive.

% We denote the intruder/adversarial species via the subscript $A$ and assume $c \alpha_{Ai} \leq 0$ and $c \alpha_{iA} \leq 0$ for all $i$ where $c \leq 0$ is the budget and $\alpha_{Aj} \geq 0, \alpha_{jA} \geq 0$.
% Also, we set $\sum_{j=1}^N \alpha_{Aj} = \sum_{j=1}^N \alpha_{jA} = 1$.
We proceed under the assumption that the system of equations in Eq.\,\eqref{eq:lotka2} has positive fixed points $x_i^*$, $y_j^*$, and $z^*$.
Linearization of Eq.\,\eqref{eq:lotka2} about the fixed points results in the Jacobian matrix $J = J_0 + P$ where 
% \begin{widetext}
% \begin{equation}
% \footnotesize
%     J_0 = \begin{bmatrix}
%     r_A-2 s_A x_A^* & 0 & 0 & \hdots & 0 \\ 
%     0 & r_1- 2 s_1 x_1^* + \sum_j \alpha_{1j} y_j^* & \alpha_{12} x_1^* & \hdots & \alpha_{1N} x_1^* \\ 
%     0 & \alpha_{21} x_2^* & r_2- 2 s_2 x_2^* + \sum_k \alpha_{2k} x_k^* + 2 \alpha_{22} x_2^* & \hdots & \alpha_{2N} x_2^* \\ 
%     \vdots & \vdots & \vdots & \ddots & \vdots \\ 
%     0 & \alpha_{N1} x_N^* & \alpha_{N2} x_N^* & \hdots & r_N- 2 s_N x_N^* + \sum_k \alpha_{Nk} x_k^* + 2 \alpha_{NN} x_N^* 
%     \end{bmatrix},
% \end{equation}
% \end{widetext}
\begin{widetext}
\begin{equation}
\footnotesize
    J_0 = \begin{bmatrix}
    w-2 u z^* & 0 & \hdots & 0 & 0 & \hdots & 0 \\ 
    0 & r_1 - 2s_1 x_1^* + \sum_j \alpha_{1j} y_j^* & \hdots & 0 & \alpha_{11} x_1^* & \hdots& \alpha_{1m} x_1^* \\ 
    \vdots & \vdots & \ddots & \vdots & \vdots & \ddots & \vdots \\ 
    0 & 0 & \hdots & r_N - 2s_N x_N^* + \sum_j \alpha_{nj} y_j^* & \alpha_{n1} x_N^* & \hdots& \alpha_{nm} x_N^* \\ 
    0 &  \beta_{11} y_1^* & \hdots& \beta_{1n} y_1^* & q_1 - 2T_1 y_1^* + \sum_i \beta_{1i} x_i^* & \hdots & 0 \\ 
    \vdots & \vdots & \ddots & \vdots & \vdots & \ddots & \vdots \\ 
    0 & \beta_{m1} y_m^* & \hdots& \beta_{mn} y_m^* & 0 & \hdots & q_m - 2T_m y_m^* + \sum_i \beta_{mi} x_i^* 
    \end{bmatrix},
\end{equation}
\begin{equation}
    P = \begin{bmatrix}
    \sum_j \alpha_{Aj} y_j^* + \sum_i \beta_{Ai} x_i^* & \beta_{A1} z^* & \hdots & \beta_{An} z^* & \alpha_{A1} z^* & \hdots & \alpha_{Am} z^* \\ 
    \beta_{1A} x_1^* & \beta_{1A} z^* & \hdots & 0 & 0 & \hdots & 0 \\ 
    \vdots & \vdots & \ddots & \vdots & \vdots & \ddots & \vdots \\ 
    \beta_{nA} x_N^* & 0 & \hdots & \beta_{nA} z^* & 0 & \hdots & 0 \\ 
    \alpha_{1A} y_1^* & 0 & \hdots & 0 & \alpha_{1A} z^* & \hdots & 0 \\ 
    \vdots & \vdots & \ddots & \vdots & \vdots & \ddots & \vdots \\ 
    \alpha_{mA} y_m^* & 0 & \hdots & 0 & 0 & \hdots & \alpha_{mA} z^* 
    \end{bmatrix}.
\end{equation}
\end{widetext}
The matrix $P$ has all non-positive entries and by construction, always has a positive eigenvalue.  \textcolor{red}{the matrix $P$ may have all negative eigenvalues as I found counter examples by generating matrices with the structure of matrix $P$. as an example:
\begin{equation}
\small
\begin{bmatrix}
-50.0472 & -0.8055 & -0.5767 & -0.1829 & -0.2399 & -0.8865 \\ 
-0.0287 & -47.1088 & 0 & 0 & 0 & 0 \\ 
-0.4899 & 0 & -5.9619 & 0 & 0 & 0 \\ 
-0.1679 & 0 & 0 & -68.1972 & 0 & 0 \\ 
-0.9787 & 0 & 0 & 0 & -4.2431 & 0 \\ 
-0.7127 & 0 & 0 & 0 & 0 & -7.1445 
\end{bmatrix}
\end{equation}
which has all negative eigenvalues.
}
When some or all the connection weights $\alpha_{kA}, \alpha_{Ak}, \beta_{k'A}, \beta_{Ak'}$ become large in magnitude, then the eigenvalues of the Jacobian matrix $J = J_0 + P$ will be dominated by those of the matrix $P$; thus destroying the stability of the positive fixed point.






% The matrix $J_0$ denotes the Jacobian before the attack (when the budget $c =0$), and $P$ denotes the effect of the attack on the Jacobian $J_0$. 


% The matrix $P$ has two non-zero eigenvalues $\lambda_{1,2} = \pm \sqrt{x_A^* \sum_{k=1}^N \alpha_{Ak} \alpha_{kA} x_k^*}$ and all other eigenvalues are zero. 
% For large budgets ($| c | \gg 0$), the eigenvalues of $J$ are dominated by those of $c P$.
% Thus, we conclude for a large budget, the Jacobian of the Lotka–Volterra model, evaluated about its fixed point, will have mixed signs eigenvalues, and as  result, the stability of the fixed point will be lost due to the intruder attack.



   
\end{comment}



\begin{remark}
    Our empirical results indicate that in the case of more than one intruder, the optimal attack, i.e., the attack that minimizes the algebraic connectivity, is the one with all the budget concentrated in one link from one of the intruders toward the node with the lowest indegree and with all the other intruders that remain disconnected.
    The particular choice of the intruder makes no difference, since all are identical before the attack.
    An extensive discussion on this point is provided in the Supplementary Note 2.
\end{remark}
\begin{remark}
    We have performed a numerical study over our randomly generated scale-free network, Fig.\ref{fig:Kuramoto_SF}A.
    We have seen that the following values of the budget can be considered `small' and `large' in our asymptotic analysis:
    \begin{itemize}
        \item the small budget should be $c < 10^{-2}$ (independent of the attacked node), 
        \item the large budget should be $c > 0.1 L_{ii}^2$ (dependent on the attacked node),
    \end{itemize}
    where $L_{ii}$ is the indegree of the attacked node.
    The details of our analysis are available in Supplementary Note 8.
\end{remark}


In Supplementary Note 9, we compare the effects of an intruder attack on two leaf nodes (nodes with degree 1): one connected to the hub node (the node with highest degree) and one not connected to the hub. By examining the dynamics immediately after the attack, we do not observe any difference between the two cases in their transient response.
\color{black}

\section{Discussion} \label{sec:conclusions}


 
 This work identifies a fundamental vulnerability of complex dynamical networks: a single node connected through adversarial links to one or more other nodes is sufficient to destabilize the network dynamics, regardless of the network size. Such an attack provides a generic mechanism for disrupting consensus in opinion dynamics, synchronization in oscillator networks, and desired formations in coupled autonomous systems. Furthermore, we find that targeting a node with the smallest indegree, {defined as the sum of the weights of all incoming connections}, typically destabilizes the network in the shortest time. In contrast, attacking hub nodes generally results in a slower onset of instability. Intuitively, this can be explained as follows: hubs respond less to the coupling with the intruder due to a large number of in-neighbors, while low-indegree nodes respond more.

 

The practical implementation of negative (adversarial) interactions depends on the application domain, but is in general well grounded across a wide range of systems. 
In technological and engineering networks, the distinction between cooperative and adversarial interactions often reduces to the sign of a feedback gain. In standard Laplacian coupling, the interaction acting on node $i$ takes the form $g (x_{\mathrm{ref}} - x_i)$, where $x_{\mathrm{ref}}$ is a reference state, $x_i$ is the state of node $i$, and $g$ is a scalar gain. Cooperative interactions correspond to $g > 0$, whereas adversarial interactions correspond to $g < 0$. Thus, changing the sign of the feedback gain naturally switches between the two regimes. In practice, such sign changes may arise from faulty or malicious agents, incorrect controller design, or cyber-attacks that invert transmitted signals. In power grid applications, for example, adversarial coupling can be realized through devices enforcing anti-phase behavior, as discussed in Subsection II.C.
In biological systems, negative interactions are ubiquitous and arise intrinsically. Examples include inhibitory synaptic coupling in neuronal networks, competitive interactions in ecological systems, and repression mechanisms in gene regulatory networks. In these settings, adversarial interactions are not externally imposed, but are a fundamental component of the system dynamics. Within this perspective, our framework captures the effect of introducing or amplifying inhibitory interactions in otherwise cooperative networks.
In social systems, negative coupling emerges through antagonistic or contrarian behavior. For instance, in opinion dynamics models, some individuals may systematically oppose the state of their neighbors (so-called ``contrarians''), effectively inducing negative interactions. Likewise, misinformation or adversarial agents in online networks can disrupt consensus by promoting opposing views. These mechanisms provide realistic pathways through which adversarial interactions can arise and influence collective dynamics.


 
 Our results are in stark contrast to the expectation that hubs are the most critical nodes of a network \cite{artime2024robustness}, and thus the ones that need to be more closely protected from attacks. 
 {This assumes particular relevance in the case of scale-free networks \cite{barabasi1999emergence}.
 A central dogma in the field of complex networks is that scale-free networks are stable against random failures (removal) of nodes but are vulnerable, in terms of structural integrity (i.e., the size of the strongly connected component) to attacks targeting the hubs. %This is well known and well characterized, for example,  in the study of percolation on networks \cite{artime2024robustness}.
We show a significantly different picture for the case of dynamical attacks considered here: while the structure of scale-free networks is more vulnerable to attacks targeting the hubs \cite{artime2024robustness}, their dynamics is more vulnerable to attacks targeting nodes with low-indegree.
In summary, while hubs exhibit high vulnerability in terms of structural integrity, they are somewhat protected in dynamical contexts, leading to a situation where low-degree nodes are most vulnerable. Reference \cite{liu2011controllability} had also shown that the structural controllability of a network was enhanced when driver nodes are not connected to high-degree nodes. {Reference  \cite{menichetti2014network} had found that the structural controllability of a network depends strongly on the fraction of
low in-degree and low out-degree nodes. } However, destabilizing the network dynamics does not require controllability, and here we provide simple sufficient conditions to induce instability of this dynamics. Our work complements a large literature on the importance of low-degree nodes and weak ties \cite{granovetter1973strength}; the relevance of low-degree nodes, under specific circumstances, is also acknowledged in the literature on percolation dynamics on networks  \cite{artime2024robustness} and {the importance of weak ties on stability for predator–prey networks has been pointed out in \cite{allesina2012stability}.} 
 }
 % {Add a sentence about the hubs being the most important nodes being the hallmark of network science?}


 Our work complements the large literature that focuses on conditions to ensure the stability of the network dynamics, see e.g., \cite{boccaletti2006complex}, while we focus on instability and on characterizing the severity of such instability. We investigate the fundamental question of how desired dynamical states can be disrupted by the purposeful insertion of one or a few intruder agents. 
Our intruders are structurally similar to the other network agents, which is consistent with the case of cyber-attacks. 
Our results are initially derived for the case of the linear consensus problem but have implications in a broad variety of network dynamics, including opinion dynamics, synchronization, formation control, and power grid balancing. We also study the case of the Kuramoto model, % on networks, 
which is relevant to synchronization of biological systems. That having been said, we do not claim that our results apply to all the possible realizations of Eq.\ \eqref{general}, but to a variety of cases of interest.
%Complex technological systems, such as coupled autonomous systems, power grids, and sensor networks are distributed cyber-physical systems. Intruder attacks happen at the cyber level (software), while leaving the physical level (hardware) unaffected.  According to our definition, an attacker has the same individual (uncoupled) dynamics as the rest of the nodes, but can intentionally establish adversarial interactions with the other nodes in order to destabilize the network dynamics. %Their distributed structure makes them prone to attacks that may propagate from one or a few targeted nodes to the network. 
 %We see that, independent of how large the network is, a single intruder coupled to a single node in a network is sufficient to completely destabilize the network dynamics. This points out to a main vulnerability of many (if not all) complex networks.

 Our work on intruder attacks in networked systems intersects with the domain of network interdiction, which has been extensively studied in operations research and optimization contexts. Smith and Song's comprehensive survey on network interdiction \cite{smith2020survey} provides a systematic overview of how adversaries can strategically disrupt network operations by targeting critical components. While traditional network interdiction focuses primarily on structural disruptions such as removing nodes or edges to maximize the disruption of flow, connectivity, or shortest paths \cite{Wood1993,IsraeliWood2002}, our approach examines dynamical disruptions through adversarial connections that destabilize the network's equilibrium states. This distinction is crucial: classical interdiction models often assume a static network structure where the adversary's goal is to disconnect components or increase path lengths \cite{SmithSong2020,CappaneraScaparra2011}, whereas our framework considers how a single intruder node can exploit the network's dynamical properties to induce instability while leaving the structural connectivity intact. %Recent extensions of interdiction theory to dynamic and stochastic settings \cite{Zenklusen2010,PanMorton2008} share conceptual similarities with our approach, particularly in modeling adaptive adversaries and time-evolving vulnerabilities. However, these works typically focus on discrete optimization problems rather than continuous dynamical systems. 
 Our results complement the interdiction literature by demonstrating that in cyber-physical systems, power grids, and consensus networks, dynamical vulnerability differs fundamentally from structural vulnerability: targeting low-indegree nodes rather than high-degree hubs identified as critical in classical interdiction models \cite{HolmeKimYoonHan2002} produces the most severe instabilities. This finding suggests that protection strategies derived from structural interdiction theory may be insufficient for dynamical networks, necessitating new defense mechanisms that account for both topological and dynamical characteristics.


We show the emergence of universal scaling properties, with small perturbations causing uniform responses and larger ones exposing specific vulnerabilities, especially at the low-indegree nodes.
When applied to real complex networks, we see that it may not always be easy to characterize the rate of instability produced by the attack in the general case of non-balanced digraphs. 
However, our theory (see also Fig.\,\ref{fig:realnet}) shows that the rate of instability generated by an attack on a single node averaged over all $N$ networks nodes, scales exactly as $-N^{-1}$,  for the case that the adversarial connection is unidirectional. This indicates that larger networks are more robust, on average, against intruder attacks.

Although our work focuses on the case that an intruder node is added to an existing network, it can be generalized to the case that an existing network node becomes an intruder. In that case, the strategy that maximizes the effects of an attack is for the intruder to target only one node with the lowest indegree, where the indegree is computed by counting all the nodes in the network, except for the intruder. 

An obvious countermeasure against intruder attacks is to disconnect the attacker(s) from the network. This is consistent, for example, with the load shedding strategy, commonly implemented for power grids in the case of a power imbalance. %, for which a typical remedy is to shed some of the loads. 
In general,   shedding the part of the network that includes the intruder provides a viable solution against intruder attacks, which is easier to do when the attacker is connected to the rest of the network by a single link compared to the case of multiple links. This observation highlights a trade-off between the severity of an attack and the recoverability from such attack.

{Extending the present framework to higher-order (many-body) interactions represents a natural and promising direction for future work. Recent studies have shown that such interactions can significantly modify synchronization properties, leading to global topological synchronization in some cases, but also to multistability and abrupt desynchronization in others \cite{carletti2023global, gallo2022synchronization, skardal2019abrupt, millan2020explosive}. This indicates that, depending on their structure, higher-order interactions may either promote coherence or give rise to additional instability mechanisms. In this context, we expect them to shift stability thresholds and modulate the system’s sensitivity to perturbations, without altering the underlying mechanism.}

Finally, a word of caution. While our results are rigorous for the case of linear dynamics (e.g., consensus), they are mostly based on extensive numerical calculations in the case of nonlinear dynamics.  We believe more work is needed to fully characterize the effects of intruder attacks in nonlinear networks, although our work provides strong indication that in many cases of interest, the rate of instability observed for linear networks is inherited by their nonlinear counterparts. 




\begin{comment}
{Biological systems can also be subject to intruder attacks, such as in the case of an invasive species introduced in an ecological network.
However, in the biological world, destabilizing effects are often desirable, as in the case of preventing abnormal synchronization in neurological diseases such as Parkinson’s disease or epilepsy \cite{louzada2012suppress}, or breaking down undesirable systemic behavior in other contexts. It is then possible that attacks 
are engineered into these systems to prevent them from relaxing onto an orderly dynamical state.}
\end{comment}



%\newpage

\section{Methods} \label{sec:methods}

\begin{comment}
\subsection{Proof of Proposition \ref{prop1}} \label{sec:prop1}

We first prove the upper bound $\lambda_{\min} (L_{aug}) < c$.
For all vectors $\bx$ of norm 1, it follows $\lambda_{\min} (L_{aug}) \leq \bx^\top L_{aug} \bx$.
By setting $\bx_1 = [1 \  0 \ \hdots 0]^\top$, it results $\bx_1^\top L_{aug} \bx_1 = c \leq 0$.
It thus follows $\lambda_{\min} (L_{aug}) \leq c$ where the equality is satisfied if $\bx_1$ is an eigenvector which only occurs if $c = 0$ and $\bb = \pmb{0}$.

Now we prove the lower bound. Based on the Gershgorin circle theorem, at least one of the eigenvalues of the matrix $L_{aug}$ lies in the dist centered at $c$ with radius $c$. Since we know $\lambda_{\min} (L_{aug}) \leq c$ from above, it thus follows $2c \leq \lambda_{\min} (L_{aug}) \leq c \leq 0$. \hfill \qed


\subsection{Proof of Proposition \ref{prop2}} \label{sec:prop2}

{To be added.}

\end{comment}




\begin{comment}

%%%%%%%%%%%%%%% old transverse reactivity section %%%%%%%%%%%%%%%%%%%%%%%%%%%%%

\subsection{Transverse reactivity of stable and unstable consensus dynamics}

\textcolor{red}{*** I wonder whether we could replace the Laplacian with the negative Laplacian. That would simplify things in this section of the methods, see Remark 2}
{I think it is better to talk about this comment.}

\begin{definition}{Spectral abscissa of a matrix.} Given a square matrix $M \in R^{N \times N}$, we define the spectral abscissa as the greatest real part of the matrix spectrum,
\begin{equation}
    \alpha(-M)=\max_i \Re(-\lambda_i),
\end{equation}
where $\lambda_1,\lambda_2,...,\lambda_N$ are the possibly complex eigenvalue of the matrix $M$ and $\Re(\cdot)$ returns the real part of its argument. 
\end{definition}

\begin{definition}{Reactivity of a matrix \cite{Neubert1997ALTERNATIVES}.} Given a square matrix $M \in R^{N \times N}$, we define the reactivity,
\begin{align} 
\label{eq:reactivity}
\begin{split}
    \xi\left( -M \right) & = \max_{  \bX \neq 0} \frac{{\bX}^\top \left( \frac{ -M-M^\top}{2} \right) \bX }{{ \bX}^\top  \bX} \\
    & =  \lambda_{\max} \left( \frac{-M-M^\top}{2} \right) \\
    & =  \lambda_{\max}(S),
\end{split}
\end{align}
where the matrix $S$ is the symmetric part of the matrix $-M$ and $\lambda_{\max}(\cdot)$ indicates the largest eigenvalue of the matrix in its argument.
\end{definition}

\begin{definition}{Transverse reactivity of a Laplacian matrix \cite{nazerian2024efficiency}.} Given a Laplacian matrix $L \in R^{N \times N}$, we define the transverse reactivity,
\begin{align} \label{eq:rhostar}
\begin{split}
    \xi^*\left( -L \right) & = \max_{\bX \neq 0, \bX \perp \bone} \frac{{\bX}^\top \left( \frac{ -L-L^\top}{2} \right) \bX }{{ \bX}^\top  \bX}.
\end{split}
\end{align}
\end{definition}

\begin{remark}
    The two definitions above are provided for the negative of a matrix because our single-integrator consensus dynamics $\eqref{consensus}$ is written in terms of the negative of the Laplacian matrix.
\end{remark}



The discussion that follows is for the case of the single-integrator consensus dynamics \eqref{consensus},
%\begin{equation}  \label{eq:lingen}
%   \dot{\bX}(t) = -L \bX(t),
%\end{equation}
in the state vector $\bX(t)=[x_1(t),x_2(t),....,x_N(t)]$.
We write the eigenvalue equation $L \bv_i=\lambda_i \bv_i$, where $\lambda_i$  ($\bv_i$) are the possibly complex eigenvalues (eigenvectors) of the Laplacian matrix $L$, $i=1,...,N$. By the property that the rows of the Laplacian all sum to zero, there exists one eigenvalue of the Laplacian equal to zero, with the corresponding eigenvector having entries that are all the same. We index this eigenvalue $j$.
 Without loss of generality, we set $\Re(\lambda_1(L)) \leq \Re(\lambda_2(L)) \leq ... \leq \Re(\lambda_N(L))$,
where $\Re(\cdot)$ returns the real part of its argument. When the Laplacian matrix is proper, $j=1$ and $\Re(\lambda_1(L))=0$, which results in an asymptotically stable consensus dynamics \eqref{consensus}. Here, we remove the assumption that the Laplacian is proper and present a general theory of transient and asymptotic stability for \eqref{consensus}.

Any point such that
\begin{equation} \label{cm}
x_1=x_2=...=x_N
\end{equation}
is an equilibrium for the dynamics \eqref{consensus}. The set of points that satisfy \eqref{cm} define the consensus manifold. We are interested in whether the consensus manifold is asymptotically stable or unstable and in characterizing the maximum rate of transient instability away from this manifold. Thus in what follows, we focus on the transverse consensus dynamics,
\begin{equation}
\label{transverse}
    \dot{\bZ}(t) = -A \bZ(t)
\end{equation}
$\bZ \in \mathbb{R}^{N-1}=V^\top \bX$, the matrix $A \in \mathbb{R}^{N-1 \times N-1} =V^\top L V$
and $V \in \mathbb{R}^{N \times N-1}$ is an orthonormal basis for the null subspace of $\bone^\top$, i.e., $V$ is a matrix whose columns are normal and orthogonal to one another and have zero column sums.

\begin{proposition}
The matrix $A$ has the same spectrum as the Laplacian matrix $L$, except for one eigenvalue of the matrix $L$ which is equal zero, i.e., $\mathcal{S}(L) = \mathcal{S}(A) \cup \{ 0 \}$ where $\mathcal{S}(\cdot)$ returns the spectrum of the matrix in its argument.
\end{proposition}



Next we comment on the relevance of both the spectral abscissa and the reactivity on the transverse consensus dynamics \eqref{transverse}. The asymptotic rate of $\bZ(t)$ in \eqref{transverse} is given by $\alpha(-A)$, i.e., the spectral abscissa of $-A$.
In general, $\alpha(-A)$ can be either positive or negative. 
On the other hand the instantaneous rate of growth of the norm of the vector $\bZ$ is given by the reactivity,
\begin{equation} \label{eq:reactivity0}
    \xi\left( -A \right) :=  \max_{\|  \bZ\| \neq 0} \left[ \frac{1}{\|  \bZ \|} \frac{d \|  \bZ\|}{d t} \right]
\end{equation}
Thus a positive (negative) reactivity $\xi(-A)$ indicates that the norm of the state $\bZ$ tends to grow (shrink) in the limit of $t \rightarrow 0$. 

\begin{proposition}
    The transverse reactivity of the Laplacian matrix in \eqref{consensus} is equal to the reactivity of the matrix $A$ in \eqref{transverse}, that is, $\xi^*(-L) \equiv \xi(-A)$.
\end{proposition}
\begin{proof}
    
    The reactivity in Eq.\,\eqref{eq:reactivity0} is
    \begin{align*}
        \xi\left( -A \right) & :=  \max_{\|  \bZ\| \neq 0} \left[ \frac{1}{\|  \bZ \|} \frac{d \|  \bZ\|}{d t} \right] \\
        & = \max_{\|  \bZ\| \neq 0} \left[ \frac{1}{\sqrt{\bZ^\top \bZ}} \frac{d \sqrt{\bZ^\top \bZ}}{d t} \right] \\
        & = \max_{\|  \bZ\| \neq 0} \left[ \frac{1}{\sqrt{\bZ^\top \bZ}} \frac{\dot{\bZ}^\top \bZ + \bZ^\top \dot{\bZ}}{2\sqrt{\bZ^\top \bZ}} \right] \\
        & = \max_{\|  \bZ\| \neq 0} \left[ \frac{\bZ^\top (-A - A^\top) \bZ}{2 \bZ^\top \bZ} \right].
    \end{align*}
    Also, $V^\top V = I$, $L = V A V^\top$ and $\bX = V \bZ$ which indicates $\bX \perp \bone$ since $\bX \in \text{range}(V)$, so
    {\footnotesize
    \begin{align*}
        \xi\left( -A \right) & = \max_{\|  \bZ\| \neq 0} \left[ \frac{\bZ^\top  \left(-(V^\top V) A (V^\top V) - (V^\top V) A^\top (V^\top V)\right) \bZ}{2 \bZ^\top \bZ} \right] \\
        & = \max_{\|  \bZ\| \neq 0} \left[ \frac{(\bZ^\top V^\top)  (- L - L^\top) (V\bZ)}{2 (\bZ^\top V^\top) (V\bZ)} \right] \\
        & = \max_{ \bX \neq 0, \bX \perp \bone} \left[ \frac{\bX^\top  (- L - L^\top) \bX}{2 \bX^\top \bX} \right] \\
        & = \xi^* (-L).
    \end{align*}}
    That concludes the proof.
\end{proof}





Thus in what follows, we refer to $\xi(-A)$ in Eq.\ \eqref{eq:reactivity0} as either the reactivity of \eqref{transverse} or the transverse reactivity of \eqref{consensus}, interchangeably.








\begin{remark}
    We assume the network has at least one directed spanning tree which ensures $\alpha(-A)\neq0$.
\end{remark}


\begin{theorem}
    The reactivity of the dynamics \eqref{transverse} $\xi(-A) \geq \alpha(-A)$.
\end{theorem}
\begin{proof}
Consider the reactivity in \eqref{eq:reactivity},
\begin{align} \label{eq:lambda}
\begin{split}
    f (-M) = \max_{ \bZ \neq 0} \frac{{\bZ}^\top S \bZ }{{ \bZ}^\top \bZ} & \geq { \bv_1}^\top S \bv_1 \\
    & =   { \bv_1}^\top \left(\frac{-M -M^\top}{2} \right) \bv_1 \\
    & =  \frac{1}{2}(-\bv_1^\top M \bv_1 - \bv_1^\top M^\top \bv_1 ) \\
    & = \frac{1}{2}(\bv_1^\top \alpha(-M) \bv_1  + \alpha(-M) \bv_1^\top \bv_1 ) \\
    & = \alpha(-M) 
\end{split}
\end{align}
This concludes the proof. 
\end{proof}

The above theorem is important as it states that the instantaneous rate of growth of the state $\bZ$ in \eqref{transverse} can be either larger than or equal to the asymptotic rate of grown of $\bZ$. This is true in both the cases of asymptotically stable and unstable dynamics. 
If $\alpha(M)>0$, we then see that $f(M) \geq \alpha(M) >0$. This indicates that for an unstable system, the rate of instantaneous growth is positive and exceeds or is equal to the rate of asymptotic growth.



In our previous work \cite{nazerian2023single}, we studied the entire consensus dynamics Eq.\,\eqref{consensus} under the assumption the consensus dynamics was stable, i.e., $\alpha(-A) < 0$.

In this case, the trajectories of the system states eventually approach the origin (asymptotic stability), while the magnitude of the states may grow at transient times due to positive reactivity, which is not desirable in various situations, such as linearized dynamics about a fixed point. This transient growth may result in the invalidity of the linearization assumption.

Unlike previous work by us and others \cite{Farrell1996Generalized,Neubert1997ALTERNATIVES,Hennequin2012Nonnormal,Tang2014Reactivity,Biancalani2017Giant,Asllani2018Structure,MUOLO2019Patterns,Gudowska2020From,Lindmark2021Centrality,Duan2022Network,nazerian2023single,Nazerian2023Reactability},  here we are mostly concerned with the reactivity of the transverse dynamics, i.e., Eq.\,\eqref{transverse} in the case that $\alpha(-A) > 0$, i.e., the system \eqref{consensus} is asymptotically unstable.





The choice of $\bZ$ that maximizes \eqref{eq:reactivity} is the eigenvector $\bw$ of the matrix $S$ associated with the eigenvalue $\lambda_{\max}(S)$.
Note that the reactivity
is computed by maximizing over $\bZ$, thus it provides a worst-case instantaneous rate of growth for $\| \bZ \|$. In the case that $\bZ$ is generically chosen, it will have a non-zero component along the  eigenvector $\bw$, thus the worst-case rate of growth for $\| \bZ \|$ will be observed with probability 1.





{
It is inferred that the inequality in \eqref{eq:lambda} is satisfied with the equal sign when the left and right eigenvectors of the matrix $M$ associated with the eigenvalue $\lambda_1$ coincide. 
This means that {$\bX^* =  \bv_1$} becomes the maximizer in \eqref{eq:lambda}. }



To illustrate the relation $f (M) \geq \alpha (M)$, we provide a numerical example of stable and unstable dynamics.
First, we randomly generate the proper Laplacian matrix
\begin{equation} \label{Lproper}
    L = \begin{bmatrix}
    1 & -1 & 0 & 0 & 0 \\ 
    0 & 1 & 0 & 0 & -1 \\ 
    -1 & -1 & 2 & 0 & 0 \\ 
    0 & 0 & 0 & 0 & 0 \\ 
    0 & 0 & 0 & -1 & 1 
    \end{bmatrix}
\end{equation}
with all eigenvalues with non-negative real-part and only one zero eigenvalue.
The matrix $A$ in Eq.\,\eqref{transverse} corresponding to the Laplacian matrix $L$ in \eqref{Lproper} is
\begin{equation}
    A = \begin{bmatrix}
    1.5427 & -0.1809 & 0.2337 & -0.9045 \\ 
    -0.1483 & 2.1281 & 0.5427 & 0.4045 \\ 
    0.5427 & -0.1809 & 0.2337 & 0.0955 \\ 
    0.5427 & -0.1809 & -0.7663 & 1.0955 
    \end{bmatrix}.
\end{equation}
Here, $f(-A) = -0.0250 \geq \alpha(-A) = -1$. 
We simulate Eq.\,\,\eqref{transverse} from the initial condition $\bZ(0) = [-0.2187,\ -0.1115,\ 0.9320 ,\ 0.2667 ]$ and plot the norm of the transverse perturbation $\| \bZ(t) \|$ in Fig.\,\ref{fig:consensus_rate} A in log scale.
As evident by the slope of $\| \bZ(t) \|$ vs $t$, the transient rate of decay $f(-A)$ is larger than the asymptotic rate of decay $\alpha(-A)$.

Next, we add an adversarial agent to the above Laplacian matrix with a budget $c = -0.001$ and obtain
\begin{equation} \label{Lnonproper}
    L_{aug} = \begin{bmatrix}
    1 & -1 & 0 & 0 & 0 & 0 \\ 
    0 & 1 & 0 & 0 & -1 & 0 \\ 
    -1 & -1 & 2 & 0 & 0 & 0 \\ 
    0 & 0 & 0 & -0.001 & 0 & 0.001 \\ 
    0 & 0 & 0 & -1 & 1 & 0 \\ 
    0 & 0 & 0 & 0.001 & 0 & -0.001 
    \end{bmatrix}
\end{equation}
which has mixed signed real-part eigenvalues.
For the Laplacian matrix \eqref{Lnonproper}, the corresponding matrix $A_{aug}$ in Eq.\,\eqref{transverse} is
\begin{equation}
    A_{aug} = \begin{bmatrix}
    1.4923 & -0.1527 & 0.2024 & -0.9160 & 0.0840 \\ 
    -0.2178 & 2.1372 & 0.4923 & 0.3739 & 0.3739 \\ 
    0.4923 & -0.1527 & 0.2014 & 0.0840 & 0.0850 \\ 
    0.4923 & -0.1527 & -0.7976 & 1.0840 & 0.0840 \\ 
    0.4923 & -0.1527 & 0.2034 & 0.0840 & 0.0830 
    \end{bmatrix}.
\end{equation}
Here, $f(-A_{aug}) = 0.0874 \geq \alpha (-A_{aug}) = 0.002$.
We simulate Eq.\,\,\eqref{transverse} from the initial condition $\bZ(0) = [0.1202,\  0.1937,\ 0.1545,\ 0.8538,\ 0.4418]$ and plot the norm of the transverse perturbation $\| \bZ(t) \|$ in Fig.\,\ref{fig:consensus_rate} B in log scale.
As evident by the slope of $\| \bZ(t) \|$ vs $t$, the transient rate of decay $f(-A)$ is larger than the asymptotic rate of decay $\alpha(-A)$.


%%%%%%%%%%%%%%% End of old transverse reactivity section %%%%%%%%%%%%%%%%%%%%%%%%%%%%%

\end{comment}



\subsection{Transverse reactivity of stable and unstable consensus dynamics}


The first part of this subsection discusses general linear dynamical systems.
The second part of this subsection focuses on linear consensus dynamics as a special case of general linear dynamical systems. 
Consider a linear dynamical system
\begin{equation} \label{eq:lti}
    \dot{\bx} (t) = M \bx(t)
\end{equation}
where $\bx(t) \in \mathbb{R}^N$ is the vector of system states and the square matrix $M \in \mathbb{R}^{N \times N}$.
\begin{definition}{Spectral abscissa of a matrix.} Given a square matrix $M \in R^{N \times N}$, we define the spectral abscissa as the greatest real part of the matrix spectrum,
\begin{equation}
    \alpha(M)=\max_i \Re(\lambda_i),
\end{equation}
where $\lambda_1,\lambda_2,...,\lambda_N$ are the possibly complex eigenvalue of the matrix $M$ and $\Re(\cdot)$ returns the real part of its argument. 
\end{definition}
The origin is a stable fixed point for the system in Eq.\,\eqref{eq:lti} if and only if $\alpha(M) < 0$.

\begin{definition}{Reactivity of a linear dynamical system \cite{Neubert1997ALTERNATIVES}.} The reactivity for the system in Eq.\,\eqref{eq:lti} is defined as,
\begin{align} 
\label{eq:reactivity}
\begin{split}
    \xi\left( M \right) & := \max_{\|  \bX\| \neq 0} \left[ \frac{1}{\|  \bX \|} \frac{d \|  \bX\|}{d t} \right] \\
    & = \max_{  \bX \neq 0} \frac{{\bX}^\top \left( \frac{ M+M^\top}{2} \right) \bX }{{ \bX}^\top  \bX} \\
    & =  \lambda_{\max} \left( \frac{M+M^\top}{2} \right) =:  \lambda_{\max}(S),
\end{split}
\end{align}
where the matrix $S = (M+M^\top)/2$ is the symmetric part of the matrix $M$ and $\lambda_{\max}(\cdot)$ indicates the largest eigenvalue of the symmetric matrix in its argument.
\end{definition}
The reactivity measures the maximal rate of instantaneous growth/decay of the norm of the state vector. 
If $\xi\left( M \right) < 0$ ($\xi\left( M \right) >0$) the norm of the state vector will decay (may grow) instantaneously.
The choice of $\bX$ that maximizes \eqref{eq:reactivity} is the eigenvector $\bv_1$ of the matrix $S$ associated with the eigenvalue $\lambda_{\max}(S)$.
Note that the reactivity is computed by maximizing over $\bX$, thus it provides a worst-case instantaneous rate of growth for $\| \bX \|$. In the case that $\bX$ is generically chosen, it will have a non-zero component along the eigenvector $\bv_1$, thus the worst-case rate of growth for $\| \bX \|$ will be observed with probability 1.
% \begin{remark}
%     The two definitions above are provided for the negative of a matrix because our single-integrator consensus dynamics $\eqref{consensus}$ is written in terms of the negative of the Laplacian matrix.
% \end{remark}

\begin{theorem} \label{theo}
    The reactivity of the dynamics \eqref{eq:lti} is greater or equal to the spectral abscissa of the matrix $M$, i.e., $\xi(M) \geq \alpha(M)$.
\end{theorem}
\begin{proof}
We set $S = (M+M^\top)/2$ and write the reactivity of the dynamics \eqref{eq:lti},
\begin{align} \label{eq:lambda}
\begin{split}
    \xi(M) = \max_{ \bZ \neq 0} \frac{{\bZ}^\top S \bZ }{{ \bZ}^\top \bZ} & \geq { \bv_1}^\top S \bv_1 \\
    & =   { \bv_1}^\top \left(\frac{M+M^\top}{2} \right) \bv_1 \\
    & =  \frac{1}{2}(\bv_1^\top M \bv_1 + \bv_1^\top M^\top \bv_1 ) \\
    & = \frac{1}{2}(\bv_1^\top \alpha(M) \bv_1  + \alpha(M) \bv_1^\top \bv_1 ) \\
    & = \alpha(M) 
\end{split}
\end{align}
This concludes the proof. 
\end{proof}

Theorem \ref{theo} is important as it states that the instantaneous rate of growth of the state $\bZ$ in \eqref{eq:lti} can be either larger than or equal to the asymptotic rate of growth of $\bZ$. This is true in both the cases of asymptotically stable and unstable dynamics. 
If $\alpha(M)>0$, we then see that $\xi(M) \geq \alpha(M) >0$. This indicates that for an unstable system, the rate of instantaneous growth is positive and exceeds or is equal to the rate of asymptotic growth.



It is inferred that the inequality in \eqref{eq:lambda} is satisfied with the equal sign when the left and right eigenvectors of the matrix $M$ associated with the eigenvalue $\lambda_1$ coincide. 
This means that $\bZ^* =  \bv_1$ becomes the maximizer in \eqref{eq:lambda}. 



The discussion that follows is for the single-integrator consensus dynamics in Eq.\,\eqref{consensus}, $\dot{\bX}(t) = -L \bX(t)$, which is a special case of the general linear dynamical system in Eq.\,\eqref{eq:lti}.
%\begin{equation}  \label{eq:lingen}
%   \dot{\bX}(t) = -L \bX(t),
%\end{equation}
% The state vector in the consensus dynamics is $\bX(t)=[x_1(t),x_2(t),....,x_N(t)]$.

We write the eigenvalue equation $L \bv_i=\lambda_i \bv_i$, where $\lambda_i$  ($\bv_i$) are the possibly complex eigenvalues (eigenvectors) of the Laplacian matrix $L$, $i=1,...,N$. By the property that the rows of the Laplacian all sum to zero, there exists one eigenvalue of the Laplacian equal to zero, with the corresponding eigenvector having entries that are all the same. We index this eigenvalue $j$.
 Without loss of generality, we set $\Re(\lambda_1(L)) \leq \Re(\lambda_2(L)) \leq ... \leq \Re(\lambda_N(L))$,
where $\Re(\cdot)$ returns the real part of its argument. When the Laplacian matrix is proper, $j=1$ and $\Re(\lambda_1(L))=0$, which results in an asymptotically stable consensus dynamics Eq.\,\eqref{consensus}. Here, we remove the assumption that the Laplacian is proper and present a general theory of transient and asymptotic stability for Eq.\,\eqref{consensus}.

% \begin{assumption}
%     The network represented by the Laplacian matrix $L$ in Eq.\,\eqref{consensus} has at least one directed spanning tree.
% \end{assumption}

Any point such that
\begin{equation} \label{cm}
x_1=x_2=...=x_N
\end{equation}
is an equilibrium for the dynamics \eqref{consensus}. The set of points that satisfy \eqref{cm} define the consensus manifold. We are interested in whether the consensus manifold is asymptotically stable or unstable and in characterizing the maximum rate of transient instability away from this manifold. Thus in what follows, we focus on the transverse consensus dynamics,
\begin{equation}
\label{transverse}
    \dot{\bZ}(t) = -A \bZ(t)
\end{equation}
$\bZ \in \mathbb{R}^{N-1}=V^\top \bX$, the matrix $A \in \mathbb{R}^{N-1 \times N-1} =V^\top L V$
and $V \in \mathbb{R}^{N \times N-1}$ is an orthonormal basis for the null subspace of $\bone^\top$, i.e., $V$ is a matrix whose columns are normal and orthogonal to one another and have zero column sums.
\begin{remark}
The matrix $A$ has the same spectrum as the Laplacian matrix $L$, except for one eigenvalue of the matrix $L$ which is equal zero, i.e., $\mathcal{S}(L) = \mathcal{S}(A) \cup \{ 0 \}$ where $\mathcal{S}(\cdot)$ returns the spectrum of the matrix in its argument.
\end{remark}

In our previous work \cite{nazerian2023single}, we studied the entire consensus dynamics Eq.\,\eqref{consensus} under the assumption the consensus dynamics was stable, i.e., $\alpha(-A) < 0$.
In this case, the trajectories of the system states eventually approach the origin (asymptotic stability), while the magnitude of the states may grow at transient times due to positive reactivity, which is not desirable in various situations, such as linearized dynamics about a fixed point. This transient growth may result in the invalidity of the linearization assumption.
Unlike previous work by us and others \cite{Farrell1996Generalized,Neubert1997ALTERNATIVES,Hennequin2012Nonnormal,Tang2014Reactivity,Biancalani2017Giant,Asllani2018Structure,MUOLO2019Patterns,Gudowska2020From,Lindmark2021Centrality,Duan2022Network,nazerian2023single,Nazerian2023Reactability},  here we are mostly concerned with the reactivity of the transverse dynamics, i.e., Eq.\,\eqref{transverse} in the case that $\alpha(-A) > 0$, i.e., the system Eq.\,\eqref{consensus} is asymptotically unstable.


Next, we comment on the relevance of both the spectral abscissa and the reactivity on the transverse consensus dynamics \eqref{transverse}. The asymptotic rate of $\bZ(t)$ in \eqref{transverse} is given by $\alpha(-A)$, i.e., the spectral abscissa of $-A$.
In general, $\alpha(-A)$ can be either positive or negative. 
On the other hand, the instantaneous rate of growth of the norm of the vector $\bZ$ is given by the reactivity,
\begin{align} \label{eq:reactivity0}
\begin{split}
    \xi\left( -A \right) & =  \max_{\|  \bZ\| \neq 0} \left[ \frac{1}{\|  \bZ \|} \frac{d \|  \bZ\|}{d t} \right] = \lambda_{\max} \left( -\frac{A+A^\top}{2} \right).
\end{split}
\end{align}
Thus a positive (negative) reactivity $\xi(-A)$ indicates that the norm of the state $\bZ$ tends to grow (shrink) in the limit of $t \rightarrow 0$. 

\begin{proposition}
    The algebraic connectivity of the graph represented by the Laplacian matrix in Eq.\,\eqref{consensus} is equal to the negative of the reactivity of the transverse dynamics in Eq.\,\eqref{transverse}, that is, $f(L) \equiv - \xi(-A)$.
\end{proposition}
\begin{proof}
    The reactivity of the dynamics in Eq.\,\eqref{eq:reactivity0} is
    \begin{align*}
        \xi\left( -A \right) & :=  \max_{\|  \bZ\| \neq 0} \left[ \frac{1}{\|  \bZ \|} \frac{d \|  \bZ\|}{d t} \right] \\
        & = \max_{\|  \bZ\| \neq 0} \left[ \frac{1}{\sqrt{\bZ^\top \bZ}} \frac{d \sqrt{\bZ^\top \bZ}}{d t} \right] \\
        & = \max_{\|  \bZ\| \neq 0} \left[ \frac{1}{\sqrt{\bZ^\top \bZ}} \frac{\dot{\bZ}^\top \bZ + \bZ^\top \dot{\bZ}}{2\sqrt{\bZ^\top \bZ}} \right] \\
        & = \max_{\|  \bZ\| \neq 0} \left[ \frac{\bZ^\top (-A - A^\top) \bZ}{2 \bZ^\top \bZ} \right].
    \end{align*}
    Also, $V^\top V = I$, $L = V A V^\top$ and $\bX = V \bZ$ which indicates $\bX \perp \bone$ since $\bX \in \text{range}(V)$, so
    {\footnotesize
    \begin{align*}
        \xi\left( -A \right) & = \max_{\|  \bZ\| \neq 0} \left[ \frac{\bZ^\top  \left(-(V^\top V) A (V^\top V) - (V^\top V) A^\top (V^\top V)\right) \bZ}{2 \bZ^\top \bZ} \right] \\
        & = \max_{\|  \bZ\| \neq 0} \left[ \frac{(\bZ^\top V^\top)  (- L - L^\top) (V\bZ)}{2 (\bZ^\top V^\top) (V\bZ)} \right] \\
        & = \max_{ \bX \neq 0, \bX \perp \bone} \left[ \frac{\bX^\top  (- L - L^\top) \bX}{2 \bX^\top \bX} \right] \\
        & = -\min_{ \bX \neq 0, \bX \perp \bone} \left[ \frac{\bX^\top  ( L + L^\top) \bX}{2 \bX^\top \bX} \right] \\
        & = -\min_{ \bX \neq 0, \bX \perp \bone} \left[ \frac{\bX^\top  L \bX}{\bX^\top \bX} \right] \\
        & = -f (L).
    \end{align*}}
    That concludes the proof.
\end{proof}





In what follows, we refer to $\xi(-A)$ in Eq.\ \eqref{eq:reactivity0} as either the reactivity of the dynamics in Eq.\,\eqref{transverse} or the transverse reactivity of the consensus dynamics in Eq.\,\eqref{consensus}, interchangeably.



\begin{corollary}
    The instantaneous rate of growth of the transverse consensus dynamics in Eq.\,\eqref{transverse} is greater or equal to the asymptotic rate of growth, i.e., $\xi (-A) \geq \alpha (-A)$.
\end{corollary}
\begin{proof}
    The relation $\xi (-A) \geq \alpha (-A)$ is obtained by placing $M = -A$ in Theorem \ref{theo}.
\end{proof}


To illustrate the relation $\xi (-A) \geq \alpha (-A)$, we provide a numerical example of stable and unstable dynamics.
First, we randomly generate the proper Laplacian matrix
\begin{equation} \label{Lproper}
    L = \begin{bmatrix}
    1 & -1 & 0 & 0 & 0 \\ 
    0 & 1 & 0 & 0 & -1 \\ 
    -1 & -1 & 2 & 0 & 0 \\ 
    0 & 0 & 0 & 0 & 0 \\ 
    0 & 0 & 0 & -1 & 1 
    \end{bmatrix}
\end{equation}
with all eigenvalues with non-negative real-part and only one zero eigenvalue.
The matrix $A$ in Eq.\,\eqref{transverse} corresponding to the Laplacian matrix $L$ in \eqref{Lproper} is
\begin{equation}
    A = \begin{bmatrix}
    1.5427 & -0.1809 & 0.2337 & -0.9045 \\ 
    -0.1483 & 2.1281 & 0.5427 & 0.4045 \\ 
    0.5427 & -0.1809 & 0.2337 & 0.0955 \\ 
    0.5427 & -0.1809 & -0.7663 & 1.0955 
    \end{bmatrix}.
\end{equation}
Here, $\xi(-A) = -0.0250 \geq \alpha(-A) = -1$. 
We simulate Eq.\,\,\eqref{transverse} from the initial condition $\bZ(0) = [-0.2187,\ -0.1115,\ 0.9320 ,\ 0.2667 ]$ and plot the norm of the transverse perturbation $\| \bZ(t) \|$ in Fig.\,\ref{fig:consensus_rate} A the in log scale.
As evident by the slope of $\| \bZ(t) \|$ vs $t$, the transient rate of decay $\xi(-A)$ is larger than the asymptotic rate of decay $\alpha(-A)$.

Next, we add an adversarial agent to the above Laplacian matrix with a budget $-c = -0.001$ and obtain
\begin{equation} \label{Lnonproper}
    L_{aug} = \begin{bmatrix}
    1 & -1 & 0 & 0 & 0 & 0 \\ 
    0 & 1 & 0 & 0 & -1 & 0 \\ 
    -1 & -1 & 2 & 0 & 0 & 0 \\ 
    0 & 0 & 0 & -0.001 & 0 & 0.001 \\ 
    0 & 0 & 0 & -1 & 1 & 0 \\ 
    0 & 0 & 0 & 0.001 & 0 & -0.001 
    \end{bmatrix}
\end{equation}
which has mixed signed real-part eigenvalues.
For the Laplacian matrix \eqref{Lnonproper}, the corresponding matrix $A_{aug}$ in Eq.\,\eqref{transverse} is
\begin{equation}
    A_{aug} = \begin{bmatrix}
    1.4923 & -0.1527 & 0.2024 & -0.9160 & 0.0840 \\ 
    -0.2178 & 2.1372 & 0.4923 & 0.3739 & 0.3739 \\ 
    0.4923 & -0.1527 & 0.2014 & 0.0840 & 0.0850 \\ 
    0.4923 & -0.1527 & -0.7976 & 1.0840 & 0.0840 \\ 
    0.4923 & -0.1527 & 0.2034 & 0.0840 & 0.0830 
    \end{bmatrix}.
\end{equation}
Here, $\xi(-A_{aug}) = 0.0874 \geq \alpha (-A_{aug}) = 0.002$.
We simulate Eq.\,\,\eqref{transverse} from the initial condition $\bZ(0) = [0.1202,\  0.1937,\ 0.1545,\ 0.8538,\ 0.4418]$ and plot the norm of the transverse perturbation $\| \bZ(t) \|$ in Fig.\,\ref{fig:consensus_rate} B in the log scale.
As evident by the slope of $\| \bZ(t) \|$ vs $t$, the transient rate of decay $\xi(-A)$ is larger than the asymptotic rate of decay $\alpha(-A)$.




\begin{figure}
    \centering
    \includegraphics[width=0.6\linewidth]{figs/Consensus_rate.pdf}
    \caption{\textbf{Transient growth and decay in consensus dynamics.}
Panels A and B show the transverse consensus dynamics governed by Eq.~\eqref{transverse} for the cases of stable and unstable dynamics, respectively. In both panels, the vertical axis is shown on a logarithmic scale, so that the slopes of the curves represent the rates of decay (panel A) and growth (panel B). In both cases, the transient rate, $\xi(-A)$, exceeds the corresponding asymptotic rate, $\alpha(-A)$, demonstrating that the initial decay or growth is faster than its long-term asymptotic behavior.}
%   Transverse consensus dynamics of A: stable dynamics, and B: unstable dynamics from Eq.\,\eqref{transverse}. The $y$ axis in both plots is in the log scale, so the slopes of the curves show the rate of growth/decay. In both panels, the transient rate of growth/decay $\xi (-A)$ is greater than the asymptotic rate of growth/decay $\alpha(-A)$. } 
    \label{fig:consensus_rate}
\end{figure}



\subsection{Proof of Proposition \ref{prop1}} \label{sec:prop1}

Given the possibly asymmetric Laplacian $L$ and the budget $-c < 0$, we solve
    % \begin{subequations} 
        \begin{align}
        \begin{split}
            \min_{\bb, \bx } \quad 
            & \bx^\top \left(V^\top \dfrac{{L_{aug}^b}+{L_{aug}^b}^\top}{2} V \right)\bx \\
            % & x_1^2 c -2x_1 \bx_r^\top \bb + \bx_r^\top L \bx_r +\bx_r^\top \diag(\bb) \bx_r \\
            \text{subject to} \quad 
            % &  \bx = \begin{bmatrix}
            %     x_1 \\ \bx_r
            % \end{bmatrix}, \\
            & \bx^\top \bx = 1 , \\
            & L_{aug}^b = \begin{bmatrix}
                L + \text{diag}(\bb) & -\bb \\
                -\bb^\top & -c
            \end{bmatrix}, \\
            & b_i \leq 0, \quad \forall i = 1, \hdots, N,\\
            & \sum_{i=1}^N b_i = -c.
        \end{split}
        \end{align}
    % \end{subequations}
    First, we set $\by:= V\bx$ which also implies $\bone^\top \by = 0$ and rewrite the optimization in $\by$ and $\bb$:
        % \begin{subequations} 
        \begin{align}
        \begin{split}
            \min_{\bb, \by } \quad 
            & \by^\top \dfrac{{L_{aug}^b}+{L_{aug}^b}^\top}{2} \by \\
            % & x_1^2 c -2x_1 \bx_r^\top \bb + \bx_r^\top L \bx_r +\bx_r^\top \diag(\bb) \bx_r \\
            \text{subject to} \quad 
            % &  \bx = \begin{bmatrix}
            %     x_1 \\ \bx_r
            % \end{bmatrix}, \\
            & \by^\top \by = 1 , \\
            & \bone^\top \by = 0, \\
            & {L_{aug}^b} = \begin{bmatrix}
                L + \text{diag}(\bb) & -\bb \\
                -\bb^\top & -c
            \end{bmatrix}, \\
            & b_i \leq 0, \quad \forall i = 1, \hdots, N,\\
            & \sum_{i=1}^N b_i = -c.
        \end{split}
        \end{align}
    % \end{subequations}
In the above optimization, we set $\by = : [\by_r^\top, y_{N+1}]^\top$ and $\by_r^\top = [y_1, \ y_2, \hdots, y_N]^\top \in \mathbb{R}^N$, i.e., we separate the last entry of the vector $\by$ from the remaining entries.
We also assume $y_i \neq y_j, \forall i\neq j, i= 2, \hdots, N, j = 2, \hdots, N.$
The objective function is
\begin{align}
\begin{split}
    \by^\top \dfrac{{L_{aug}^b}+{L_{aug}^b}^\top}{2} \by = & -y_{N+1}^2 c+ \by_r^\top \frac{L+L^\top}{2} \by_r \\ 
    &  -2y_{N+1} \by_r^\top \bb +\by_r^\top \diag(\bb) \by_r.
\end{split}
\end{align}
The two terms $-2y_{N+1} \by_r^\top \bb +\by_r^\top \diag(\bb) \by_r$ in the above equation that explicitly depend on $\bb$ are written as
\begin{align}
\begin{split}
    -2y_{N+1} \by_r^\top \bb + \by_r^\top \diag(\bb) \by_r & = \sum_{i=1}^N -2 y_{N+1} y_i b_i  + \sum_{i=1}^N y_i^2 b_i \\
    & = \sum_{i=1}^N \left( y_i^2 -2 y_{N+1} y_i  \right) b_i \\
    & =: \sum_{i=1}^N \beta_i b_i
\end{split}
\end{align}
where $\beta_i =  y_i^2 -2 y_{N+1} y_i, \ i = 1, 2, \hdots, N$.
Since these two terms in the objective function depend linearly on $\bb$ and $\sum_j b_j = -c$ and $b_j \leq 0, \forall j$, thus one focuses all budget on the node with the maximum $\beta_i$, i.e., $b_{i^*} = -c$, $b_j = 0, \ \forall j \neq i^*$, and $i^* = \arg \max_i \beta_i$. 
That concludes the proof. \hfill \qed

{We remark that Proposition 1 simply predicts under generic assumptions the particular structure of $b^*$, the optimal $b$, namely that all the entries of $b^*$ will be zero except for one, but it does not predict which one.}

\subsection{Proof of Proposition \ref{prop2}} \label{sec:prop2}

We evaluate the algebraic connectivity using \eqref{eq:algeb} and take the vector $\bX_0 = [N, \ -1, \allowbreak \ -1, \allowbreak \cdots, \allowbreak  -1]^\top$ such that its entries sum to 0.
Evaluating the Rayleigh quotient for $L_{aug}$ from \eqref{eq:Lnewbi} and $\bX_0$ yields
\begin{equation*}
    f \leq \dfrac{\bX_0^\top L_{aug} \bX_0}{\bX_0^\top \bX_0} = -\dfrac{(N+1)^2 c}{N^2 + N} \leq 0.
\end{equation*}
Thus $f \leq 0$. \hfill \qed



\subsection{Proof of Proposition \ref{prop3}} \label{sec:prop3}

Given the possibly asymmetric Laplacian $L$ and the budget $-c < 0$, we solve
    % \begin{subequations} 
        \begin{align}
        \begin{split}
            \min_{\bb, \bx } \quad 
            & \bx^\top \left(V^\top \dfrac{{L_{aug}^u}+{L_{aug}^u}^\top}{2} V \right)\bx \\
            % & x_1^2 c -2x_1 \bx_r^\top \bb + \bx_r^\top L \bx_r +\bx_r^\top \diag(\bb) \bx_r \\
            \text{subject to} \quad 
            % &  \bx = \begin{bmatrix}
            %     x_1 \\ \bx_r
            % \end{bmatrix}, \\
            & \bx^\top \bx = 1 , \\
            & L_{aug}^u = \begin{bmatrix}
                L + \text{diag}(\bb) & -\bb \\
                \pmb{0}^\top & 0
            \end{bmatrix}, \\
            & b_i \leq 0, \quad \forall i = 1, \hdots, N,\\
            & \sum_{i=1}^N b_i = -c.
        \end{split}
        \end{align}
    % \end{subequations}
    First, we set $\by:= V\bx$ which also implies $\bone^\top \by = 0$ and rewrite the optimization in $\by$ and $\bb$:
        % \begin{subequations} 
        \begin{align}
        \begin{split}
            \min_{\bb, \by } \quad 
            & \by^\top \dfrac{{L_{aug}^u}+{L_{aug}^u}^\top}{2} \by \\
            % & x_1^2 c -2x_1 \bx_r^\top \bb + \bx_r^\top L \bx_r +\bx_r^\top \diag(\bb) \bx_r \\
            \text{subject to} \quad 
            % &  \bx = \begin{bmatrix}
            %     x_1 \\ \bx_r
            % \end{bmatrix}, \\
            & \by^\top \by = 1 , \\
            & \bone^\top \by = 0, \\
            & L_{aug}^u = \begin{bmatrix}
                L + \text{diag}(\bb) & -\bb \\
                \pmb{0}^\top & 0
            \end{bmatrix}, \\
            & b_i \leq 0, \quad \forall i = 1, \hdots, N,\\
            & \sum_{i=1}^N b_i = -c.
        \end{split}
        \end{align}
    % \end{subequations}
In the above optimization, we set $\by = : [\by_r^\top, y_{N+1}]^\top$ and $\by_r^\top = [y_1, \ y_2, \hdots, y_N]^\top \in \mathbb{R}^N$, i.e., we separate the last entry of the vector $\by$ from the remaining entries.
We also assume $y_i \neq y_j, \forall i\neq j, i= 2, \hdots, N, j = 2, \hdots, N.$
The objective function is
\begin{align}
\begin{split}
    \by^\top \dfrac{{L_{aug}^u}+{L_{aug}^u}^\top}{2} \by = & -y_{N+1} \by_r^\top \bb + \by_r^\top \frac{L+L^\top}{2} \by_r \\ 
    & +\by_r^\top \diag(\bb) \by_r
\end{split}
\end{align}
The two terms $-y_{N+1} \by_r^\top \bb +\by_r^\top \diag(\bb) \by_r$ in the above equation that explicitly depend on $\bb$ are written as
\begin{align}
\begin{split}
    -y_{N+1} \by_r^\top \bb + \by_r^\top \diag(\bb) \by_r & = \sum_{i=1}^N - y_{N+1} y_i b_i  + \sum_{i=1}^N y_i^2 b_i \\
    & = \sum_{i=1}^N \left( y_i^2 - y_{N+1} y_i  \right) b_i \\
    & =: \sum_{i=1}^N \beta_i b_i
\end{split}
\end{align}
where $\beta_i =  y_i^2 - y_{N+1} y_i, \ i = 1, 2, \hdots, N$.
Since these two terms in the objective function depend linearly on $\bb$ and $\sum_j b_j = -c$ and $b_j \leq 0, \forall j$, thus one focuses all budget on the node with the maximum $\beta_i$, i.e., $b_{i^*} = - c$, $b_j = 0, \ \forall j \neq i^*$, and $i^* = \arg \max_i \beta_i$. 
That concludes the proof. \hfill \qed


\subsection{Proof of Proposition \ref{prop4}} \label{sec:prop4}

    We evaluate the algebraic connectivity using \eqref{eq:algeb} and take the vector $\bX_0 = [N, \ -1, \allowbreak \ -1, \allowbreak \cdots, \allowbreak  -1]^\top$ such that its entries sum to 0.
    Evaluating the Rayleigh quotient for $L_{aug}$ from \eqref{eq:Lnewuni} and $\bx_0$ yields
    \begin{equation*}
        f \leq \dfrac{\bX_0^\top L_{aug} \bX_0}{\bX_0^\top \bX_0} = -\dfrac{(N+1) c}{N^2 + N} \leq 0.
    \end{equation*}
    Thus $f \leq 0$, and the proof is complete. \hfill \qed



\subsection{Digraphs with a Unidirectional Connection} \label{sec:uni}

Here we use matrix perturbation theory \cite{bamieh2022tutorial} to predict the changes in the eigenvalues of a given matrix $A_0$ subject to a small perturbation $\epsilon A_1$, i.e., $\Lambda_\epsilon = \Lambda_0 + \epsilon \Lambda_1 + \epsilon^2 \Lambda_2 + \cdots$ where $\Lambda_\epsilon$ is the estimated eigenvalues of the matrix $A_\epsilon = A_0 + \epsilon A_1$, $\Lambda_0$ is the exact eigenvalues of $A_0$, and $\epsilon^k \Lambda_k, \ \forall k \geq 1$, are the $k$th order approximation of changes in the eigenvalues $\Lambda_0$.


In the following, we focus on the case of directed networks, i.e., for which the matrix $A_0$ is asymmetric and set $k = 1$.
Thus, $\Lambda_1 = \diag(W_0^* A_1 V_0)$, where $W_0$ and $V_0$ are the matrices of the left and the right eigenvectors of $A_0$ such that $W_0^* V_0 = I$, and the superscript ${}^*$ denotes conjugate transpose.


\textbf{The limit of small $c$ and unidirectional connection:} In this case, we set $\epsilon = c$ and
\begin{equation}
    A_0 =  V^\top \dfrac{\tilde{L} + \tilde{L}^\top}{2} V, \quad A_1 = V^\top \dfrac{P_i + P_i ^\top}{2} V
\end{equation}
 where 
 % \begin{subequations}
\begin{align} \label{eq:lagematrix2}
    \tilde{L} = 
    % \begin{bmatrix}
    %     L & 0 \\
    %     0 & 0
    % \end{bmatrix}, 
    \begin{bmatrix}
    & & & & & & & 0\\ 
    & & & & & & & \vdots \\ 
    & & & & & & & 0 \\ 
    & & & L & & & & 0 \\ 
    & & & & & & & 0 \\ 
    & & & & & & & \vdots \\ 
    & & & & & & & 0 \\ 
    0 & \hdots & 0 & 0 & 0 & \hdots & 0 & 0
    \end{bmatrix}, \\
    P_i = 
    % \begin{bmatrix}
    %     & & & \\
    %     & -1 & & 1 \\
    %     & & & \\
    %     & 1 & & -1
    % \end{bmatrix}
    \begin{bmatrix}
    0 & \hdots & 0 & 0 & 0 & \hdots & 0 & 0 \\ 
    \vdots & \ddots & \vdots & \vdots & \vdots & \ddots & \vdots & \vdots \\ 
    0 & \hdots & 0 & 0 & 0 & \hdots & 0 & 0 \\ 
    0 & \hdots & 0 & -1 & 0 & \hdots & 0 & 1 \\ 
    0 & \hdots & 0 & 0 & 0 & \hdots & 0 & 0 \\ 
    \vdots & \ddots & \vdots & \vdots & \vdots & \ddots & \vdots & \vdots \\ 
    0 & \hdots & 0 & 0 & 0 & \hdots & 0 & 0 \\ 
    0 & \hdots & 0 & 0 & 0 & \hdots & 0 & 0 
    \end{bmatrix}
\end{align}
 % \end{subequations}
where $P_i$ has the entry -1 on the main diagonal on the $i$th position.
The first order matrix perturbation theory is then
\begin{align}
\begin{split}
    f_b & = f_0 + c \dfrac{df}{db} \\
    & = \bv_0^\top A_0 \bv_0 + c \bv_0^\top A_1 \bv_0 \\
    & = \bv_0^\top V^\top \dfrac{\tilde{L} + \tilde{L}^\top}{2} V \bv_0 + c \bv_0^\top V^\top \dfrac{P_i + P_i ^\top}{2} V \bv_0 \\
    & =: \by_0^\top \dfrac{\tilde{L} + \tilde{L}^\top}{2} \by_0 + c \by_0^\top \dfrac{P_i + P_i ^\top}{2} \by_0 \\
    & =: f_0 + c \by_0^\top \dfrac{P_i + P_i ^\top}{2} \by_0
\end{split}
\end{align}
where $\by_0 := V \bv_0$ and $\by_0^\top\left(\frac{\tilde{L} + \tilde{L}^\top}{2} \right) \by_0 = f_0 $. 
Note that since the graph is balanced, the Laplacian matrix $L$ has both row and column sums equal to zero. 
Also, considering the block diagonal structure of the matrix $\tilde{L}$ results in having two zero eigenvalues, one corresponding to the eigenvector with all entries equal ones and the other eigenvalue $f_0 = 0$ corresponding to the eigenvector $\by_0 = \frac{1}{\sqrt{N(N+1)}}[1, \ 1,\ \hdots 1, \ -N] \in \mathbb{R}^{N+1}$.

We are interested in the change in the zero eigenvalue corresponding to the eigenvector $\by_0$.
It thus follows
\begin{align}
\begin{split} \label{eq:smallb2}
    f_b & = f_0 + c \by_0^\top \dfrac{P_i + P_i ^\top}{2} \by_0 = 0 + c (-0.1) = -0.1 c.
\end{split}
\end{align}

\textbf{The limit of large $c$ and unidirectional connection:}
Here, we set $\epsilon = 1$ and
\begin{equation}
    A_0 = c V^\top \frac{P_i+P_i^\top}{2}V = c V^\top P_i V, \quad A_1 = V^\top \dfrac{\tilde{L} + \tilde{L}^\top}{2} V
\end{equation}
where the matrices $P_i$ and $\tilde{L}$ are defined in Eq.\, \eqref{eq:lagematrix2}.
Now, the eigenvalue of interest is the most negative nonzero eigenvalue of $\frac{P_i+P_i^\top}{2}$ which is $f_0 =  -\frac{\sqrt{2}+1}{2}$.
The corresponding eigenvector is $\by_0^i = [0 \ \cdots \ 0 \ \ {\sqrt{2}+1} \ \ 0 \ \cdots \ 0 \ \ -1]^\top$ where the $i$th and the last entry of this eigenvectors have nonzero elements.
We set $\by_0^i = V\bv_0^i$ as the effective vector for the first-order approximation of matrix perturbation theory.
The first order change in the eigenvalue $f_0$ is ${\by_0^i}^\top A_1 \by_0^i = {\by_0^i}^\top \frac{\tilde{L} + \tilde{L}^\top}{2} \by_0^i = (3+2\sqrt{2})  L_{ii} > 0$.
Therefore, the perturbed eigenvalue after pinning node $i$ is
\begin{equation}
    f_b = f_0 + \epsilon {\by_0^i}^\top A_1 \by_0^i = -\frac{\sqrt{2}+1}{2}c + (3+2\sqrt{2})  L_{ii}.
\end{equation}






\subsection{Digraphs with a Bidirectional Connection} \label{sec:bi}


\textbf{The limit of small $c$ and bidirectional connection:} all the derivations are similar to the case of unidirectional connection except for the matrix
\begin{equation} \label{eq:pbi}
    P_i = 
    \begin{bmatrix}
    0 & \hdots & 0 & 0 & 0 & \hdots & 0 & 0 \\ 
    \vdots & \ddots & \vdots & \vdots & \vdots & \ddots & \vdots & \vdots \\ 
    0 & \hdots & 0 & 0 & 0 & \hdots & 0 & 0 \\ 
    0 & \hdots & 0 & -1 & 0 & \hdots & 0 & 1 \\ 
    0 & \hdots & 0 & 0 & 0 & \hdots & 0 & 0 \\ 
    \vdots & \ddots & \vdots & \vdots & \vdots & \ddots & \vdots & \vdots \\ 
    0 & \hdots & 0 & 0 & 0 & \hdots & 0 & 0 \\ 
    0 & \hdots & 0 & 1 & 0 & \hdots & 0 & -1
    \end{bmatrix}.
\end{equation}
This results in
\begin{align}
\begin{split} 
    f_b & = f_0 + c \by_0^\top \dfrac{P_i + P_i ^\top}{2} \by_0 = 0 + c (-1.1) = -1.1 c.
\end{split}
\end{align}




\textbf{The limit of large $c$ and bidirectional connection:}
Here, we set $\epsilon = 1$ and
\begin{equation}
    A_0 = c V^\top \frac{P_i+P_i^\top}{2}V = c V^\top P_i V, \quad A_1 = V^\top \dfrac{\tilde{L} + \tilde{L}^\top}{2} V
\end{equation}
where the matrices $P_i$ and $\tilde{L}$ are defined in Eqs.\,\eqref{eq:pbi} and \eqref{eq:lagematrix2}.
Now, the eigenvalue of interest is the one and only nonzero eigenvalue of $V^\top P_i V$ which is $f_0 = -2$.
Since $V$ is a similarity transformation for $N$ number of the eigenvalues, the matrix $P_i$ has the same eigenvalue $f_0 = -2$ and the corresponding eigenvector is $\by_0^i = [0 \ \cdots \ 0 \ \ {-1/\sqrt{2}} \ \ 0 \ \cdots \ 0 \ \ 1/\sqrt{2} ]^\top$ where the $i$th and the last entry of this eigenvectors have nonzero elements.
We define $\by_0^i := V\bv_0^i$ as the effective vector for the first-order approximation of matrix perturbation theory.
The first order change in the eigenvalue $f_0$ is ${\by_0^i}^\top A_1 \by_0^i = {\by_0^i}^\top \frac{\tilde{L} + \tilde{L}^\top}{2} \by_0^i = L_{ii} /2 > 0$.
Therefore, the perturbed eigenvalue after pinning node $i$ is
\begin{equation}
    f_b = f_0 + \epsilon {\by_0^i}^\top A_1 \by_0^i = -2c + \dfrac{L_{ii}}{2}.
\end{equation}


\subsection{Proof of Mean-Slope Relation} \label{sec:meanslope}

We show this relation $< df / d c > = -1/N$ in what follows.
{
We emphasize that we proceed without introducing the assumption that the network is  balanced.
}
Based on first order perturbation theory, $df_i / dc = \by_0^\top \frac{P_i + P_i^\top}{2} \by_0$ where $\by_0 = V \bv_0$ where $\bv_0$ is eigenvector corresponding to the smallest eigenvalue of the matrix $V^\top \frac{\tilde{L} + \tilde{L}^\top}{2}V$ and the matrices $P_i$ and $\tilde{L}$ are defined in Eq.\,\eqref{eq:lagematrix2}.
Note that $\by_0$ is in the range of the matrix $V$, resulting in the zero-sum of the entries of the vector $\by_0$.
It follows
\begin{align}
\begin{split}
    \left< \frac{d f}{d c} \right> = \dfrac{1}{N} \sum_{i=1}^N \frac{d f_i}{d c} & = \dfrac{1}{N} \by_0^\top \left( \sum_{i=1}^N  \dfrac{P_i + P_i^\top}{2}\right) \by_0 \\
    & =: \dfrac{1}{N} \by_0^\top \left( \bar{P}\right) \by_0,
\end{split}
\end{align}
where the matrix $\bar{P}$ has the following form
\begin{equation}
    \bar{P} = \sum_{i=1}^N  \dfrac{P_i + P_i^\top}{2} = \begin{bmatrix}
-1 &  &  &  & \frac{1}{2} \\ 
 & -1 & &  & \frac{1}{2} \\ 
 &  & \ddots &  & \vdots \\ 
 &  &  & -1 & \frac{1}{2} \\ 
\frac{1}{2} & \frac{1}{2} & \hdots & \frac{1}{2} & 0 
\end{bmatrix}.
\end{equation}
It follows that the matrix $\bar{P}$ has $N-1$ eigenvalues which are equal to $-1$ and the corresponding eigenvectors have their sum of the entries equal to zero. 
This suggests that these eigenvectors form a basis for the vector $\by_0$ and thus, the product $\by_0^\top \bar{P} \by_0 = -1$. 
Therefore, $< df / d c > = -1/N$. 
\hfill \qed




\subsection{Asymptotic Stability Analysis for the Case of Synchronization Dynamics \label{Sec:MAI}}

Next we present standard derivations from \cite{fujisaka1983stability,pecora1998master}.
A particular solution for the synchronization dynamics  \eqref{eq:synch} is the synchronous solution $\bx_1(t)=\bx_2(t)=...=\bx_N(t)=\bs(t)$, where $\bs(t)$, evolves in time based on the equation:
\begin{equation} \label{eq:s}
    \dot{\bs} (t) = \bF(\bs(t)).
\end{equation}
The synchronization manifold is the sets of phase space points that satisfy  $\bx_1 = \bx_2 = \hdots \bx_N$.

We study the time evolution of infinitesimal perturbations about the synchronization manifold, $\delta \bx_i (t) \coloneq \bx_i (t) - \bs(t)$. 
The synchronization is stable if $\delta \bx_i (t) \to 0$ as $t \to \infty$, $\forall i$, which implies $\bx_i (t) \to \bs(t)$ as $t \to \infty$, $\forall i$. 

To study the asymptotic stability of the perturbations $\delta \bx_i (t)$, we linearize Eq.\,\eqref{eq:synch} about $\bs(t)$ and derive the following linearized sets of dynamical equations:
\begin{equation} \label{eq:dx}
    \delta \dot{\bx}_i (t) = D\bF(\bs(t)) \delta \bx_i(t) - \sigma \sum_{j = 1}^{N} L_{ij} D\bH (\bs(t)) \delta \bx_j (t),
\end{equation}
$i = 1, \hdots, N$, where $D\bF (\bs(t))$ and $D \bH(\bs(t))$ are the Jacobians of $\bF$ and $\bH$ evaluated at $\bs(t)$, respectively.
A compact form of the above equations can be obtained by first concatenating the perturbation vectors as $\delta \bX(t) \coloneq \left[\delta \bx_1 (t)^\top, \hdots, \delta \bx_N (t)^\top \right]^\top \in \mathbb{R}^{mN}$, and then writing:
\begin{equation} \label{eq:compact}
    \delta \dot{\bX} (t) = \Bigg(I_N \otimes D \bF(s) - \sigma L \otimes D\bH(s)\Bigg) \delta \bX(t),  
\end{equation}
where $\otimes$ denotes the Kronecker product. We call $V$ the matrix having for columns the eigenvectors of $L$ and $\Lambda$ the associated diagonal matrix, having the eigenvalues of $L$ on the main diagonal.
Equation \eqref{eq:compact} can be decoupled through the eigen-decomposition of the Laplacian matrix $L$ by applying the similarity transformation $T = V \otimes I_m \in \mathbb{R}^{mN \times mN}$ as
\begin{equation} \label{eq:compact1}
    \delta \dot{\bZ} (t) = \Bigg(I_N \otimes D \bF(s) - \sigma \Lambda \otimes D\bH(s)\Bigg) \delta \bZ(t),  
\end{equation}
where $\delta \bZ (t) = T^{-1} \delta \bX (t) \in \mathbb{R}^{mN}$. 
By setting $[\delta \bz_1(t)^{\top}, \hdots , \delta \bz_N (t)^\top]^\top \coloneq \delta \bZ (t)$, with $\delta \bz_i (t) \in \mathbb{R}^m, \forall i$, we can write the decomposed equation of transformed perturbations as
\begin{equation} \label{dzi}
    \delta \dot{\bz}_i (t) = \Bigg(D\bF(\bs(t)) - \sigma \lambda_i D\bH(\bs (t)) \Bigg) \delta \bz_i (t),
\end{equation}
$i = 1, \hdots, N$. By construction, one eigenvalue $\lambda_1=0$, which is associated to stability of the dynamics on the synchronization manifold. As such it is excluded from the analysis that follows, that focus on the stability of the dynamics transverse to this manifold.
The synchronization is stable if $\delta \dot{\bz}_i (t) \to 0$ as $t \to 0$, $\forall i$.

From Eq.\ \eqref{dzi} we see that the stability of the synchronization is dependent on the Jacobians $D\bF(\bs(t))$ and $D\bH(\bs(t))$, and the scalar $\sigma \lambda_i$.
We introduce the parameter $K = \sigma \lambda_i$ and determine the stability of
\begin{equation} \label{eq:decomposed}
    \dot{\by} (t) = \Bigg(D\bF(\bs(t)) - K D\bH(\bs (t)) \Bigg) \by (t),
\end{equation}
where $\by (t) \in \mathbb{R}^m$ is a generic vector representing perturbations.
The stability is assessed through the evaluation of the Maximum Lyapunov Exponent (MLE),
\begin{equation}
    \text{MLE} = \lim_{t\to \infty} \dfrac{1}{t} \log \dfrac{\| \by(t) \|}{\| \by (0) \|}.
\end{equation}
If the MLE is negative (positive), the above system is stable (unstable).
Now, assume the set of parameters $\mathcal{K}$ for which the system \eqref{eq:decomposed} is stable (MLE $<0$).
If $\sigma \lambda_i \in \mathcal{K}$, $i=2,...,N$, then the synchronization dynamics \eqref{eq:synch} is stable and $\bx_i (t) \to \bs(t)$ when $t \to \infty$.
The master stability function (MSF) maps the parameter $K$ to the maximum Lyapunov exponent (MLE) of Eq.~\eqref{eq:decomposed}.

In the following, we present MSF plots for real values of $K$ for three well-known oscillators, the Lorenz, R\"ossler, and Chua systems, and for several choices of the output function $\mathbf{H}(\mathbf{x})$. For all combinations of $\mathbf{F}(\mathbf{x})$ and $\mathbf{H}(\mathbf{x})$ considered in this study, the MSF is positive for negative values of $K$. Under the assumption that this property holds in general, it follows that  a sufficient condition for the synchronous state to become unstable is that at least one of the eigenvalues of the Laplacian matrix $L$ has negative real part.

\textbf{Lorenz MSF plots}

The state vector, the local dynamics, and the coupling function for the Lorenz system are 
\begin{align}
    \bx = \begin{bmatrix}
        x \\ y \\ z
    \end{bmatrix}, \quad \bF (\bx) = \begin{bmatrix}
        10( y-x ) \\
        x( 28-z ) - y \\
        xy - \frac{8}{3} z
    \end{bmatrix}, \quad \bH(\bx) = H \bx,
\end{align}
where the matrix $H$ is the linear coupling matrix. 
We provide 9 MSF plots for the Lorenz system, where for each plot, a different entry of $H$ is equal to 1, and all others are zero.
It  follows that the Jacobian $D\bH(\bs) = H$.
The Jacobian of the local dynamics is
\begin{equation}
    \bs = \begin{bmatrix}
        x_s \\ y_s \\ z_s
    \end{bmatrix}, \quad D \bF (\bs) = \begin{bmatrix}
        -10 & 10 &  0 \\
        28-z_s& -1& -x_s \\
        y_s & x_s & -\frac{8}{3}
    \end{bmatrix}
\end{equation}

\begin{figure}
    \centering
    \includegraphics[width=0.99\linewidth]{figs/merged_MSF_figure_vertical.pdf}
    \caption{\textbf{Master stability functions for both positive and negative values of the function argument.}
    Master stability plots for the Lorenz system (A), Roessler system (B), and Chua System (C). Plots show the Maximum Lyapunov Exponent (MLE) as the argument $K$ of Eq.\,\eqref{eq:decomposed} is varied.
    In each subfigure, panel in row $i$ (from the top) and column $j$ (from the left) shows the case when $H_{ij} = 1$, and all other entries of the matrix $H$ are zero.
    }
    \label{fig:MSFLorenz}
\end{figure}

Figure \ref{fig:MSFLorenz} shows that the MSF is positive in all different coupling setups for the negative range of the MSF argument $K$. 
Thus, it shows that the synchronization is locally unstable for negative ranges of the parameter $K$.


\textbf{R\"ossler MSF plots}

The state vector, the local dynamics, and the coupling function for the Rossler system are 
\begin{align}
    \bx = \begin{bmatrix}
        x \\ y \\ z
    \end{bmatrix}, \quad \bF (\bx) = \begin{bmatrix}
        -y - z \\
        x + 0.1y \\
        0.1 + ( x-35 )z
    \end{bmatrix}, \quad \bH(\bx) = H \bx,
\end{align}
where the matrix $H$ is the linear coupling matrix. 
We provide 9 MSF plots for the Rossler system, where for each plot, a different entry of $H$ is equal to 1, and all others are zero.
It  follows that the Jacobian $D\bH(\bs) = H$.
The Jacobian of the local dynamics is
\begin{equation}
    \bs = \begin{bmatrix}
        x_s \\ y_s \\ z_s
    \end{bmatrix}, \quad D \bF (\bs) = \begin{bmatrix}
        0 & -1 & -1 \\
        1 & 0.1 & 0 \\
        z_s & 0 & x_s-35
    \end{bmatrix}
\end{equation}

%\begin{figure}
%    \centering
%    \includegraphics[width=0.99\linewidth]{figs/Rossler_MSF_fig.pdf}
%    \caption{
%    The master stability plots of the Rossler system. Plots show the Maximum Lyapunov Exponent (MLE) as the argument $K$ of Eq.\,\eqref{eq:decomposed} is varied.
%    Panel in row $i$ (from the top) and column $j$ (from the left) shows the case when $H_{ij} = 1$, and all other entries of the matrix $H$ are zero.
%    }
%    \label{fig:MSFRossler}
%\end{figure}

Figure \ref{fig:MSFRossler} shows that the MSF is positive in all different coupling setups for the negative range of the MSF argument $K$. 
Thus, it shows that the synchronization is locally unstable for negative ranges of the parameter $K$.


\textbf{Chua's circuit MSF plots}

The state vector, the local dynamics, and the coupling function for the Chua's circuit system are 
\begin{align}
    \bx = \begin{bmatrix}
        x \\ y \\ z
    \end{bmatrix}, \quad \bF (\bx) = \begin{bmatrix}
        10( y - x + g(x) ) \\
        x - y + z \\
        -17.85y]
    \end{bmatrix}, \quad \bH(\bx) = H \bx,
\end{align}
where the matrix $H$ is the linear coupling matrix, and the function 
\begin{equation}
    g(x) = \begin{cases}
        - bx - a + b, & \quad x > 1, \\
        - ax, & \quad  |x| < 1, \\
        - bx + a - b, & \quad x < -1,
    \end{cases}, \quad \begin{cases}
        a = -1.2, \\
        b = -0.7.
    \end{cases}
\end{equation}
We provide 9 MSF plots for the Chua's circuit system, where for each plot, a different entry of $H$ is equal to 1, and all others are zero.
It  follows that the Jacobian $D\bH(\bs) = H$.
The Jacobian of the local dynamics is
\begin{equation}
    \bs = \begin{bmatrix}
        x_s \\ y_s \\ z_s
    \end{bmatrix}, \quad D \bF (\bs) = \begin{bmatrix}
        -10-10 d(x_s) & 10 & 0 \\
    1 & -1 & 1 \\
    0 & -17.85 & 0
    \end{bmatrix}
\end{equation}
where $d(x_s) = b$ when $|x_s| > 1$, and $d(x_s) = a$, otherwise.


%\begin{figure}
%    \centering
%    \includegraphics[width=0.99\linewidth]{figs/Chua circuit_MSF_fig.pdf}
%    \caption{
%    The master stability plots of the Chua's circuit system. Plots show the Maximum Lyapunov Exponent (MLE) as the argument $K$ of Eq.\,\eqref{eq:decomposed} is varied.
%    Panel in row $i$ (from the top) and column $j$ (from the left) shows the case when $H_{ij} = 1$, and all other entries of the matrix $H$ are zero.
%    }
%    \label{fig:MSFChua}
%\end{figure}

Figure \ref{fig:MSFChua} shows that the MSF is positive in all different coupling setups for the negative range of the MSF argument $K$. 
Thus, it shows that the synchronization is locally unstable for negative ranges of the parameter $K$.







\subsection{Asymptotic Stability Analysis for the Case of Power Grid Dynamics \label{Sec:MAII}}

We linearize Eq.\  \eqref{eq:swingsimple} about the synchronous operating state $\theta_i^*= \omega_s t+ \theta_i^0$,
\begin{equation} \label{Eq:J}
\begin{aligned}
     \delta \ddot{\vartheta}_i (t) &= - {\gamma} {\delta \dot\vartheta}_i (t) - \sum_{j= 1, j \neq i}^N A_{ij} \cos(\theta_i^0 - \theta_j^0) (\delta \vartheta_i-\delta \vartheta_j)\\
     &= - {\gamma} {\delta \dot\vartheta}_i (t) - \sum_{j= 1, j \neq i}^N A'_{ij}  (\delta \vartheta_i-\delta \vartheta_j)\\
     &= - {\gamma} {\delta \dot\vartheta}_i (t) - \sum_{j= 1}^N L'_{ij}  \delta \vartheta_j
     \end{aligned}
     \end{equation}
$i=1,...,N$, where $\gamma>0$, $A=[A_{ij}]$ is the original adjacency matrix,  i.e.,  $A_{ij}=A_{ji}=1$ ($0$) if there is (is not) a transmission line between between nodes $i$ and $j$, $A'=[A'_{ij}]$ is the weighted adjacency matrix,
with entries $A'_{ij}= A_{ij} \cos(\theta_i^0 - \theta_j^0)$, and $L'=[L'_{ij}]$ is the associated Laplacian matrix, with entries $L'_{ij}= (\delta_{ij} \sum_j A'_{ij}  - A'_{ij})$. We proceed under the assumption that, prior to the intervention of the intruder, the differences $\theta_i^0 - \theta_j^0$ are small, so that  $A'_{ij}>0$ iff $A_{ij}=1$, for which the Laplacian matrix $L'$ is proper.

By diagonalizing the Laplacian matrix $L'$, the last line of Eq.\ \eqref{Eq:J} can be transformed into the following independent equations,
\begin{equation} \label{Eq:JJ}
    \ddot{y}_k(t)= - {\gamma} {\dot y}_k (t) - \lambda'_{k}  y_k(t),
\end{equation}
$k=1,..,N$, where $\lambda'_{k}$, $k=1,..,N$,  are the eigenvalues of the matrix $L'$. By construction there is one eigenvalue $\lambda'_{k}=0$, which corresponds to marginal stability of the synchronous operating state $\theta_i^*= \omega_s t+ \theta_i^0$.
Transverse stability is determined by the remaining eigenvalues of $L'$, as we explain below. % and is achieved (is not achieved) when these are all positive (when at least one of these is negative.)

In the case that $\lambda'_k$ is real, Eq.\ \eqref{Eq:JJ} is: (i) stable for $\lambda'_{k}>0$, (ii) unstable for $\lambda'_{k}<0$, and (iii) marginally stable for $\lambda'_{k}=0$. 

A more interesting case is for a complex eigenvalue of the Laplacian $\lambda'_k=\alpha+ j \beta$, for which the characteristic equation is,
\begin{equation} \label{pol}
    r^2+\gamma r +(\alpha+j \beta)=0,
\end{equation}
where here $r$ is a complex number and we recall that $\gamma$ is real and positive. The largest real part root of Eq.\ \eqref{pol} is
\begin{equation} \label{largest_real_root}
    -\frac{\gamma}{2} + \frac{1}{2} \sqrt{\frac{|D|+\gamma^2 - 4 \alpha}{2}},
\end{equation}
where the modulus of the complex discriminant $|D|=\sqrt{(\gamma^2-4 \alpha)^2 + 16 \beta^2}$. Eq.\ \eqref{largest_real_root} is negative if 
\begin{equation} \label{crilu}
\beta^2< \gamma^2 \alpha.
\end{equation}
For stability of the synchronous operating state, criterion \eqref{crilu} needs to be satisfied for all the eigenvalues $\lambda'_k=\alpha+ j \beta$ of $L'$ (with the exception of the one zero eigenvalue.) It is easy to see that \eqref{crilu} is never satisfied for $\alpha<0$. 

We thus conclude that a sufficient condition for the synchronous operating state to be unstable is that at least one of the eigenvalues of the  the matrix $L'$ has negative real part (however, instability is also possible when all the eigenvalues of $L'$ have negative or zero real part.)




\subsection{Asymptotic Stability Analysis for the Case of Formation Control Dynamics \label{Sec:MAIII}}

In the case of Eq.\ \eqref{Eq_formation_control}, the dynamics is linear and stability depends uniquely on the spectrum of the matrix $F(L)$. In the case considered in this paper that $F(L) = I_N - \exp (-3L)$, each eigenvalues of the matrix $F(L)$ is equal to $\mu_k=1- \exp (-3 \lambda_k)$, where $\lambda_k$ is an eigenvalue of $L$, $k=1,..,N$. 

It follows that the transformation that maps $\lambda_k$ into  $\mu_k$ preserves the origin and the sign of the real part of the eigenvalues. Specifically,
\begin{equation}
\lambda_k = 0 \iff \mu_k = 0,
\end{equation}
and
\begin{equation}
\operatorname{sgn}\!\left(\Re(\mu_k)\right)
=
\operatorname{sgn}\!\left(\Re(\lambda_k)\right),
\qquad k = 1,\ldots,N,
\end{equation}
where with $\operatorname{sgn}\!(x)$ we indicate the sign of $x$.

We conclude that if the Laplacian matrix $L$ is (is not) proper, the formation control state is stable (unstable.) 


\subsection{Asymptotic Stability Analysis for the Case of the Kuramoto Dynamics \label{Sec:MAIIII}}

We fix the adjacency matrix $A$ and linearize Eq.\  \eqref{eq:swingsimple} about the phase locked state $\theta_i^*= \omega_s t+ \theta_i^0$,
\begin{equation} \label{Eq:K}
\begin{aligned}
     \delta \dot{\vartheta}_i (t) &= -\dfrac{1}{N+1} \sum_{j=1, j \neq i}^{N+1}  A_{ij} \cos(\theta_i^0 - \theta_j^0) (\delta \vartheta_i-\delta \vartheta_j)\\
     &=  -\dfrac{1}{N+1} \sum_{j=1, j \neq i}^{N+1}  A'_{ij}  (\delta \vartheta_i-\delta \vartheta_j)\\
     &= -\dfrac{1}{N+1}  \sum_{j= 1, j \neq i}^{N+1} L'_{ij}  \delta \vartheta_j
     \end{aligned}
     \end{equation}
$i=1,...,N$, where  $A=[A_{ij}]$ is the fixed adjacency matrix defined in Eq. \eqref{eq:Akuramoto} either before or after the attack, $A'=[A'_{ij}]$ is the weighted adjacency matrix,
with entries $A'_{ij}= A_{ij} \cos(\theta_i^0 - \theta_j^0)$, and $L'=[L'_{ij}]$ is the associated Laplacian matrix, with entries $L'_{ij}= (\delta_{ij} \sum_j A'_{ij}  - A'_{ij})$. We proceed under the assumption that, prior to the intervention of the intruder, the differences $\theta_i^0 - \theta_j^0$ are small, so that  $A'_{ij}>0$ iff $A_{ij}=1$, for which the Laplacian matrix $L'$ is proper.

By diagonalizing the Laplacian matrix $L'$, the last line of Eq.\ \eqref{Eq:K} can be transformed as follows,
\begin{equation} \label{Eq:KK}
    \dot{y}_k= - \lambda'_{k}  y_k,
\end{equation}
$k=1,..,N$, where $\lambda'_{k}$, $k=1,..,N$,  are the eigenvalues of the matrix $L'$. Equation \eqref{Eq:KK} is stable if the real part of $\lambda_k'$ is positive, unstable if it is negative, and marginally stable if $\lambda_k'=0$. By construction there is one eigenvalue $\lambda'_{k}=0$, which corresponds to marginal stability of the phase locked state $\theta_i^*= \omega_s t+ \theta_i^0$. Transverse stability is achieved if all the other eigenvalues of $L'$ have positive real part. 

We conclude that a sufficient condition for the phase-locked state to become unstable is that at least one of the eigenvalues of the  the matrix $L'$ has negative real part.




\section{Reproducibility Code Description}

The MATLAB code used to generate Fig.~\ref{fig:synch_node}B performs the following steps. 
First, the adjacency matrix $A$ is loaded, and the graph Laplacian matrix $L$ is constructed.
For each selected target node $i$, an attack vector $\bb$ is defined with $b_i=-c$ and all other entries equal to zero, with $c=1$ as the budget for the attack. 
The unidirectional augmented Laplacian after the attack is then constructed as in Eq.~\eqref{eq:Lnewuni}.
The simulation uses the original Laplacian before the attack time $t_s=4 s$ and switches to the augmented Laplacian for $t \geq t_s$. 
A network of coupled Lorenz oscillators is then simulated using MATLAB’s ODE solver. 
The norm of the transverse perturbations $\| \delta \bX (t) \|$ is evaluated as in Eq.~\eqref{eq:trans_synch}. 
Finally, the code plots $\| \delta \bX (t) \|$ in time for different attacked nodes, and saves a figure similar to Fig.~\ref{fig:synch_node}B.

\color{black}

\section*{Data availability}
All data generated or analyzed during this study are included in this published article (and its supplementary information files).





\section*{Code availability}


The MATLAB code used in this study is openly available as a reproducible compute capsule on Code Ocean at \url{https://codeocean.com/capsule/0183273/tree/v1} under an MIT license.


\color{black}

\section*{References}
\begin{thebibliography}{10}
\expandafter\ifx\csname url\endcsname\relax
  \def\url#1{\texttt{#1}}\fi
\expandafter\ifx\csname urlprefix\endcsname\relax\def\urlprefix{URL }\fi
\providecommand{\bibinfo}[2]{#2}
\providecommand{\eprint}[2][]{\url{#2}}

\bibitem{6341809}
\bibinfo{author}{Wang, J.} \& \bibinfo{author}{Xin, M.}
\newblock \bibinfo{title}{Integrated optimal formation control of multiple unmanned aerial vehicles}.
\newblock \emph{\bibinfo{journal}{IEEE Transactions on Control Systems Technology}} \textbf{\bibinfo{volume}{21}}, \bibinfo{pages}{1731--1744} (\bibinfo{year}{2013}).

\bibitem{9356608}
\bibinfo{author}{Yang, Y.}, \bibinfo{author}{Xiao, Y.} \& \bibinfo{author}{Li, T.}
\newblock \bibinfo{title}{A survey of autonomous underwater vehicle formation: Performance, formation control, and communication capability}.
\newblock \emph{\bibinfo{journal}{IEEE Communications Surveys \& Tutorials}} \textbf{\bibinfo{volume}{23}}, \bibinfo{pages}{815--841} (\bibinfo{year}{2021}).

\bibitem{1545539}
\bibinfo{author}{Chen, Y.~Q.} \& \bibinfo{author}{Wang, Z.}
\newblock \bibinfo{title}{Formation control: a review and a new consideration}.
\newblock In \emph{\bibinfo{booktitle}{2005 IEEE/RSJ International Conference on Intelligent Robots and Systems}}, \bibinfo{pages}{3181--3186} (\bibinfo{year}{2005}).

\bibitem{artime2024robustness}
\bibinfo{author}{Artime, O.} \emph{et~al.}
\newblock \bibinfo{title}{Robustness and resilience of complex networks}.
\newblock \emph{\bibinfo{journal}{Nature Reviews Physics}} \textbf{\bibinfo{volume}{6}}, \bibinfo{pages}{114--131} (\bibinfo{year}{2024}).

\bibitem{granovetter1978threshold}
\bibinfo{author}{Granovetter, M.}
\newblock \bibinfo{title}{Threshold models of collective behavior}.
\newblock \emph{\bibinfo{journal}{American journal of sociology}} \textbf{\bibinfo{volume}{83}}, \bibinfo{pages}{1420--1443} (\bibinfo{year}{1978}).

\bibitem{pastor2001epidemic}
\bibinfo{author}{Pastor-Satorras, R.} \& \bibinfo{author}{Vespignani, A.}
\newblock \bibinfo{title}{Epidemic spreading in scale-free networks}.
\newblock \emph{\bibinfo{journal}{Physical review letters}} \textbf{\bibinfo{volume}{86}}, \bibinfo{pages}{3200} (\bibinfo{year}{2001}).

\bibitem{castellano2009statistical}
\bibinfo{author}{Castellano, C.}, \bibinfo{author}{Fortunato, S.} \& \bibinfo{author}{Loreto, V.}
\newblock \bibinfo{title}{{Statistical physics of social dynamics}}.
\newblock \emph{\bibinfo{journal}{Reviews of modern physics}} \textbf{\bibinfo{volume}{81}}, \bibinfo{pages}{591} (\bibinfo{year}{2009}).
\newblock \urlprefix\url{https://journals.aps.org/rmp/abstract/10.1103/RevModPhys.81.591}.

\bibitem{centola2010spread}
\bibinfo{author}{Centola, D.}
\newblock \bibinfo{title}{The spread of behavior in an online social network experiment}.
\newblock \emph{\bibinfo{journal}{science}} \textbf{\bibinfo{volume}{329}}, \bibinfo{pages}{1194--1197} (\bibinfo{year}{2010}).

\bibitem{PhysRevE.66.065102}
\bibinfo{author}{Motter, A.~E.} \& \bibinfo{author}{Lai, Y.-C.}
\newblock \bibinfo{title}{Cascade-based attacks on complex networks}.
\newblock \emph{\bibinfo{journal}{Phys. Rev. E}} \textbf{\bibinfo{volume}{66}}, \bibinfo{pages}{065102} (\bibinfo{year}{2002}).

\bibitem{Buldyrev_2010}
\bibinfo{author}{Buldyrev, S.~V.}, \bibinfo{author}{Parshani, R.}, \bibinfo{author}{Paul, G.}, \bibinfo{author}{Stanley, H.~E.} \& \bibinfo{author}{Havlin, S.}
\newblock \bibinfo{title}{Catastrophic cascade of failures in interdependent networks}.
\newblock \emph{\bibinfo{journal}{Nature}} \textbf{\bibinfo{volume}{464}}, \bibinfo{pages}{1025–1028} (\bibinfo{year}{2010}).

\bibitem{PhysRevE.69.045104}
\bibinfo{author}{Crucitti, P.}, \bibinfo{author}{Latora, V.} \& \bibinfo{author}{Marchiori, M.}
\newblock \bibinfo{title}{Model for cascading failures in complex networks}.
\newblock \emph{\bibinfo{journal}{Phys. Rev. E}} \textbf{\bibinfo{volume}{69}}, \bibinfo{pages}{045104} (\bibinfo{year}{2004}).

\bibitem{motter2004cascade}
\bibinfo{author}{Motter, A.~E.}
\newblock \bibinfo{title}{{Cascade control and defense in complex networks}}.
\newblock \emph{\bibinfo{journal}{Physical Review Letters}} \textbf{\bibinfo{volume}{93}}, \bibinfo{pages}{98701} (\bibinfo{year}{2004}).

\bibitem{DENG2007714}
\bibinfo{author}{Deng, K.}, \bibinfo{author}{Zhao, H.} \& \bibinfo{author}{Li, D.}
\newblock \bibinfo{title}{Effect of node deleting on network structure}.
\newblock \emph{\bibinfo{journal}{Physica A: Statistical Mechanics and its Applications}} \textbf{\bibinfo{volume}{379}}, \bibinfo{pages}{714--726} (\bibinfo{year}{2007}).

\bibitem{requiao2015fast}
\bibinfo{author}{Requi{\~a}o~da Cunha, B.}, \bibinfo{author}{Gonz{\'a}lez-Avella, J.~C.} \& \bibinfo{author}{Gon{\c{c}}alves, S.}
\newblock \bibinfo{title}{Fast fragmentation of networks using module-based attacks}.
\newblock \emph{\bibinfo{journal}{PloS one}} \textbf{\bibinfo{volume}{10}}, \bibinfo{pages}{e0142824} (\bibinfo{year}{2015}).

\bibitem{Altafini2013Consensus}
\bibinfo{author}{Altafini, C.}
\newblock \bibinfo{title}{Consensus problems on networks with antagonistic interactions}.
\newblock \emph{\bibinfo{journal}{IEEE Transactions on Automatic Control}} \textbf{\bibinfo{volume}{58}}, \bibinfo{pages}{935--946} (\bibinfo{year}{2013}).

\bibitem{Chen2017spectral}
\bibinfo{author}{Chen, W.}, \bibinfo{author}{Wang, D.}, \bibinfo{author}{Liu, J.}, \bibinfo{author}{Başar, T.} \& \bibinfo{author}{Qiu, L.}
\newblock \bibinfo{title}{On spectral properties of signed laplacians for undirected graphs}.
\newblock In \emph{\bibinfo{booktitle}{2017 IEEE 56th Annual Conference on Decision and Control (CDC)}}, \bibinfo{pages}{1999--2002} (\bibinfo{year}{2017}).

\bibitem{ahmadizadeh2017eigenvalues}
\bibinfo{author}{Ahmadizadeh, S.}, \bibinfo{author}{Shames, I.}, \bibinfo{author}{Martin, S.} \& \bibinfo{author}{Ne{\v{s}}i{\'c}, D.}
\newblock \bibinfo{title}{On eigenvalues of laplacian matrix for a class of directed signed graphs}.
\newblock \emph{\bibinfo{journal}{Linear Algebra and its Applications}} \textbf{\bibinfo{volume}{523}}, \bibinfo{pages}{281--306} (\bibinfo{year}{2017}).

\bibitem{reis2014avoiding}
\bibinfo{author}{Reis, S.~D.} \emph{et~al.}
\newblock \bibinfo{title}{Avoiding catastrophic failure in correlated networks of networks}.
\newblock \emph{\bibinfo{journal}{Nature Physics}} \textbf{\bibinfo{volume}{10}}, \bibinfo{pages}{762--767} (\bibinfo{year}{2014}).

\bibitem{albert2000error}
\bibinfo{author}{Albert, R.}, \bibinfo{author}{Jeong, H.} \& \bibinfo{author}{Barab{\'{a}}si, A.-L.}
\newblock \bibinfo{title}{{Error and attack tolerance of complex networks}}.
\newblock \emph{\bibinfo{journal}{Nature}} \textbf{\bibinfo{volume}{406}}, \bibinfo{pages}{378--382} (\bibinfo{year}{2000}).
\newblock \urlprefix\url{https://www.nature.com/articles/35019019 http://www.nature.com/doifinder/10.1038/35019019}.

\bibitem{pagani2013power}
\bibinfo{author}{Pagani, G.~A.} \& \bibinfo{author}{Aiello, M.}
\newblock \bibinfo{title}{{The power grid as a complex network: a survey}}.
\newblock \emph{\bibinfo{journal}{Physica A: Statistical Mechanics and its Applications}} \textbf{\bibinfo{volume}{392}}, \bibinfo{pages}{2688--2700} (\bibinfo{year}{2013}).

\bibitem{ren2007information}
\bibinfo{author}{Ren, W.}, \bibinfo{author}{Beard, R.~W.} \& \bibinfo{author}{Atkins, E.~M.}
\newblock \bibinfo{title}{Information consensus in multivehicle cooperative control}.
\newblock \emph{\bibinfo{journal}{IEEE Control systems magazine}} \textbf{\bibinfo{volume}{27}}, \bibinfo{pages}{71--82} (\bibinfo{year}{2007}).

\bibitem{hearnshaw2013complex}
\bibinfo{author}{Hearnshaw, E.~J.} \& \bibinfo{author}{Wilson, M.~M.}
\newblock \bibinfo{title}{A complex network approach to supply chain network theory}.
\newblock \emph{\bibinfo{journal}{International Journal of Operations \& Production Management}} \textbf{\bibinfo{volume}{33}}, \bibinfo{pages}{442--469} (\bibinfo{year}{2013}).

\bibitem{zhang2011optimal}
\bibinfo{author}{Zhang, H.}, \bibinfo{author}{Lewis, F.~L.} \& \bibinfo{author}{Das, A.}
\newblock \bibinfo{title}{Optimal design for synchronization of cooperative systems: state feedback, observer and output feedback}.
\newblock \emph{\bibinfo{journal}{IEEE Transactions on Automatic Control}} \textbf{\bibinfo{volume}{56}}, \bibinfo{pages}{1948--1952} (\bibinfo{year}{2011}).

\bibitem{pasqualetti2013attack}
\bibinfo{author}{Pasqualetti, F.}, \bibinfo{author}{Dorfler, F.} \& \bibinfo{author}{Bullo, F.}
\newblock \bibinfo{title}{{Attack Detection and Identification in Cyber-Physical Systems}}.
\newblock \emph{\bibinfo{journal}{IEEE Transactions on Automatic Control}} \textbf{\bibinfo{volume}{58}}, \bibinfo{pages}{2715--2729} (\bibinfo{year}{2013}).

\bibitem{KuraBOOK}
\bibinfo{author}{Kuramoto, Y.}
\newblock \bibinfo{title}{Chemical turbulence}.
\newblock In \emph{\bibinfo{booktitle}{Chemical oscillations, waves, and turbulence}}, \bibinfo{pages}{111--140} (\bibinfo{publisher}{Springer}, \bibinfo{year}{1984}).

\bibitem{johnson2017looplessness}
\bibinfo{author}{Johnson, S.} \& \bibinfo{author}{Jones, N.~S.}
\newblock \bibinfo{title}{Looplessness in networks is linked to trophic coherence}.
\newblock \emph{\bibinfo{journal}{Proceedings of the National Academy of Sciences}} \textbf{\bibinfo{volume}{114}}, \bibinfo{pages}{5618--5623} (\bibinfo{year}{2017}).

\bibitem{Asllani2018Structure}
\bibinfo{author}{Asllani, M.}, \bibinfo{author}{Lambiotte, R.} \& \bibinfo{author}{Carletti, T.}
\newblock \bibinfo{title}{Structure and dynamical behavior of non-normal networks}.
\newblock \emph{\bibinfo{journal}{Science Advances}} \textbf{\bibinfo{volume}{4}}, \bibinfo{pages}{eaau9403} (\bibinfo{year}{2018}).

\bibitem{asllani2018topological}
\bibinfo{author}{Asllani, M.} \& \bibinfo{author}{Carletti, T.}
\newblock \bibinfo{title}{Topological resilience in non-normal networked systems}.
\newblock \emph{\bibinfo{journal}{Physical Review E}} \textbf{\bibinfo{volume}{97}}, \bibinfo{pages}{042302} (\bibinfo{year}{2018}).

\bibitem{MUOLO2019Patterns}
\bibinfo{author}{Muolo, R.}, \bibinfo{author}{Asllani, M.}, \bibinfo{author}{Fanelli, D.}, \bibinfo{author}{Maini, P.~K.} \& \bibinfo{author}{Carletti, T.}
\newblock \bibinfo{title}{Patterns of non-normality in networked systems}.
\newblock \emph{\bibinfo{journal}{Journal of Theoretical Biology}} \textbf{\bibinfo{volume}{480}}, \bibinfo{pages}{81--91} (\bibinfo{year}{2019}).

\bibitem{muolo2020synchronization}
\bibinfo{author}{Muolo, R.}, \bibinfo{author}{Carletti, T.}, \bibinfo{author}{Gleeson, J.~P.} \& \bibinfo{author}{Asllani, M.}
\newblock \bibinfo{title}{Synchronization dynamics in non-normal networks: the trade-off for optimality}.
\newblock \emph{\bibinfo{journal}{Entropy}} \textbf{\bibinfo{volume}{23}}, \bibinfo{pages}{36} (\bibinfo{year}{2020}).

\bibitem{o2021hierarchical}
\bibinfo{author}{O'Brien, J.~D.}, \bibinfo{author}{Oliveira, K.~A.}, \bibinfo{author}{Gleeson, J.~P.} \& \bibinfo{author}{Asllani, M.}
\newblock \bibinfo{title}{Hierarchical route to the emergence of leader nodes in real-world networks}.
\newblock \emph{\bibinfo{journal}{Physical Review Research}} \textbf{\bibinfo{volume}{3}}, \bibinfo{pages}{023117} (\bibinfo{year}{2021}).

\bibitem{Duan2022Network}
\bibinfo{author}{Duan, C.}, \bibinfo{author}{Nishikawa, T.}, \bibinfo{author}{Eroglu, D.} \& \bibinfo{author}{Motter, A.~E.}
\newblock \bibinfo{title}{Network structural origin of instabilities in large complex systems}.
\newblock \emph{\bibinfo{journal}{Science Advances}} \textbf{\bibinfo{volume}{8}}, \bibinfo{pages}{1--12} (\bibinfo{year}{2022}).

\bibitem{nazerian2023commphys}
\bibinfo{author}{Nazerian, A.}, \bibinfo{author}{Panahi, S.} \& \bibinfo{author}{Sorrentino, F.}
\newblock \bibinfo{title}{Synchronization in networked systems with large parameter heterogeneity}.
\newblock \emph{\bibinfo{journal}{Communications Physics}} \textbf{\bibinfo{volume}{6}}, \bibinfo{pages}{253} (\bibinfo{year}{2023}).

\bibitem{Ramon}
\bibinfo{author}{Nartallo-Kaluarachchi, R.} \emph{et~al.}
\newblock \bibinfo{title}{Broken detailed balance and entropy production in directed networks}.
\newblock \emph{\bibinfo{journal}{Phys. Rev. E}} \textbf{\bibinfo{volume}{110}}, \bibinfo{pages}{034313} (\bibinfo{year}{2024}).

\bibitem{nazerian2024efficiency}
\bibinfo{author}{Nazerian, A.}, \bibinfo{author}{Hart, J.~D.}, \bibinfo{author}{Lodi, M.} \& \bibinfo{author}{Sorrentino, F.}
\newblock \bibinfo{title}{The efficiency of synchronization dynamics and the role of network syncreactivity}.
\newblock \emph{\bibinfo{journal}{Nature Communications}} \textbf{\bibinfo{volume}{15}}, \bibinfo{pages}{9003} (\bibinfo{year}{2024}).

\bibitem{muolo2024persistence}
\bibinfo{author}{Muolo, R.}, \bibinfo{author}{O’Brien, J.~D.}, \bibinfo{author}{Carletti, T.} \& \bibinfo{author}{Asllani, M.}
\newblock \bibinfo{title}{Persistence of chimera states and the challenge for synchronization in real-world networks}.
\newblock \emph{\bibinfo{journal}{The European Physical Journal B}} \textbf{\bibinfo{volume}{97}}, \bibinfo{pages}{6} (\bibinfo{year}{2024}).

\bibitem{nazerian2026frequency}
\bibinfo{author}{Nazerian, A.}, \bibinfo{author}{Asllani, M.}, \bibinfo{author}{Tyloo, M.}, \bibinfo{author}{Ku, W.~L.} \& \bibinfo{author}{Sorrentino, F.}
\newblock \bibinfo{title}{The frequency response of networks as open systems}.
\newblock \emph{\bibinfo{journal}{Nature Communications}} \textbf{\bibinfo{volume}{17}} (\bibinfo{year}{2026}).

\bibitem{johnson2020digraphs}
\bibinfo{author}{Johnson, S.}
\newblock \bibinfo{title}{Digraphs are different: Why directionality matters in complex systems}.
\newblock \emph{\bibinfo{journal}{Journal of Physics: Complexity}} \textbf{\bibinfo{volume}{1}}, \bibinfo{pages}{015003} (\bibinfo{year}{2020}).

\bibitem{wu2007synchronization}
\bibinfo{author}{Wu, C.~W.}
\newblock \emph{\bibinfo{title}{Synchronization in complex networks of nonlinear dynamical systems}} (\bibinfo{publisher}{World scientific}, \bibinfo{year}{2007}).

\bibitem{fiedler1973algebraic}
\bibinfo{author}{Fiedler, M.}
\newblock \bibinfo{title}{Algebraic connectivity of graphs}.
\newblock \emph{\bibinfo{journal}{Czechoslovak mathematical journal}} \textbf{\bibinfo{volume}{23}}, \bibinfo{pages}{298--305} (\bibinfo{year}{1973}).

\bibitem{olfati2007consensus}
\bibinfo{author}{Olfati-Saber, R.}, \bibinfo{author}{Fax, J.~A.} \& \bibinfo{author}{Murray, R.~M.}
\newblock \bibinfo{title}{{Consensus and cooperation in networked multi-agent systems}}.
\newblock \emph{\bibinfo{journal}{Proceedings of the IEEE}} \textbf{\bibinfo{volume}{95}}, \bibinfo{pages}{215--233} (\bibinfo{year}{2007}).

\bibitem{trefethen1991pseudospectra}
\bibinfo{author}{Trefethen, L.~N.}
\newblock \bibinfo{title}{Pseudospectra of matrices}.
\newblock \emph{\bibinfo{journal}{Numerical analysis}} \textbf{\bibinfo{volume}{91}}, \bibinfo{pages}{234--266} (\bibinfo{year}{1991}).

\bibitem{chung1997spectral}
\bibinfo{author}{Chung, F.~R.}
\newblock \emph{\bibinfo{title}{Spectral graph theory}}, vol.~\bibinfo{volume}{92} (\bibinfo{publisher}{American Mathematical Soc.}, \bibinfo{year}{1997}).

\bibitem{konect}
\bibinfo{author}{Kunegis, J.}
\newblock \bibinfo{title}{{KONECT} -- {The} {Koblenz} {Network} {Collection}}.
\newblock In \emph{\bibinfo{booktitle}{Proc. Int. Conf. on World Wide Web Companion}}, \bibinfo{pages}{1343--1350} (\bibinfo{year}{2013}).
\newblock \urlprefix\url{http://dl.acm.org/citation.cfm?id=2488173}.

\bibitem{fujisaka1983stability}
\bibinfo{author}{Fujisaka, H.} \& \bibinfo{author}{Yamada, T.}
\newblock \bibinfo{title}{Stability theory of synchronized motion in coupled-oscillator systems}.
\newblock \emph{\bibinfo{journal}{Progress of theoretical physics}} \textbf{\bibinfo{volume}{69}}, \bibinfo{pages}{32--47} (\bibinfo{year}{1983}).

\bibitem{pecora1998master}
\bibinfo{author}{Pecora, L.~M.} \& \bibinfo{author}{Carroll, T.~L.}
\newblock \bibinfo{title}{{Master stability functions for synchronized coupled systems}}.
\newblock \emph{\bibinfo{journal}{Physical review letters}} \textbf{\bibinfo{volume}{80}}, \bibinfo{pages}{2109} (\bibinfo{year}{1998}).
\newblock \urlprefix\url{https://journals.aps.org/prl/abstract/10.1103/PhysRevLett.80.2109}.

\bibitem{sorrentino2007controllability}
\bibinfo{author}{Sorrentino, F.}, \bibinfo{author}{Di~Bernardo, M.}, \bibinfo{author}{Garofalo, F.} \& \bibinfo{author}{Chen, G.}
\newblock \bibinfo{title}{Controllability of complex networks via pinning}.
\newblock \emph{\bibinfo{journal}{Physical Review E—Statistical, Nonlinear, and Soft Matter Physics}} \textbf{\bibinfo{volume}{75}}, \bibinfo{pages}{046103} (\bibinfo{year}{2007}).

\bibitem{goh2001universal}
\bibinfo{author}{Goh, K.-I.}, \bibinfo{author}{Kahng, B.} \& \bibinfo{author}{Kim, D.}
\newblock \bibinfo{title}{{Universal Behavior of Load Distribution in Scale-Free Networks}}.
\newblock \emph{\bibinfo{journal}{Physical Review Letters}} \textbf{\bibinfo{volume}{87}}, \bibinfo{pages}{278701} (\bibinfo{year}{2001}).

\bibitem{nishikawa2015comparative}
\bibinfo{author}{Nishikawa, T.} \& \bibinfo{author}{Motter, A.~E.}
\newblock \bibinfo{title}{{Comparative analysis of existing models for power-grid synchronization}}.
\newblock \emph{\bibinfo{journal}{New Journal of Physics}} \textbf{\bibinfo{volume}{17}}, \bibinfo{pages}{015012} (\bibinfo{year}{2015}).

\bibitem{bhatta2021modal}
\bibinfo{author}{Bhatta, K.}, \bibinfo{author}{Hayat, M.~M.} \& \bibinfo{author}{Sorrentino, F.}
\newblock \bibinfo{title}{Modal decomposition of the linear swing equation in networks with symmetries}.
\newblock \emph{\bibinfo{journal}{IEEE Transactions on Network Science and Engineering}} \textbf{\bibinfo{volume}{8}}, \bibinfo{pages}{2482--2494} (\bibinfo{year}{2021}).

\bibitem{oh2015survey}
\bibinfo{author}{Oh, K.-K.}, \bibinfo{author}{Park, M.-C.} \& \bibinfo{author}{Ahn, H.-S.}
\newblock \bibinfo{title}{A survey of multi-agent formation control}.
\newblock \emph{\bibinfo{journal}{Automatica}} \textbf{\bibinfo{volume}{53}}, \bibinfo{pages}{424--440} (\bibinfo{year}{2015}).

\bibitem{wang2021formation}
\bibinfo{author}{Wang, Y.}, \bibinfo{author}{Yue, Y.}, \bibinfo{author}{Shan, M.}, \bibinfo{author}{He, L.} \& \bibinfo{author}{Wang, D.}
\newblock \bibinfo{title}{Formation reconstruction and trajectory replanning for multi-uav patrol}.
\newblock \emph{\bibinfo{journal}{IEEE/ASME Transactions on Mechatronics}} \textbf{\bibinfo{volume}{26}}, \bibinfo{pages}{719--729} (\bibinfo{year}{2021}).

\bibitem{yu2022decentralized}
\bibinfo{author}{Yu, X.} \& \bibinfo{author}{Su, R.}
\newblock \bibinfo{title}{Decentralized circular formation control of nonholonomic mobile robots under a directed sensor graph}.
\newblock \emph{\bibinfo{journal}{IEEE Transactions on Automatic Control}} \textbf{\bibinfo{volume}{68}}, \bibinfo{pages}{3656--3663} (\bibinfo{year}{2022}).

\bibitem{Morbidi2022Functions}
\bibinfo{author}{Morbidi, F.}
\newblock \bibinfo{title}{Functions of the laplacian matrix with application to distributed formation control}.
\newblock \emph{\bibinfo{journal}{IEEE Transactions on Control of Network Systems}} \textbf{\bibinfo{volume}{9}}, \bibinfo{pages}{1459--1467} (\bibinfo{year}{2022}).

\bibitem{barabasi1999emergence}
\bibinfo{author}{Barab{\'{a}}si, A.-L.} \& \bibinfo{author}{Albert, R.}
\newblock \bibinfo{title}{{Emergence of scaling in random networks}}.
\newblock \emph{\bibinfo{journal}{Science}} \textbf{\bibinfo{volume}{286}}, \bibinfo{pages}{509--512} (\bibinfo{year}{1999}).
\newblock \urlprefix\url{http://science.sciencemag.org/content/286/5439/509}.

\bibitem{liu2011controllability}
\bibinfo{author}{Liu, Y.-Y.}, \bibinfo{author}{Slotine, J.-J.} \& \bibinfo{author}{Barab{\'{a}}si, A.-L.}
\newblock \bibinfo{title}{{Controllability of complex networks}}.
\newblock \emph{\bibinfo{journal}{Nature}} \textbf{\bibinfo{volume}{473}}, \bibinfo{pages}{167--173} (\bibinfo{year}{2011}).

\bibitem{menichetti2014network}
\bibinfo{author}{Menichetti, G.}, \bibinfo{author}{Dall’Asta, L.} \& \bibinfo{author}{Bianconi, G.}
\newblock \bibinfo{title}{Network controllability is determined by the density of low in-degree and out-degree nodes}.
\newblock \emph{\bibinfo{journal}{Physical review letters}} \textbf{\bibinfo{volume}{113}}, \bibinfo{pages}{078701} (\bibinfo{year}{2014}).

\bibitem{granovetter1973strength}
\bibinfo{author}{Granovetter, M.~S.}
\newblock \bibinfo{title}{The strength of weak ties}.
\newblock \emph{\bibinfo{journal}{American journal of sociology}} \textbf{\bibinfo{volume}{78}}, \bibinfo{pages}{1360--1380} (\bibinfo{year}{1973}).

\bibitem{allesina2012stability}
\bibinfo{author}{Allesina, S.} \& \bibinfo{author}{Tang, S.}
\newblock \bibinfo{title}{Stability criteria for complex ecosystems}.
\newblock \emph{\bibinfo{journal}{Nature}} \textbf{\bibinfo{volume}{483}}, \bibinfo{pages}{205--208} (\bibinfo{year}{2012}).

\bibitem{boccaletti2006complex}
\bibinfo{author}{Boccaletti, S.}, \bibinfo{author}{Latora, V.}, \bibinfo{author}{Moreno, Y.}, \bibinfo{author}{Chavez, M.} \& \bibinfo{author}{Hwang, D.-U.}
\newblock \bibinfo{title}{Complex networks: Structure and dynamics}.
\newblock \emph{\bibinfo{journal}{Physics reports}} \textbf{\bibinfo{volume}{424}}, \bibinfo{pages}{175--308} (\bibinfo{year}{2006}).

\bibitem{smith2020survey}
\bibinfo{author}{Smith, J.~C.} \& \bibinfo{author}{Song, Y.}
\newblock \bibinfo{title}{A survey of network interdiction models and algorithms}.
\newblock \emph{\bibinfo{journal}{European Journal of Operational Research}} \textbf{\bibinfo{volume}{283}}, \bibinfo{pages}{797--811} (\bibinfo{year}{2020}).

\bibitem{Wood1993}
\bibinfo{author}{Wood, R.~K.}
\newblock \bibinfo{title}{Deterministic network interdiction}.
\newblock \emph{\bibinfo{journal}{Mathematical and Computer Modelling}} \textbf{\bibinfo{volume}{17}}, \bibinfo{pages}{1--18} (\bibinfo{year}{1993}).

\bibitem{IsraeliWood2002}
\bibinfo{author}{Israeli, E.} \& \bibinfo{author}{Wood, R.~K.}
\newblock \bibinfo{title}{Shortest-path network interdiction}.
\newblock \emph{\bibinfo{journal}{Networks: An International Journal}} \textbf{\bibinfo{volume}{40}}, \bibinfo{pages}{97--111} (\bibinfo{year}{2002}).

\bibitem{SmithSong2020}
\bibinfo{author}{Smith, J.~C.} \& \bibinfo{author}{Song, Y.}
\newblock \bibinfo{title}{A survey of network interdiction models and algorithms}.
\newblock \emph{\bibinfo{journal}{European Journal of Operational Research}} \textbf{\bibinfo{volume}{283}}, \bibinfo{pages}{797--811} (\bibinfo{year}{2020}).

\bibitem{CappaneraScaparra2011}
\bibinfo{author}{Cappanera, P.} \& \bibinfo{author}{Scaparra, M.~P.}
\newblock \bibinfo{title}{Optimal allocation of protective resources in shortest-path networks}.
\newblock \emph{\bibinfo{journal}{Transportation Science}} \textbf{\bibinfo{volume}{45}}, \bibinfo{pages}{64--80} (\bibinfo{year}{2011}).

\bibitem{HolmeKimYoonHan2002}
\bibinfo{author}{Holme, P.}, \bibinfo{author}{Kim, B.~J.}, \bibinfo{author}{Yoon, C.~N.} \& \bibinfo{author}{Han, S.~K.}
\newblock \bibinfo{title}{Attack vulnerability of complex networks}.
\newblock \emph{\bibinfo{journal}{Physical review E}} \textbf{\bibinfo{volume}{65}}, \bibinfo{pages}{056109} (\bibinfo{year}{2002}).

\bibitem{carletti2023global}
\bibinfo{author}{Carletti, T.}, \bibinfo{author}{Giambagli, L.} \& \bibinfo{author}{Bianconi, G.}
\newblock \bibinfo{title}{Global topological synchronization on simplicial and cell complexes}.
\newblock \emph{\bibinfo{journal}{Physical review letters}} \textbf{\bibinfo{volume}{130}}, \bibinfo{pages}{187401} (\bibinfo{year}{2023}).

\bibitem{gallo2022synchronization}
\bibinfo{author}{Gallo, L.} \emph{et~al.}
\newblock \bibinfo{title}{Synchronization induced by directed higher-order interactions}.
\newblock \emph{\bibinfo{journal}{Communications Physics}} \textbf{\bibinfo{volume}{5}}, \bibinfo{pages}{263} (\bibinfo{year}{2022}).

\bibitem{skardal2019abrupt}
\bibinfo{author}{Skardal, P.~S.} \& \bibinfo{author}{Arenas, A.}
\newblock \bibinfo{title}{Abrupt desynchronization and extensive multistability in globally coupled oscillator simplexes}.
\newblock \emph{\bibinfo{journal}{Physical review letters}} \textbf{\bibinfo{volume}{122}}, \bibinfo{pages}{248301} (\bibinfo{year}{2019}).

\bibitem{millan2020explosive}
\bibinfo{author}{Mill{\'a}n, A.~P.}, \bibinfo{author}{Torres, J.~J.} \& \bibinfo{author}{Bianconi, G.}
\newblock \bibinfo{title}{Explosive higher-order kuramoto dynamics on simplicial complexes}.
\newblock \emph{\bibinfo{journal}{Physical Review Letters}} \textbf{\bibinfo{volume}{124}}, \bibinfo{pages}{218301} (\bibinfo{year}{2020}).

\bibitem{Neubert1997ALTERNATIVES}
\bibinfo{author}{Neubert, M.~G.} \& \bibinfo{author}{Caswell, H.}
\newblock \bibinfo{title}{Alternatives to resilience for measuring the responses of ecological systems to perturbations}.
\newblock \emph{\bibinfo{journal}{Ecology}} \textbf{\bibinfo{volume}{78}}, \bibinfo{pages}{653--665} (\bibinfo{year}{1997}).

\bibitem{nazerian2023single}
\bibinfo{author}{Nazerian, A.}, \bibinfo{author}{Phillips, D.}, \bibinfo{author}{Makse, H.~A.} \& \bibinfo{author}{Sorrentino, F.}
\newblock \bibinfo{title}{Single-integrator consensus dynamics over minimally reactive networks}.
\newblock \emph{\bibinfo{journal}{IEEE Control Systems Letters}}  (\bibinfo{year}{2023}).

\bibitem{Farrell1996Generalized}
\bibinfo{author}{Farrell, B.~F.} \& \bibinfo{author}{Ioannou, P.~J.}
\newblock \bibinfo{title}{Generalized stability theory. part i: Autonomous operators}.
\newblock \emph{\bibinfo{journal}{Journal of Atmospheric Sciences}} \textbf{\bibinfo{volume}{53}}, \bibinfo{pages}{2025 -- 2040} (\bibinfo{year}{1996}).

\bibitem{Hennequin2012Nonnormal}
\bibinfo{author}{Hennequin, G.}, \bibinfo{author}{Vogels, T.~P.} \& \bibinfo{author}{Gerstner, W.}
\newblock \bibinfo{title}{Non-normal amplification in random balanced neuronal networks}.
\newblock \emph{\bibinfo{journal}{Phys. Rev. E}} \textbf{\bibinfo{volume}{86}}, \bibinfo{pages}{011909} (\bibinfo{year}{2012}).

\bibitem{Tang2014Reactivity}
\bibinfo{author}{Tang, S.} \& \bibinfo{author}{Allesina, S.}
\newblock \bibinfo{title}{Reactivity and stability of large ecosystems}.
\newblock \emph{\bibinfo{journal}{Frontiers in Ecology and Evolution}} \textbf{\bibinfo{volume}{2}} (\bibinfo{year}{2014}).

\bibitem{Biancalani2017Giant}
\bibinfo{author}{Biancalani, T.}, \bibinfo{author}{Jafarpour, F.} \& \bibinfo{author}{Goldenfeld, N.}
\newblock \bibinfo{title}{Giant amplification of noise in fluctuation-induced pattern formation}.
\newblock \emph{\bibinfo{journal}{Phys. Rev. Lett.}} \textbf{\bibinfo{volume}{118}}, \bibinfo{pages}{018101} (\bibinfo{year}{2017}).

\bibitem{Gudowska2020From}
\bibinfo{author}{Gudowska-Nowak, E.}, \bibinfo{author}{Nowak, M.~A.}, \bibinfo{author}{Chialvo, D.~R.}, \bibinfo{author}{Ochab, J.~K.} \& \bibinfo{author}{Tarnowski, W.}
\newblock \bibinfo{title}{{From Synaptic Interactions to Collective Dynamics in Random Neuronal Networks Models: Critical Role of Eigenvectors and Transient Behavior}}.
\newblock \emph{\bibinfo{journal}{Neural Computation}} \textbf{\bibinfo{volume}{32}}, \bibinfo{pages}{395--423} (\bibinfo{year}{2020}).

\bibitem{Lindmark2021Centrality}
\bibinfo{author}{Lindmark, G.} \& \bibinfo{author}{Altafini, C.}
\newblock \bibinfo{title}{Centrality measures and the role of non-normality for network control energy reduction}.
\newblock \emph{\bibinfo{journal}{IEEE Control Systems Letters}} \textbf{\bibinfo{volume}{5}}, \bibinfo{pages}{1013--1018} (\bibinfo{year}{2021}).

\bibitem{Nazerian2023Reactability}
\bibinfo{author}{Nazerian, A.}, \bibinfo{author}{Phillips, D.}, \bibinfo{author}{Frasca, M.} \& \bibinfo{author}{Sorrentino, F.}
\newblock \bibinfo{title}{The reactability of discrete time systems}.
\newblock \emph{\bibinfo{journal}{IEEE Control Systems Letters}} \textbf{\bibinfo{volume}{7}}, \bibinfo{pages}{3657--3662} (\bibinfo{year}{2023}).

\bibitem{bamieh2022tutorial}
\bibinfo{author}{Bamieh, B.}
\newblock \bibinfo{title}{A tutorial on matrix perturbation theory (using compact matrix notation)} (\bibinfo{year}{2022}).
\newblock \urlprefix\url{https://arxiv.org/abs/2002.05001}.
\newblock \eprint{2002.05001}.

\end{thebibliography}


\section*{Funding Statement}
FS acknowledges support from grants AFOSR FA9550-24-1-0214 and Oak Ridge National Laboratory 006321-00001A. MA acknowledges funding support from FSU CRS-SEED program, `Structure and Dynamics of Non-Normal Networks.' HAM acknowledges funding  from NSF-HNDS Grant 2214217. 



\section*{Author Contributions}
AN worked on the theory and numerical simulations. ST worked on the numerical simulations and the figures. MA worked on the theory and contributed to writing the paper. DP and HAM provided essential input and guidance.    FS worked on the theory and supervised the research. 


\section*{Competing Interests Statement}
The authors declare no competing interests


\end{document}
